\part{自测与参考}

\chapter{测验题与详细解答}
\label{quiz-questions-detailed-answers}

% This chapter provides a comprehensive set of questions designed to test and reinforce your understanding of the material covered throughout this guide. Each question targets a key concept, algorithm, or system design decision---the kind of knowledge that distinguishes surface-level familiarity from genuine expertise. Use these questions for self-assessment: attempt your own answer before reading the detailed solution. The questions progress from foundational concepts (LLM architecture, reinforcement learning basics) through core algorithms (PPO, DPO, GRPO) to advanced system design and agentic AI topics.
本章提供一整套题目，用于检验和巩固你对本书所讲内容的理解。每道题目都瞄准一个关键概念、算法或系统设计决策——正是这种知识把表面熟悉与真正的专长区分开来。请将这些题目用于自我检验：先尝试自己作答，再阅读详细解答。题目从基础概念（LLM 架构、强化学习基础）逐步推进到核心算法（PPO、DPO、GRPO），最终覆盖高级系统设计与智能体 AI 主题。

\section{基础题}
\label{foundations-questions}

% \begin{questionbox}[Q0a: What is the role of the attention mechanism in a decoder-only Transformer? Why is it causal?]
\begin{questionbox}[Q0a：在 decoder-only Transformer 中，attention 机制的作用是什么？为什么它是因果的？]
% \textbf{Answer}: The attention mechanism allows each token to attend to (i.e., compute a weighted combination of) representations from other tokens. In a decoder-only Transformer, attention is \textbf{causal} (also called \emph{autoregressive}): token $t$ can only attend to tokens $1, \ldots, t$, never to future tokens $t+1, \ldots, T$.
\textbf{答}：attention 机制允许每个 Token 关注（即对其表示做加权组合）其他 Token 的表示。在 decoder-only Transformer 中，attention 是\textbf{因果的}（也称为\emph{自回归}）：Token $t$ 只能关注 Token $1, \ldots, t$，绝不能关注未来 Token $t+1, \ldots, T$。

% \textbf{Why causal?} Because the model generates text left-to-right. At inference time, future tokens literally do not exist yet. The causal mask during training simulates this constraint so that the model learns to predict each token using only its left context. Mathematically, the attention matrix is masked:
\textbf{为什么是因果的？}因为模型从左到右生成文本。推理时，未来 Token 字面意义上还不存在。训练时的因果掩码模拟了这一约束，使模型学会仅用左侧上下文预测每个 Token。在数学上，attention 矩阵被掩盖：
\[
\text{Attention}(Q, K, V) = \text{softmax}\!\left(\frac{QK^\top}{\sqrt{d_k}} + M\right) V
\]
% where $M_{ij} = -\infty$ for $j > i$ (future positions), forcing those attention weights to zero.
其中对于 $j > i$（未来位置）有 $M_{ij} = -\infty$，强制将这些 attention 权重置零。

% \textbf{Practical implication}: This enables the KV-cache optimization at inference---since past tokens’ keys and values never change, they can be cached and reused, reducing generation from $O(T^2)$ to $O(T)$ per new token.
\textbf{实践含义}：这使得推理时的 KV-cache 优化成为可能——由于过去 Token 的 keys 和 values 永远不变，可被缓存复用，将每个新 Token 的生成代价从 $O(T^2)$ 降至 $O(T)$。

% \emph{\textbf{Review:} Chapter~1 (LLM Architecture and Optimization Methods).}
\emph{\textbf{复习：} 第~1 章（LLM 架构与优化方法）。}
\end{questionbox}

% \begin{questionbox}[Q0b: Explain FlashAttention. What problem does it solve and how?]
\begin{questionbox}[Q0b：解释 FlashAttention。它解决什么问题，又是怎么解决的？]
% \textbf{Answer}: Standard attention computes the full $T \times T$ attention matrix, which requires $O(T^2)$ memory and is \textbf{memory-bandwidth bound}---the GPU spends most of its time moving data between HBM (slow, large) and SRAM (fast, small), not doing actual computation.
\textbf{答}：标准 attention 会计算完整的 $T \times T$ attention 矩阵，需要 $O(T^2)$ 显存，并且是\textbf{受显存带宽约束的}——GPU 大部分时间都花在 HBM（慢、大）与 SRAM（快、小）之间搬运数据，而不是做真正的计算。

% \textbf{FlashAttention’s insight}: Never materialize the full attention matrix in HBM. Instead, tile the computation into blocks that fit in SRAM, compute attention \emph{block by block} using an online softmax algorithm, and write only the final output to HBM.
\textbf{FlashAttention 的洞见}：永远不要在 HBM 中物化完整的 attention 矩阵。相反，把计算切分成可放进 SRAM 的分块（tile），用在线 Softmax（Online Softmax）算法\emph{逐块}计算 attention，并只把最终输出写入 HBM。

% \textbf{Key techniques}:
\textbf{关键技术}：

\begin{enumerate}
  % \item \textbf{Tiling}: Split Q, K, V into blocks of size $B_r \times B_c$ that fit in SRAM
  \item \textbf{分块（Tiling）}：将 Q、K、V 拆成大小为 $B_r \times B_c$、可装入 SRAM 的块。
  % \item \textbf{Online softmax}: Track running max and sum to compute softmax incrementally without the full row
  \item \textbf{在线 Softmax（Online Softmax）}：维护一个滑动的 max 和 sum，以增量方式计算 softmax，无需访问完整一行。
  % \item \textbf{Recomputation}: In the backward pass, recompute attention from Q, K, V (cheap) rather than storing the $T \times T$ matrix (expensive)
  \item \textbf{重算（Recomputation）}：在反向传播时，从 Q、K、V 重新计算 attention（开销小），而不存储 $T \times T$ 矩阵（开销大）。
\end{enumerate}

% \textbf{Result}: $O(T)$ HBM memory (instead of $O(T^2)$), 2--4$\times$ wall-clock speedup, exact same numerical output (not an approximation).
\textbf{结果}：HBM 显存复杂度为 $O(T)$（而非 $O(T^2)$），墙钟（wall-clock）时间加速 2--4$\times$，数值输出完全一致（并非近似）。

% \emph{\textbf{Review:} Chapters~1--2 (LLM Architecture; Systems Foundations).}
\emph{\textbf{复习：} 第~1--2 章（LLM 架构；系统基础）。}
\end{questionbox}

% \begin{questionbox}[Q0c: What is the difference between SFT, RLHF, and DPO at a high level? When do you use each?]
\begin{questionbox}[Q0c：从高层次看，SFT、RLHF 与 DPO 有什么区别？分别在什么场景下使用？]
\textbf{答}：

\begin{itemize}
  % \item \textbf{SFT (Supervised Fine-Tuning)}: Train the model to imitate high-quality demonstrations. Loss: next-token prediction on curated data. Teaches \emph{format} and \emph{style}.
  \item \textbf{SFT（监督微调，Supervised Fine-Tuning）}：训练模型去模仿高质量示范。Loss：在精选数据上做下一个 Token 预测。教会模型\emph{格式}与\emph{风格}。
  % \item \textbf{RLHF}: Train a reward model from human preferences, then optimize the policy against it using RL (PPO). The model explores beyond the demonstration data. Teaches \emph{what humans prefer}.
  \item \textbf{RLHF}：从人类偏好训练一个 reward 模型，然后用 RL（PPO）针对它优化 policy。模型会探索示范数据之外的空间。教会模型\emph{人类偏好什么}。
  % \item \textbf{DPO}: Skip the reward model. Directly optimize the policy on preference pairs $(y_w, y_l)$ using a contrastive loss. Same goal as RLHF but simpler pipeline.
  \item \textbf{DPO}：跳过 reward 模型。直接在偏好对 $(y_w, y_l)$ 上用对比 loss 优化 policy。目标与 RLHF 相同，但流水线更简单。
\end{itemize}

% \textbf{Typical pipeline}: SFT first (gives the model a good starting point), then RLHF or DPO (refines preferences). SFT alone tends to produce verbose, hedge-heavy responses. RLHF/DPO makes outputs more direct and aligned with human intent.
\textbf{典型流水线}：先做 SFT（提供良好起点），再做 RLHF 或 DPO（细化偏好）。仅做 SFT 倾向产出冗长、过度模棱两可的回答。RLHF/DPO 让输出更直接、更贴合人类意图。

% \textbf{When to use each}: SFT when you have gold-standard outputs. DPO when you have preference pairs but limited compute. RLHF (PPO) when you need maximum quality and can afford the infrastructure.
\textbf{各自适用场景}：有黄金标准输出 $\rightarrow$ SFT。有偏好对但算力有限 $\rightarrow$ DPO。需要极致质量且能承担基础设施成本 $\rightarrow$ RLHF（PPO）。

% \emph{\textbf{Review:} Chapters~5, 6, and~10 (PPO; DPO; SFT Best Practices).}
\emph{\textbf{复习：} 第~5、6 与~10 章（PPO；DPO；SFT 最佳实践）。}
\end{questionbox}

% \begin{questionbox}[Q0d: What is a reward model? How is it trained and what can go wrong?]
\begin{questionbox}[Q0d：什么是 reward 模型？它如何训练，可能出什么问题？]
% \textbf{Answer}: A reward model (RM) is a neural network that takes a (prompt, response) pair and outputs a scalar score indicating quality. It is trained on human preference data: given pairs $(y_w, y_l)$ where $y_w$ is preferred, the RM learns to assign $R(y_w) > R(y_l)$.
\textbf{答}：reward 模型（Reward Model, RM）是一个神经网络，输入 (prompt, response) 对，输出一个表示质量的标量分数。它在人类偏好数据上训练：给定 $(y_w, y_l)$ 对（其中 $y_w$ 被偏好），RM 学习赋予 $R(y_w) > R(y_l)$。

% \textbf{Training}: Bradley-Terry loss: $\mathcal{L} = -\log\sigma(R(y_w) - R(y_l))$. Architecture: typically the same transformer as the policy, with the LM head replaced by a scalar projection.
\textbf{训练}：Bradley-Terry loss：$\mathcal{L} = -\log\sigma(R(y_w) - R(y_l))$。架构：通常与 policy 共用同一个 Transformer，将 LM head 替换为标量投影。

% \textbf{What can go wrong}:
\textbf{可能出什么问题}：

\begin{enumerate}
  % \item \textbf{Reward hacking}: The policy finds outputs that score high on the RM but are actually low quality (e.g., excessively long, repetitive, or containing specific phrases the RM was biased toward)
  \item \textbf{Reward hacking}：policy 找到了在 RM 上得分很高但实际质量很差的输出（例如过度冗长、重复，或包含 RM 存在偏好的特定短语）。
  % \item \textbf{Distribution shift}: The RM was trained on outputs from an earlier policy. As training progresses, the current policy generates out-of-distribution outputs the RM cannot score accurately
  \item \textbf{分布漂移（Distribution Shift）}：RM 是在更早的 policy 输出上训练的。随着训练推进，当前 policy 生成的是 RM 无法准确打分的分布外输出。
  % \item \textbf{Label noise}: Human annotators disagree, are tired, or apply inconsistent criteria. This noise propagates into RM predictions
  \item \textbf{标签噪声}：人类标注者意见不一致、疲劳，或使用不一致的标准。这种噪声会传播到 RM 的预测中。
  % \item \textbf{Overconfidence}: The RM assigns extreme scores to outputs it has never seen, providing misleading gradient signal
  \item \textbf{过度自信}：RM 对从未见过的输出赋予极端分数，提供具有误导性的梯度信号。
\end{enumerate}

% \emph{\textbf{Review:} Chapter~9 (Reward Model Training).}
\emph{\textbf{复习：} 第~9 章（Reward 模型训练）。}
\end{questionbox}

% \begin{questionbox}[Q0e: Explain the exploration-exploitation trade-off in RL. How does it manifest in LLM training?]
\begin{questionbox}[Q0e：解释 RL 中的探索-利用（exploration-exploitation）权衡。它在 LLM 训练中如何体现？]
% \textbf{Answer}: In RL, an agent must balance:
\textbf{答}：在 RL 中，Agent 必须在以下两者间取得平衡：

\begin{itemize}
  % \item \textbf{Exploitation}: choosing actions known to yield high reward (greedy behavior)
  \item \textbf{利用（Exploitation）}：选择已知能产生高 reward 的动作（贪婪行为）。
  % \item \textbf{Exploration}: trying new actions that might yield even higher reward (but might also fail)
  \item \textbf{探索（Exploration）}：尝试可能带来更高 reward（但也可能失败）的新动作。
\end{itemize}

% \textbf{In LLM training}: The policy is the language model. ``Actions'' are token choices. ``Exploitation'' means generating responses similar to what already scored well. ``Exploration'' means trying novel phrasings, structures, or reasoning paths.
\textbf{在 LLM 训练中}：policy 就是语言模型。“动作”是 Token 的选择。“利用”指生成与已得高分相似的回答；“探索”指尝试新颖的措辞、结构或推理路径。

% \textbf{How it manifests}:
\textbf{表现形式}：

\begin{itemize}
  % \item \textbf{Temperature during generation}: Higher temperature = more exploration. GRPO uses temperature 1.0 to get diverse samples within each group.
  \item \textbf{生成时的温度}：温度越高 = 探索越多。GRPO 使用温度 1.0 以在每个 group 内获得多样的样本。
  % \item \textbf{KL penalty}: Acts as an anti-exploration brake---prevents the policy from straying too far from the reference model. Without it, the policy might collapse to a single high-reward template (mode collapse).
  \item \textbf{KL 惩罚}：作为一种反探索的刹车——防止 policy 偏离参考模型过远。否则 policy 可能会塌缩为单一的高 reward 模板（mode collapse）。
  % \item \textbf{Group sampling in GRPO}: Generating $G$ responses per prompt explicitly explores the output space, then reinforces above-average responses.
  \item \textbf{GRPO 中的 group 采样}：每个 prompt 生成 $G$ 个回答以显式探索输出空间，然后强化高于平均水平的回答。
\end{itemize}

% \textbf{The tension}: Too little exploration $\rightarrow$ the model gets stuck in local optima (always giving the same safe answer). Too much $\rightarrow$ training is unstable, quality fluctuates wildly.
\textbf{张力}：探索过少 $\rightarrow$ 模型陷入局部最优（永远给同一个安全答案）。探索过多 $\rightarrow$ 训练不稳定，质量剧烈波动。

% \emph{\textbf{Review:} Chapters~3 and~7 (Introduction to RL; GRPO).}
\emph{\textbf{复习：} 第~3 与~7 章（RL 简介；GRPO）。}
\end{questionbox}

\section{核心算法题}
\label{core-algorithm-questions}

% \begin{questionbox}[Q1: Explain PPO’s clipped objective. Why does it work better than vanilla PG?]
\begin{questionbox}[Q1：解释 PPO 的裁剪目标。为什么它比朴素 policy gradient 更有效？]
% \textbf{Answer}: Vanilla policy gradient: $\nabla J = \mathbb{E}[\nabla\log\pi(a|s) \cdot \hat{A}]$. Problem: one lucky/unlucky sample can produce a huge gradient $\rightarrow$ policy jumps to a bad region $\rightarrow$ generates garbage $\rightarrow$ next gradient makes it worse $\rightarrow$ unrecoverable ``death spiral.''
\textbf{答}：朴素 policy gradient：$\nabla J = \mathbb{E}[\nabla\log\pi(a|s) \cdot \hat{A}]$。问题：一个幸运/不幸的样本就能产生巨大的梯度 $\rightarrow$ policy 跳到糟糕区域 $\rightarrow$ 生成乱码 $\rightarrow$ 下一个梯度让情况更糟 $\rightarrow$ 进入无法挽回的“死亡螺旋”。

% \textbf{PPO’s solution}: Clip the probability ratio $r = \pi_\text{new}/\pi_\text{old}$ to $[0.8, 1.2]$.
\textbf{PPO 的方案}：将概率比 $r = \pi_\text{new}/\pi_\text{old}$ 裁剪到 $[0.8, 1.2]$。

% \textbf{Mechanics}: For good actions ($\hat{A}>0$): objective is $\min(r\hat{A}, 1.2\hat{A})$. Once $r$ exceeds 1.2, no further benefit --- stops the policy from over-committing to one example. For bad actions ($\hat{A}<0$): objective is $\min(r\hat{A}, 0.8\hat{A})$. Once $r$ drops below 0.8, penalty stops growing --- prevents catastrophic forgetting.
\textbf{机制}：对好的动作（$\hat{A}>0$）：目标为 $\min(r\hat{A}, 1.2\hat{A})$。一旦 $r$ 超过 1.2，便不再有额外收益——防止 policy 在单个样本上过度承诺。对坏的动作（$\hat{A}<0$）：目标为 $\min(r\hat{A}, 0.8\hat{A})$。一旦 $r$ 降至 0.8 以下，惩罚便不再增长——防止灾难性遗忘（Catastrophic Forgetting）。

% \textbf{Key insight}: It’s a first-order approximation of TRPO’s KL constraint, but without expensive second-order optimization. Each update changes the policy by at most $\pm$20\%.
\textbf{关键洞见}：它是对 TRPO 的 KL 约束的一阶近似，但无需昂贵的二阶优化。每次更新对 policy 的改变最多 $\pm$20\%。

% \textbf{For LLMs specifically}: The token-level ratio $r_t = \pi_\theta(y_t|y_{<t})/\pi_\text{old}(y_t|y_{<t})$ prevents any single token’s probability from changing too drastically, preserving coherent generation.
\textbf{对 LLM 而言}：Token 级比率 $r_t = \pi_\theta(y_t|y_{<t})/\pi_\text{old}(y_t|y_{<t})$ 防止任何单个 Token 的概率发生过大变化，从而维持生成的连贯性。

% \emph{\textbf{Review:} Chapter~5 (PPO).}
\emph{\textbf{复习：} 第~5 章（PPO）。}
\end{questionbox}

% \begin{questionbox}[Q2: Derive DPO from first principles. What assumptions does it make?]
\begin{questionbox}[Q2：从第一原理推导 DPO。它做了哪些假设？]
% \textbf{Answer}: Start with RLHF objective: $\max_\pi \mathbb{E}[r(x,y)] - \beta D_\text{KL}[\pi\|\pi_\text{ref}]$.
\textbf{答}：从 RLHF 目标出发：$\max_\pi \mathbb{E}[r(x,y)] - \beta D_\text{KL}[\pi\|\pi_\text{ref}]$。

% \textbf{Step 1}: Write the KKT conditions. Optimal policy has closed form: $\pi^*(y|x) \propto \pi_\text{ref}(y|x)\exp(r(x,y)/\beta)$.
\textbf{步骤 1}：写出 KKT 条件。最优 policy 有闭式解：$\pi^*(y|x) \propto \pi_\text{ref}(y|x)\exp(r(x,y)/\beta)$。

% \textbf{Step 2}: Invert to express reward: $r(x,y) = \beta\log(\pi^*/\pi_\text{ref}) + \beta\log Z(x)$.
\textbf{步骤 2}：反解出 reward 的表达：$r(x,y) = \beta\log(\pi^*/\pi_\text{ref}) + \beta\log Z(x)$。

% \textbf{Step 3}: Substitute into Bradley-Terry model $P(y_w \succ y_l) = \sigma(r(y_w) - r(y_l))$. The partition function $Z(x)$ cancels (same prompt).
\textbf{步骤 3}：代入 Bradley-Terry 模型 $P(y_w \succ y_l) = \sigma(r(y_w) - r(y_l))$。配分函数 $Z(x)$ 抵消（同一 prompt）。

% \textbf{Step 4}: Replace $\pi^*$ with $\pi_\theta$ (parameterized policy we’re training): $\mathcal{L} = -\mathbb{E}[\log\sigma(\beta\log\frac{\pi_\theta(y_w)}{\pi_\text{ref}(y_w)} - \beta\log\frac{\pi_\theta(y_l)}{\pi_\text{ref}(y_l)})]$.
\textbf{步骤 4}：将 $\pi^*$ 替换为 $\pi_\theta$（我们要训练的参数化 policy）：$\mathcal{L} = -\mathbb{E}[\log\sigma(\beta\log\frac{\pi_\theta(y_w)}{\pi_\text{ref}(y_w)} - \beta\log\frac{\pi_\theta(y_l)}{\pi_\text{ref}(y_l)})]$。

% \textbf{Assumptions}:
\textbf{假设}：

\begin{enumerate}
  % \item Bradley-Terry preference model (pairwise, no ties, transitive)
  \item Bradley-Terry 偏好模型（成对、无平局、可传递）。
  % \item Optimal policy is achievable by $\pi_\theta$ (sufficient capacity)
  \item 最优 policy 可被 $\pi_\theta$ 实现（容量足够）。
  % \item Preferences are generated from the same distribution as training data (no distribution shift)
  \item 偏好与训练数据来自同一分布（无分布漂移）。
  % \item Reference model is fixed and reasonable
  \item 参考模型是固定且合理的。
\end{enumerate}

% \textbf{When assumptions break}: Real preferences aren’t transitive, data shifts during training, labels are noisy $\rightarrow$ that’s why Online DPO and IPO exist.
\textbf{当假设失效时}：真实偏好并不可传递，数据在训练中漂移，标签有噪声 $\rightarrow$ 这就是 Online DPO 和 IPO 存在的原因。

% \emph{\textbf{Review:} Chapter~6 (DPO).}
\emph{\textbf{复习：} 第~6 章（DPO）。}
\end{questionbox}

% \begin{questionbox}[Q3: GRPO vs PPO — when would you choose each? What’s the trade-off?]
\begin{questionbox}[Q3：GRPO vs PPO——你会在什么场景下选择哪一个？权衡何在？]
\textbf{答}：

% \textbf{GRPO advantages}:
\textbf{GRPO 的优势}：

\begin{itemize}
  % \item No value function needed: saves one model’s worth of memory and complexity
  \item 无需 value function：节省了一个模型规模的显存与复杂度。
  % \item Simpler: fewer hyperparameters, more intuitive (above-mean = good, below-mean = bad)
  \item 更简单：超参更少，更直观（高于均值 = 好，低于均值 = 坏）。
  % \item Better for verifiable rewards: math/code where $r \in \{0, 1\}$ gives crisp signal
  \item 更适合可验证奖励：数学/代码中 $r \in \{0, 1\}$ 给出清晰信号。
  % \item DeepSeek-R1 proved it can teach emergent reasoning with just binary rewards
  \item DeepSeek-R1 证明，仅用二元 reward 即可教出涌现的推理能力。
\end{itemize}

% \textbf{PPO advantages}:
\textbf{PPO 的优势}：

\begin{itemize}
  % \item Per-token credit assignment: value function assigns reward to each token, not just sequence-level
  \item 逐 Token 信用分配：value function 为每个 Token 分配 reward，而非仅在序列级。
  % \item More sample efficient: GAE uses value predictions to estimate advantage without generating $G$ samples
  \item 样本效率更高：GAE 用 value 预测估计 advantage，无需生成 $G$ 个样本。
  % \item Better for nuanced rewards: when reward is continuous and varies significantly across tokens
  \item 更适合细腻的 reward：当 reward 是连续的、在 Token 间显著变化时。
  % \item More mature: battle-tested at OpenAI, Anthropic, etc.
  \item 更成熟：在 OpenAI、Anthropic 等公司经过实战检验。
\end{itemize}

% \textbf{Rule of thumb}: If rewards are verifiable (right/wrong) $\rightarrow$ GRPO. If rewards are nuanced (RM scores) and you need max quality $\rightarrow$ PPO. If compute is limited $\rightarrow$ GRPO (no critic training).
\textbf{经验法则}：若 reward 可验证（对/错） $\rightarrow$ GRPO。若 reward 细腻（RM 分数）且追求极致质量 $\rightarrow$ PPO。若算力有限 $\rightarrow$ GRPO（无需训练 critic）。

% \textbf{Compute comparison}: GRPO generates $G$ responses per prompt (8$\times$ more generation), but skips value function training. Net: similar total compute but distributed differently (more gen, less training).
\textbf{算力对比}：GRPO 为每个 prompt 生成 $G$ 个回答（生成量增加 8$\times$），但跳过 value function 训练。净算力相近，但分布不同（生成更多，训练更少）。

% \emph{\textbf{Review:} Chapters~5 and~7 (PPO; GRPO).}
\emph{\textbf{复习：} 第~5 与~7 章（PPO；GRPO）。}
\end{questionbox}

% \begin{questionbox}[Q4: How does GAE work? Walk through a concrete example for LLMs.]
\begin{questionbox}[Q4：GAE 如何工作？为 LLM 走一个具体例子。]
% \textbf{Answer}: GAE = weighted sum of $n$-step TD errors: $\hat{A}_t = \sum_{l=0}^{T-t} (\gamma\lambda)^l \delta_{t+l}$.
\textbf{答}：GAE = $n$-步 TD 误差的加权和：$\hat{A}_t = \sum_{l=0}^{T-t} (\gamma\lambda)^l \delta_{t+l}$。

% \textbf{Concrete example}: Response has 5 tokens. Reward only at end ($r_5 = 0.8$). Value predictions: $V_1=0.5, V_2=0.55, V_3=0.6, V_4=0.65, V_5=0.7$.
\textbf{具体例子}：回答有 5 个 Token。仅在末端有 reward（$r_5 = 0.8$）。Value 预测：$V_1=0.5, V_2=0.55, V_3=0.6, V_4=0.65, V_5=0.7$。

% TD errors ($\gamma=1$): $\delta_1 = 0 + V_2 - V_1 = 0.05$, $\delta_2 = 0 + V_3 - V_2 = 0.05$, ..., $\delta_5 = 0.8 + 0 - 0.7 = 0.1$.
TD 误差（$\gamma=1$）：$\delta_1 = 0 + V_2 - V_1 = 0.05$，$\delta_2 = 0 + V_3 - V_2 = 0.05$，……，$\delta_5 = 0.8 + 0 - 0.7 = 0.1$。

% With $\lambda = 0.95$: $\hat{A}_5 = 0.1$ (just the final TD error), $\hat{A}_4 = 0.05 + 0.95 \times 0.1 = 0.145$, $\hat{A}_3 = 0.05 + 0.95 \times 0.145 = 0.188$, etc.
取 $\lambda = 0.95$：$\hat{A}_5 = 0.1$（仅是最后的 TD 误差），$\hat{A}_4 = 0.05 + 0.95 \times 0.1 = 0.145$，$\hat{A}_3 = 0.05 + 0.95 \times 0.145 = 0.188$，依此类推。

% \textbf{Interpretation}: Token 3 gets advantage 0.188 because it contributed to a sequence that got higher reward than expected. Earlier tokens get credit through the exponential decay.
\textbf{解读}：Token 3 获得 advantage 0.188，因为它对一个得到高于预期 reward 的序列做出了贡献。更早的 Token 通过指数衰减获得信用。

% \textbf{For LLMs}: $\gamma=1.0$ (all tokens matter, finite horizon). The advantage at token $t$ answers: ``given what followed this token, was this token choice better or worse than expected?''
\textbf{对 LLM}：$\gamma=1.0$（所有 Token 都重要，且为有限视野）。Token $t$ 的 advantage 回答：“鉴于此 Token 之后发生的事，这个 Token 的选择比预期更好还是更差？”

% \emph{\textbf{Review:} Chapter~5 (PPO).}
\emph{\textbf{复习：} 第~5 章（PPO）。}
\end{questionbox}

% \begin{questionbox}[Q5: How do you prevent reward hacking? Give a layered defense strategy.]
\begin{questionbox}[Q5：如何防止 reward hacking？给出一种分层防御策略。]
% \textbf{Answer}: \textbf{Detection signals}: RM score rising but win-rate flat/declining. Response length growing monotonically. KL divergence $>$ 15. Diversity (unique n-grams) dropping. Reading high-reward outputs reveals exploits.
\textbf{答}：\textbf{检测信号}：RM 分数上升但胜率持平或下降；回答长度单调增长；KL 散度 $>$ 15；多样性（unique n-grams）下降；阅读高 reward 输出可发现作弊模式。

% \textbf{Layered defense (in priority order)}:
\textbf{分层防御（按优先级排序）}：

\begin{enumerate}
  % \item \textbf{KL penalty} (primary): Adaptive controller targets KL $\approx$ 6. If KL rises, $\beta$ increases automatically. Prevents drifting too far from reference.
  \item \textbf{KL 惩罚}（首要）：自适应控制器以 KL $\approx$ 6 为目标。若 KL 上升，$\beta$ 自动增大。防止过度偏离参考模型。
  % \item \textbf{Reward model ensemble} (3--5 models): Use min or mean of scores. Individual models have different blind spots --- exploits that fool one rarely fool all.
  \item \textbf{Reward 模型集成}（3--5 个模型）：取分数的 min 或 mean。各模型有不同的盲区——能骗过一个的漏洞很少能骗过所有。
  % \item \textbf{Length penalty}: $r' = r - c \cdot \max(0, \text{length} - L_\text{target})$. Prevents ``just generate longer = higher score'' exploit.
  \item \textbf{长度惩罚}：$r' = r - c \cdot \max(0, \text{length} - L_\text{target})$。防止“只要更长就得高分”的作弊。
  % \item \textbf{Periodic RM refresh}: Every 2000 steps, generate data from current policy, relabel, add to RM training set. Closes the exploit as model finds it.
  \item \textbf{周期性 RM 刷新}：每 2000 步，从当前 policy 生成数据、重新标注，并加入 RM 训练集。模型一找到作弊路径就立即堵上。
  % \item \textbf{Win-rate based stopping}: Track win-rate against SFT baseline. If RM score rises but win-rate stalls for 200+ steps, stop training. The model is exploiting, not improving.
  \item \textbf{基于胜率的停训}：跟踪相对 SFT baseline 的胜率。如果 RM 分数上升但胜率停滞 200+ 步，立即停训。模型在作弊，而非进步。
\end{enumerate}

% \textbf{Post-detection recovery}: Roll back to last ``clean'' checkpoint. Increase $\beta$ by 2$\times$. Add the discovered exploit to the RM’s negative examples.
\textbf{检测后恢复}：回滚到最近一个“干净”的 checkpoint。将 $\beta$ 提升 2$\times$。把发现的作弊模式加入 RM 的负样本。

% \emph{\textbf{Review:} Chapters~9 and~11 (Reward Model Training; System Architecture).}
\emph{\textbf{复习：} 第~9 与~11 章（Reward 模型训练；系统架构）。}
\end{questionbox}

\section{系统设计题}
\label{system-design-questions}

% \begin{questionbox}[Q6: Design an RLHF system for training a 70B model. Walk through every component.]
\begin{questionbox}[Q6：为训练 70B 模型设计一个 RLHF 系统。逐组件讲解。]
% \textbf{Answer}: Three-cluster decoupled architecture on 72 A100-80GB GPUs:
\textbf{答}：基于 72 张 A100-80GB GPU 的三集群解耦架构：

% \textbf{Cluster 1 --- Generation (32 GPUs)}:
\textbf{集群 1 ——生成（32 GPU）}：

\begin{itemize}
  % \item 8 vLLM instances, TP=4 each. PagedAttention + speculative decoding (1B draft model).
  \item 8 个 vLLM 实例，每个 TP=4。PagedAttention + 投机解码（speculative decoding，1B 草稿模型）。
  % \item Continuous batching, max 256 sequences in flight. INT8 weights for bandwidth.
  \item 连续批处理，最多 256 条在飞序列。INT8 权重以省带宽。
  % \item Output: (prompt, response, per-token log-probs). Throughput: $\sim$500 responses/minute.
  \item 输出：(prompt, response, 逐 Token log-probs)。吞吐量：$\sim$500 回答/分钟。
  % \item Stateless: just loads latest weights from shared store. Trivial recovery.
  \item 无状态：仅从共享存储加载最新权重。恢复极简。
\end{itemize}

% \textbf{Cluster 2 --- Scoring (8 GPUs)}:
\textbf{集群 2 ——打分（8 GPU）}：

\begin{itemize}
  % \item Reward model (70B, INT8 = 70GB) on 4 GPUs (TP=4).
  \item Reward 模型（70B，INT8 = 70GB），4 GPU（TP=4）。
  % \item Reference model (70B, INT8) on 4 GPUs (TP=4). Computes per-token log-probs for KL.
  \item 参考模型（70B，INT8），4 GPU（TP=4）。计算逐 Token 的 log-probs 以算 KL。
  % \item Output: reward scores + KL per token. Lightweight, batch inference.
  \item 输出：reward 分数 + 逐 Token 的 KL。轻量级，批量推理。
\end{itemize}

% \textbf{Cluster 3 --- Training (32 GPUs)}:
\textbf{集群 3 ——训练（32 GPU）}：

\begin{itemize}
  % \item Policy model with FSDP (ZeRO-3). FlashAttention 2. Gradient checkpointing.
  \item Policy 模型用 FSDP（ZeRO-3）。FlashAttention 2。梯度检查点。
  % \item Consumes scored experiences from buffer. PPO update: 4 epochs on mini-batch of 16.
  \item 从 buffer 消费已打分的 experience。PPO 更新：在 mini-batch=16 上跑 4 个 epoch。
  % \item Pushes updated weights to shared store every 50 steps (async, background transfer).
  \item 每 50 步将更新后的权重推送到共享存储（异步、后台传输）。
  % \item Async checkpoint every 100 steps (non-blocking write to NVMe + S3 backup).
  \item 每 100 步做一次异步 checkpoint（非阻塞写入 NVMe + S3 备份）。
\end{itemize}

% \textbf{Connection fabric}:
\textbf{连接结构}：

\begin{itemize}
  % \item Gen $\rightarrow$ Score: Ray/Redis queue ($\sim$10 MB per batch: token IDs + log-probs)
  \item 生成 $\rightarrow$ 打分：Ray/Redis 队列（每 batch $\sim$10 MB：Token ID + log-probs）。
  % \item Score $\rightarrow$ Train: Experience buffer (circular, holds last 500 steps)
  \item 打分 $\rightarrow$ 训练：经验缓冲（循环式，保存最近 500 步）。
  % \item Train $\rightarrow$ Gen: Weight store on shared parallel FS (Lustre/GPFS). 140GB push every 50 steps, async.
  \item 训练 $\rightarrow$ 生成：共享并行文件系统（Lustre/GPFS）上的权重存储。每 50 步异步推送 140GB。
\end{itemize}

% \textbf{Overlap}: While training processes step $N$, generation is already producing data for step $N+1$. This hides generation latency, achieving 1.3--1.5$\times$ throughput vs monolithic.
\textbf{重叠}：当训练处理第 $N$ 步时，生成已经在为第 $N+1$ 步生产数据。这隐藏了生成延迟，相对单体架构吞吐提升 1.3--1.5$\times$。

% \emph{\textbf{Review:} Chapter~11 (System Architecture \& Infrastructure at Scale).}
\emph{\textbf{复习：} 第~11 章（系统架构与大规模基础设施）。}
\end{questionbox}

% \begin{questionbox}[Q7: How do you handle the generation bottleneck? Quantify the solutions.]
\begin{questionbox}[Q7：如何处理生成瓶颈？量化各种方案的收益。]
% \textbf{Answer}: Generation = 60--70\% of total RLHF wall-clock time. Root cause: autoregressive decoding is memory-bandwidth bound (arithmetic intensity $\approx 1$ FLOP/byte vs roofline of 156 FLOP/byte on A100).
\textbf{答}：生成占 RLHF 总墙钟时间的 60--70\%。根本原因：自回归解码受显存带宽约束（算术强度 $\approx 1$ FLOP/byte，而 A100 的 roofline 是 156 FLOP/byte）。

% \textbf{Solutions ranked by impact}:
\textbf{按影响力排序的方案}：

% \textbf{1. Decouple gen from training} (1.3--1.5$\times$ end-to-end): Run generation on separate hardware, overlap with training. The single biggest architectural win.
\textbf{1. 解耦生成与训练}（端到端 1.3--1.5$\times$）：在独立硬件上运行生成，与训练重叠。最大的单一架构性胜利。

% \textbf{2. vLLM with PagedAttention} (2--4$\times$): Eliminates 60--80\% KV cache memory waste from internal fragmentation. Enables 3--4$\times$ larger batches = better bandwidth utilization.
\textbf{2. vLLM + PagedAttention}（2--4$\times$）：消除内部碎片造成的 60--80\% KV cache 显存浪费。可承载 3--4$\times$ 更大的 batch = 更好的带宽利用。

% \textbf{3. Continuous batching} (1.5--2$\times$): Don’t wait for longest sequence to finish. Start new sequences immediately in freed slots. Keeps GPUs busy.
\textbf{3. 连续批处理}（1.5--2$\times$）：不等最长序列结束。立即在腾出的槽位开新序列。保持 GPU 繁忙。

% \textbf{4. Speculative decoding} (2--3$\times$): 1B draft model proposes 5 tokens. 70B model verifies all 5 in one forward pass (parallel!). Accept 3--4 on average $\rightarrow$ 3--4 tokens per forward pass instead of 1.
\textbf{4. 投机解码（Speculative Decoding）}（2--3$\times$）：1B 草稿模型提议 5 个 Token。70B 模型在一次前向中并行验证全部 5 个！平均接受 3--4 个 $\rightarrow$ 每次前向产出 3--4 个 Token 而非 1 个。

% \textbf{5. INT8/FP8 for generation weights} (2$\times$): Halves the 140GB weight read per token. Quality loss is minimal because (a) we’re sampling with temperature anyway, and (b) only generation uses INT8, training stays BF16.
\textbf{5. 生成权重用 INT8/FP8}（2$\times$）：将每 Token 的 140GB 权重读取减半。质量损失极小，因为：(a) 我们本就在用温度采样，(b) 仅生成使用 INT8，训练仍是 BF16。

% \textbf{6. CUDA graphs + kernel fusion} (1.1--1.3$\times$): Eliminate Python/CUDA launch overhead. Fuse layernorm+attention+MLP into fewer kernels.
\textbf{6. CUDA graphs + kernel fusion}（1.1--1.3$\times$）：消除 Python/CUDA 启动开销。把 layernorm+attention+MLP 融合成更少的 kernel。

% \textbf{Combined}: $1.5 \times 3 \times 1.5 \times 2.5 \times 2 \times 1.2 = 40\times$ over naive. In practice, diminishing returns limit total to $\sim$10--20$\times$ over naive implementation.
\textbf{组合后}：$1.5 \times 3 \times 1.5 \times 2.5 \times 2 \times 1.2 = 40\times$ 优于朴素实现。实际中由于边际收益递减，整体大约为朴素实现的 10--20$\times$。

% \emph{\textbf{Review:} Chapters~2 and~11 (Systems Foundations; System Architecture).}
\emph{\textbf{复习：} 第~2 与~11 章（系统基础；系统架构）。}
\end{questionbox}

% \begin{questionbox}[Q8: Explain weight synchronization in a decoupled system. How much staleness can you tolerate?]
\begin{questionbox}[Q8：解释解耦系统中的权重同步。能容忍多少陈旧度（staleness）？]
\textbf{答}：

% \textbf{The problem}: Generation cluster uses policy weights to produce responses. Training cluster updates those weights. They’re on different hardware. How to keep them in sync?
\textbf{问题}：生成集群使用 policy 权重生成回答，训练集群更新这些权重。它们位于不同硬件上。如何保持同步？

% \textbf{Why perfect sync is wasteful}: Full sync of 70B BF16 = 140GB. At InfiniBand 400Gb/s (50GB/s): 2.8s per sync. If you sync every step (every 50--90s), you spend 3--5\% of time on weight transfer. Acceptable, but unnecessary.
\textbf{为什么完美同步是浪费}：70B BF16 全量同步 = 140GB。InfiniBand 400Gb/s（50GB/s）下，每次同步 2.8s。若每步（每 50--90s）都同步，权重传输占总时间的 3--5\%。可接受，但没必要。

% \textbf{Staleness tolerance analysis}:
\textbf{陈旧度容忍分析}：

\begin{itemize}
  % \item Per-step policy change: $\sim$0.1\% (measured by mean param delta)
  \item 每步 policy 变化：$\sim$0.1\%（按参数均值差衡量）。
  % \item 50 steps: $\sim$5\% cumulative drift
  \item 50 步：$\sim$5\% 累计漂移。
  % \item PPO clip range: handles up to 20\% probability ratio deviation
  \item PPO 裁剪范围：可处理高达 20\% 的概率比偏离。
  % \item Empirical: 50-step staleness $\rightarrow$ $<$2\% quality degradation (measured by win-rate)
  \item 经验：50 步陈旧度 $\rightarrow$ 质量下降 $<$2\%（按胜率衡量）。
\end{itemize}

% \textbf{Production strategy}:
\textbf{生产策略}：

\begin{enumerate}
  % \item Every 50 training steps: push full BF16 checkpoint to shared store (2.8s transfer)
  \item 每 50 训练步：向共享存储推送一次 BF16 完整 checkpoint（传输 2.8s）。
  % \item Generation cluster: non-blocking weight reload between batches
  \item 生成集群：在 batch 之间做非阻塞的权重重载。
  % \item Delta compression (optional): Only send changed parameters (INT8 delta $\approx$ 5GB), apply as offset. 10$\times$ less bandwidth.
  \item 增量压缩（可选）：仅发送变更参数（INT8 delta $\approx$ 5GB），作为偏移应用。带宽减少 10$\times$。
  % \item For very large scale (256+ GPUs): streaming sync --- continuously send small chunks in background. Average staleness: 5--10 steps.
  \item 超大规模（256+ GPU）：流式同步——后台连续发送小块。平均陈旧度：5--10 步。
\end{enumerate}

% \textbf{Important subtlety}: Log-probs computed during generation use stale weights. The PPO ratio $\pi_\text{new}/\pi_\text{old}$ computes $\pi_\text{old}$ using these stale log-probs. This is fine because PPO is designed for off-policy corrections.
\textbf{关键细节}：生成阶段计算的 log-probs 使用的是陈旧权重。PPO 比率 $\pi_\text{new}/\pi_\text{old}$ 用这些陈旧 log-probs 作为 $\pi_\text{old}$。这没问题，因为 PPO 本就被设计为处理 off-policy 修正。

% \emph{\textbf{Review:} Chapter~11 (System Architecture \& Infrastructure at Scale).}
\emph{\textbf{复习：} 第~11 章（系统架构与大规模基础设施）。}
\end{questionbox}

% \begin{questionbox}[Q9: How would you make this fault-tolerant at 512 GPUs?]
\begin{questionbox}[Q9：在 512 GPU 规模下如何做到容错？]
% \textbf{Answer}: At 512 GPUs, MTBF is 4--8 hours. A 5-day training run will see 15--30 failures.
\textbf{答}：在 512 GPU 规模下，MTBF 为 4--8 小时。5 天的训练会经历 15--30 次故障。

% \textbf{Architecture-level resilience}:
\textbf{架构级韧性}：

\begin{itemize}
  % \item Generation cluster = stateless. Failed instance restarts in $<$60s (just load weights, no state).
  \item 生成集群 = 无状态。失败实例在 $<$60s 内重启（只需加载权重，无状态）。
  % \item Training cluster = stateful. Needs checkpoint-based recovery.
  \item 训练集群 = 有状态。需要基于 checkpoint 的恢复。
  % \item Scoring cluster = stateless. Same as generation.
  \item 打分集群 = 无状态。与生成相同。
\end{itemize}

% \textbf{Checkpointing strategy}:
\textbf{Checkpoint 策略}：

\begin{itemize}
  % \item Frequency: Every 50--100 steps (5--10 min of training).
  \item 频率：每 50--100 步（5--10 分钟训练）。
  % \item Method: Async (non-blocking). Background thread writes while next step proceeds. Uses FSDP’s distributed save (each rank saves its shard in parallel).
  \item 方法：异步（非阻塞）。后台线程在下一步进行时写入。使用 FSDP 的分布式保存（每个 rank 并行保存自己的分片）。
  % \item Contents: Model weights, optimizer states (Adam m/v), LR scheduler, RNG states, KL adaptive coefficient, global step counter, replay buffer pointer.
  \item 内容：模型权重、优化器状态（Adam m/v）、LR scheduler、RNG 状态、KL 自适应系数、全局 step 计数器、回放缓冲指针。
  % \item Storage: Local NVMe (fast, 30s for 70B) + async copy to S3/shared FS (durable).
  \item 存储：本地 NVMe（快，70B 仅 30s）+ 异步拷贝到 S3 / 共享 FS（持久化）。
  % \item Retention: Keep last 3 checkpoints. Auto-delete older ones.
  \item 保留：保留最近 3 个 checkpoint。自动删除更早的。
\end{itemize}

% \textbf{Detection and recovery flow}:
\textbf{检测与恢复流程}：

\begin{enumerate}
  % \item NCCL collective timeout (60s) or heartbeat miss (10s) $\rightarrow$ failure detected.
  \item NCCL 集合通信超时（60s）或心跳丢失（10s）$\rightarrow$ 检测到故障。
  % \item Identify failed node(s) via NVML health check.
  \item 通过 NVML 健康检查定位故障节点。
  % \item Option A (fast, $<$2 min): Torch Elastic shrinks world size, redistributes shards, continue with $N-1$ nodes. Request replacement in background.
  \item 方案 A（快，$<$2 min）：Torch Elastic 缩减 world size、重分配分片，以 $N-1$ 节点继续。后台请求替换。
  % \item Option B (clean, $\sim$5 min): Bring up replacement node, rebuild process group, load last checkpoint, resume.
  \item 方案 B（干净，$\sim$5 min）：拉起替换节点、重建进程组、加载最近 checkpoint、恢复。
  % \item Experience buffer is persisted --- no regeneration needed.
  \item 经验缓冲已持久化——无需重新生成。
\end{enumerate}

% \textbf{Prevention}: Pre-screening stress test (GEMM, memory, NVLink). ECC error monitoring (preemptive migration if errors spike). Hot spare nodes (pre-loaded with environment). Dual-rail InfiniBand for network redundancy.
\textbf{预防}：上线前的压力测试（GEMM、显存、NVLink）。ECC 错误监控（错误激增时抢占式迁移）。热备节点（环境预先加载）。双轨 InfiniBand 提供网络冗余。

% \emph{\textbf{Review:} Chapter~11 (System Architecture \& Infrastructure at Scale).}
\emph{\textbf{复习：} 第~11 章（系统架构与大规模基础设施）。}
\end{questionbox}

% \begin{questionbox}[Q10: How do you scale from 7B to 70B to 405B? What changes at each scale?]
\begin{questionbox}[Q10：如何从 7B 扩展到 70B 再到 405B？每个量级要变什么？]
\textbf{答}：

% \textbf{7B (single 8-GPU node, hours)}:
\textbf{7B（单 8-GPU 节点，数小时）}：

\begin{itemize}
  % \item Architecture: Monolithic (TRL default). All models on same GPUs.
  \item 架构：单体（TRL 默认）。所有模型在同一组 GPU 上。
  % \item Memory: LoRA + INT8 ref/RM fits in 8$\times$80GB easily.
  \item 显存：LoRA + INT8 的 ref/RM 在 8$\times$80GB 内轻松容纳。
  % \item Parallelism: DP=8 or FSDP within node. No network communication.
  \item 并行：节点内 DP=8 或 FSDP。无网络通信。
  % \item Hyperparams: LR=$5\times10^{-6}$, aggressive $\beta$=0.02, 50K--100K steps.
  \item 超参：LR=$5\times10^{-6}$，激进的 $\beta$=0.02，50K--100K 步。
  % \item Time: 4--12 hours per run. Fast iteration.
  \item 时间：每次运行 4--12 小时。快速迭代。
\end{itemize}

% \textbf{70B (32--64 GPUs, 2--5 days)}:
\textbf{70B（32--64 GPU，2--5 天）}：

\begin{itemize}
  % \item Architecture: Semi-decoupled. vLLM generation + FSDP training.
  \item 架构：半解耦。vLLM 生成 + FSDP 训练。
  % \item Memory: ZeRO-3 essential. Gradient checkpointing. INT8 ref/RM.
  \item 显存：ZeRO-3 必不可少。梯度检查点。INT8 ref/RM。
  % \item Parallelism: TP=8 intra-node (gen), FSDP inter-node (training).
  \item 并行：节点内 TP=8（生成），节点间 FSDP（训练）。
  % \item Hyperparams: LR=$1.5\times10^{-6}$, moderate $\beta$=0.05, 10K--30K steps.
  \item 超参：LR=$1.5\times10^{-6}$，温和的 $\beta$=0.05，10K--30K 步。
  % \item Fault tolerance: Async checkpoints, monitoring, but manageable manually.
  \item 容错：异步 checkpoint、监控，但仍可人工管理。
\end{itemize}

% \textbf{405B (256--512 GPUs, 1--3 weeks)}:
\textbf{405B（256--512 GPU，1--3 周）}：

\begin{itemize}
  % \item Architecture: Fully decoupled. Separate clusters. Weight store + queue.
  \item 架构：完全解耦。独立集群。权重存储 + 队列。
  % \item Memory: ZeRO-3 + TP=8 + PP=2 for training. INT4 generation.
  \item 显存：训练用 ZeRO-3 + TP=8 + PP=2。生成用 INT4。
  % \item Parallelism: 3D parallelism (TP$\times$PP$\times$DP = 8$\times$2$\times$16 = 256 GPUs training).
  \item 并行：3D 并行（TP$\times$PP$\times$DP = 8$\times$2$\times$16 = 256 GPU 训练）。
  % \item Hyperparams: LR=$5\times10^{-7}$, very conservative $\beta$=0.1, 2K--5K steps.
  \item 超参：LR=$5\times10^{-7}$，非常保守的 $\beta$=0.1，2K--5K 步。
  % \item Fault tolerance: Mandatory. Elastic training, redundant checkpoints, hot spares.
  \item 容错：必备。弹性训练、冗余 checkpoint、热备节点。
  % \item Key change: Much less RL training needed (model is already very capable from pretraining). But each step is 50$\times$ more expensive, so instability is catastrophic.
  \item 关键变化：所需 RL 训练量大幅减少（模型预训练后已非常强）。但每步贵 50$\times$，因此不稳定性是灾难性的。
\end{itemize}

% \textbf{Paradox}: Larger models are actually \emph{easier to RL-train per step} (more stable, smoother loss landscape). But the cost of instability scales with model size --- a bad run at 405B wastes \$100K+ in compute.
\textbf{悖论}：更大的模型实际上\emph{逐步 RL 训练更容易}（更稳定，loss 地形更平滑）。但不稳定性的代价随模型规模放大——405B 上一次失败的训练浪费 \$100K+ 的算力。

% \emph{\textbf{Review:} Chapter~11 (System Architecture \& Infrastructure at Scale).}
\emph{\textbf{复习：} 第~11 章（系统架构与大规模基础设施）。}
\end{questionbox}

\section{实战与调试题}
\label{practical-and-debugging-questions}

% \begin{questionbox}[Q11: Reward score is increasing but model quality is declining. Diagnose and fix.]
\begin{questionbox}[Q11：reward 分数上升但模型质量下降。诊断并修复。]
% \textbf{Answer}: Classic \textbf{reward hacking / Goodhart’s Law}.
\textbf{答}：典型的\textbf{reward hacking / 古德哈特定律（Goodhart's Law）}。

% \textbf{Diagnostic protocol}:
\textbf{诊断流程}：

\begin{enumerate}
  % \item \textbf{Check response length}: Plot mean length over training. Growing monotonically? = Length exploit (RM gives higher scores to longer responses).
  \item \textbf{检查回答长度}：绘制训练过程中的平均长度。单调增长？= 长度作弊（RM 给更长的回答更高分）。
  % \item \textbf{Check KL divergence}: $>$15? = Policy has diverged too far from reference. Lost capabilities.
  \item \textbf{检查 KL 散度}：$>$15？= policy 偏离参考模型过远，丧失了能力。
  % \item \textbf{Check diversity}: Unique trigrams per response. Dropping? = Mode collapse (repeating same high-reward pattern).
  \item \textbf{检查多样性}：每条回答的 unique trigrams。下降？= mode collapse（重复同一高 reward 模式）。
  % \item \textbf{Manual inspection}: Read 20 highest-reward responses. What pattern do they share? (e.g., all start with ``Great question!'', all use bullet points, excessive hedging).
  \item \textbf{人工检查}：阅读 20 个最高 reward 的回答。它们共享什么模式？（例如：都以"Great question!"开头、都用项目符号、过度模棱两可）。
  % \item \textbf{Win-rate}: Evaluate against SFT baseline on held-out prompts. If flat/declining while RM rises = confirmed exploit.
  \item \textbf{胜率}：在留出 prompt 上与 SFT baseline 比较。若 RM 上升而胜率持平/下降 = 确认存在作弊。
\end{enumerate}

% \textbf{Immediate fixes}:
\textbf{即时修复}：

\begin{itemize}
  % \item Roll back to last checkpoint where win-rate was improving
  \item 回滚到胜率还在提升的最近 checkpoint。
  % \item Increase $\beta$ by 2--3$\times$ (stronger KL penalty)
  \item 将 $\beta$ 提升 2--3$\times$（更强的 KL 惩罚）。
  % \item Add explicit length penalty: $r' = r - 0.001 \cdot \max(0, \text{len} - 500)$
  \item 添加显式长度惩罚：$r' = r - 0.001 \cdot \max(0, \text{len} - 500)$。
\end{itemize}

% \textbf{Structural fixes (prevent recurrence)}:
\textbf{结构性修复（防止复发）}：

\begin{itemize}
  % \item RM ensemble: Train 3--5 RMs on different data splits. Use min or mean. Exploits are model-specific.
  \item RM 集成：在不同数据切分上训练 3--5 个 RM。取 min 或 mean。作弊往往针对特定模型。
  % \item RM refresh: Every 2000 steps, generate from current policy, get human labels, retrain RM.
  \item RM 刷新：每 2000 步，从当前 policy 生成、获取人工标签、重训 RM。
  % \item Multi-objective reward: Combine helpfulness + harmlessness + conciseness with separate RMs.
  \item 多目标 reward：用独立 RM 组合“有用性 + 无害性 + 简洁性”。
  % \item Early stopping on win-rate (not RM score). The metric you optimize should differ from training signal.
  \item 基于胜率（而非 RM 分数）早停。你优化的指标应与训练信号不同。
\end{itemize}

% \emph{\textbf{Review:} Chapters~9 and~11 (Reward Model Training; System Architecture).}
\emph{\textbf{复习：} 第~9 与~11 章（Reward 模型训练；系统架构）。}
\end{questionbox}

% \begin{questionbox}[Q12: How do you decide the prompt distribution for RL training?]
\begin{questionbox}[Q12：如何决定 RL 训练的 prompt 分布？]
% \textbf{Answer}: Prompt quality is the most underrated factor in RLHF. Bad prompts = no learning signal.
\textbf{答}：prompt 质量是 RLHF 中最被低估的因素。差的 prompt = 没有学习信号。

% \textbf{Composition (my default mix)}:
\textbf{组成（我默认的配比）}：

\begin{itemize}
  % \item 40\% real user traffic (represents actual use cases)
  \item 40\% 真实用户流量（代表实际用例）。
  % \item 30\% synthetic (LLM-generated, fills gaps in coverage --- rare topics, edge cases)
  \item 30\% 合成（LLM 生成，填补覆盖空白——稀有主题、边角情况）。
  % \item 20\% curriculum (graduated difficulty --- start easy, increase complexity as model improves)
  \item 20\% 课程式（渐进难度——先易，随模型变强提升复杂度）。
  % \item 10\% adversarial (red-team prompts, jailbreak attempts, ambiguous instructions)
  \item 10\% 对抗式（red-team prompt、越狱尝试、含糊指令）。
\end{itemize}

% \textbf{Critical: The Goldilocks Filter}:
\textbf{关键：Goldilocks 过滤器}：

\begin{enumerate}
  % \item For each candidate prompt, generate 4--8 responses with current model.
  \item 对每个候选 prompt，用当前模型生成 4--8 个回答。
  % \item Score with RM. Compute pass rate (fraction above threshold).
  \item 用 RM 打分。计算通过率（超过阈值的比例）。
  % \item Keep only prompts with 20--80\% pass rate:
  \item 仅保留通过率在 20--80\% 的 prompt：

\begin{itemize}
  % \item $<$20\%: Too hard. Model almost always fails $\rightarrow$ all negative advantages $\rightarrow$ no useful gradient.
  \item $<$20\%：太难。模型几乎总失败 $\rightarrow$ 全是负 advantage $\rightarrow$ 无有效梯度。
  % \item $>$80\%: Too easy. Model almost always succeeds $\rightarrow$ all positive advantages $\rightarrow$ no contrast.
  \item $>$80\%：太易。模型几乎总成功 $\rightarrow$ 全是正 advantage $\rightarrow$ 无对比。
  % \item 20--80\%: Perfect. Mix of successes and failures $\rightarrow$ clear signal about what works.
  \item 20--80\%：完美。成功与失败兼有 $\rightarrow$ 关于“什么有效”的清晰信号。
\end{itemize}
  % \item Re-filter every 500 training steps (model improves, difficulty distribution shifts).
  \item 每 500 训练步重新过滤（模型在进步，难度分布在漂移）。
\end{enumerate}

% \textbf{Topic diversity}: Ensure no single topic dominates ($<$10\% per category). Use embedding clustering to verify coverage. Otherwise the model over-optimizes for dominant topics.
\textbf{主题多样性}：确保没有单一主题主导（每类 $<$10\%）。用 Embedding 聚类验证覆盖度。否则模型会过度优化主导主题。

% \emph{\textbf{Review:} Chapters~7 and~9 (GRPO; Reward Model Training).}
\emph{\textbf{复习：} 第~7 与~9 章（GRPO；Reward 模型训练）。}
\end{questionbox}

% \begin{questionbox}[Q13: LoRA vs full fine-tuning for RLHF. When would you use each?]
\begin{questionbox}[Q13：RLHF 中 LoRA vs 全量微调。各自何时使用？]
\textbf{答}：

% \textbf{LoRA} (Low-Rank Adaptation, $r$=64, $\alpha$=16):
\textbf{LoRA}（Low-Rank Adaptation，$r$=64，$\alpha$=16）：

\begin{itemize}
  % \item Trainable params: $\sim$0.2\% of model (200M for 70B)
  \item 可训练参数：模型的 $\sim$0.2\%（70B 对应 200M）。
  % \item Memory savings: No separate reference model needed (base model = reference)! Saves 140GB.
  \item 显存节省：无需独立参考模型（base 模型 = 参考）！节省 140GB。
  % \item Stability: Inherently more stable (low-rank constraint limits how far policy can drift)
  \item 稳定性：天然更稳定（低秩约束限制了 policy 的漂移幅度）。
  % \item Speed: Faster per-step (fewer params to update), but may need more steps
  \item 速度：每步更快（更新参数更少），但可能需要更多步。
  % \item Quality ceiling: 90--95\% of full FT quality typically
  \item 质量天花板：通常达到全量微调的 90--95\%。
\end{itemize}

% \textbf{Full fine-tuning}:
\textbf{全量微调}：

\begin{itemize}
  % \item All parameters updated. Maximum expressiveness.
  \item 所有参数都更新。表达力最强。
  % \item Needs separate reference model copy (140GB for 70B). Or very frequent checkpointing to ``anchor.''
  \item 需要独立的参考模型副本（70B 为 140GB）。或者通过非常频繁的 checkpoint 作为“锚”。
  % \item Higher risk of catastrophic forgetting. Need lower LR ($3\times$ lower than LoRA) and stronger $\beta$.
  \item 灾难性遗忘的风险更高。需要更低的 LR（比 LoRA 低 $3\times$）和更强的 $\beta$。
  % \item Better when: large distributional shift needed (new language, very different style), LoRA hits capacity limit.
  \item 更适合：需要大幅分布漂移（新语言、迥异风格），LoRA 触及容量上限的情况。
\end{itemize}

% \textbf{My decision framework}:
\textbf{我的决策框架}：

\begin{enumerate}
  % \item Start with LoRA ($r$=64). It’s 3$\times$ cheaper and more stable.
  \item 先用 LoRA（$r$=64）。便宜 3$\times$ 且更稳定。
  % \item Monitor gradient norms on LoRA matrices. If consistently $>$1.0 (high relative to parameter count): LoRA is capacity-limited.
  \item 监控 LoRA 矩阵的梯度范数。如果持续 $>$1.0（相对于参数量很大）：LoRA 已到容量上限。
  % \item Switch to full FT only when: win-rate plateaus AND gradient analysis suggests capacity limitation.
  \item 仅当胜率停滞 \emph{且} 梯度分析指示容量受限时，才切换到全量微调。
  % \item For full FT: use $\text{LR}/3$, $\beta \times 2$, more frequent checkpointing, and early stopping based on win-rate.
  \item 对全量微调：使用 $\text{LR}/3$、$\beta \times 2$、更频繁的 checkpoint，以及基于胜率的早停。
\end{enumerate}

% \textbf{Hybrid}: LoRA for alignment/safety (small behavioral shift) + Full FT for capabilities/reasoning (large shift needed).
\textbf{混合}：LoRA 用于对齐/安全（小行为漂移）+ 全量微调用于能力/推理（需要大漂移）。

% \emph{\textbf{Review:} Chapters~1 and~10 (LLM Architecture; SFT Best Practices).}
\emph{\textbf{复习：} 第~1 与~10 章（LLM 架构；SFT 最佳实践）。}
\end{questionbox}

% \begin{questionbox}[Q14: Process Reward Models (PRM) vs Outcome Reward Models (ORM). Design a PRM system.]
\begin{questionbox}[Q14：过程奖励模型（Process Reward Model, PRM）vs 结果奖励模型（Outcome Reward Model, ORM）。设计一个 PRM 系统。]
\textbf{答}：

% \textbf{ORM}: Scores the final answer only. ``Is the response good overall?'' Simple but can’t identify \emph{where} reasoning went wrong.
\textbf{ORM}：只给最终答案打分。“整体回答好不好？”简单但无法定位推理\emph{在何处}出错。

% \textbf{PRM}: Scores each intermediate step. ``Is step 3 of this derivation correct?'' Much more informative but harder to train.
\textbf{PRM}：为每个中间步骤打分。“这个推导的第 3 步对吗？”信息量大得多，但更难训练。

% \textbf{PRM advantages for reasoning}:
\textbf{PRM 对推理任务的优势}：

\begin{itemize}
  % \item Identifies exactly where reasoning fails (step-level credit assignment)
  \item 准确定位推理失败的位置（步骤级信用分配）。
  % \item Enables tree search: expand only branches with high step scores
  \item 支持树搜索：只展开步骤分数高的分支。
  % \item Less reward hacking: can’t get high score from wrong steps + lucky final answer
  \item reward hacking 更难：错误步骤 + 侥幸正确答案无法拿到高分。
  % \item PRM + best-of-N beats ORM + best-of-N by 10--20\% on MATH benchmarks
  \item 在 MATH 基准上，PRM + Best-of-N 比 ORM + Best-of-N 高 10--20\%。
\end{itemize}

% \textbf{Training a PRM}:
\textbf{训练 PRM}：

\begin{enumerate}
  % \item \textbf{Data collection} (Monte Carlo approach):
  \item \textbf{数据收集}（蒙特卡洛方法）：

\begin{itemize}
  % \item For each problem, generate reasoning trace step by step
  \item 对每道题，逐步生成推理轨迹。
  % \item At each step $k$, complete the solution $M$ times ($M=32$) from that point
  \item 在每一步 $k$，从该步出发把解完整补全 $M$ 次（$M=32$）。
  % \item Step score = fraction of completions that reach correct answer
  \item 步骤分数 = 抵达正确答案的补全比例。
  % \item Steps where score drops significantly = the ``mistake'' steps
  \item 分数明显下降的步骤 = “出错”的步骤。
\end{itemize}
  % \item \textbf{Labeling}: Step is ``correct'' if its completion rate $>$ 50\%, ``incorrect'' if $<$ 20\%.
  \item \textbf{标注}：补全率 $>$ 50\% 的步骤标为“正确”，$<$ 20\% 的标为“错误”。
  % \item \textbf{Model}: Same architecture as base model + classification head per token position. Train with binary cross-entropy on step labels.
  \item \textbf{模型}：与 base 模型同架构 + 每个 Token 位置上的分类头。用步骤标签做二元交叉熵训练。
  % \item \textbf{Inference}: Score each step. If any step scores $<$0.3, flag the trace as flawed.
  \item \textbf{推理}：为每步打分。若任一步骤分数 $<$0.3，则将该轨迹标记为有缺陷。
\end{enumerate}

% \textbf{Using PRM in RLHF}: Per-token rewards from PRM feed directly into GAE. Each token gets immediate feedback, not just end-of-sequence. This dramatically improves credit assignment for long reasoning chains.
\textbf{在 RLHF 中使用 PRM}：PRM 给出的逐 Token reward 直接喂入 GAE。每个 Token 都获得即时反馈，而非仅在序列末端。这极大改善了长推理链上的信用分配。

% \emph{\textbf{Review:} Chapters~9 and~13 (Reward Model Training; RL for Large Reasoning Models).}
\emph{\textbf{复习：} 第~9 与~13 章（Reward 模型训练；大型推理模型的 RL）。}
\end{questionbox}

% \begin{questionbox}[Q15: How do you evaluate whether RL actually improved the model?]
\begin{questionbox}[Q15：如何评估 RL 是否真的改进了模型？]
% \textbf{Answer}: Multi-faceted evaluation (no single metric captures ``quality''):
\textbf{答}：多面向评估（单一指标无法捕捉“质量”）：

% \textbf{1. Win-rate} (most important, most reliable):
\textbf{1. 胜率}（最重要、最可靠）：

\begin{itemize}
  % \item 500+ diverse prompts. LLM judge (GPT-4 or Claude) picks winner in blind A/B comparison vs SFT baseline.
  \item 500+ 多样 prompt。用 LLM judge（GPT-4 或 Claude）在与 SFT baseline 的盲对比 A/B 中选胜者。
  % \item Target: $>$55\% win-rate = meaningful improvement. $>$65\% = strong improvement.
  \item 目标：$>$55\% 胜率 = 有意义的改进；$>$65\% = 强改进。
  % \item Use position-debiasing (swap A/B order, average). Report confidence intervals.
  \item 使用位置去偏（互换 A/B 顺序、取平均）。报告置信区间。
\end{itemize}

% \textbf{2. Capability benchmarks} (regression detection):
\textbf{2. 能力基准}（回退检测）：

\begin{itemize}
  % \item MMLU (knowledge), HumanEval (code), MATH (reasoning), MT-Bench (multi-turn).
  \item MMLU（知识）、HumanEval（代码）、MATH（推理）、MT-Bench（多轮）。
  % \item Any $>$2\% drop = concerning alignment tax. Investigate which categories degraded.
  \item 任意 $>$2\% 的下降 = 值得警惕的对齐税。深入排查哪些类别退化。
\end{itemize}

% \textbf{3. Category-specific evals}:
\textbf{3. 类别专项评测}：

\begin{itemize}
  % \item Safety: refusal rate on harmful prompts (should increase)
  \item 安全：对有害 prompt 的拒绝率（应上升）。
  % \item Truthfulness: TruthfulQA score (should increase or stay flat)
  \item 真实性：TruthfulQA 分数（应上升或持平）。
  % \item Helpfulness: task completion rate on instruction-following benchmarks
  \item 有用性：指令跟随基准上的任务完成率。
\end{itemize}

% \textbf{4. Distributional metrics}:
\textbf{4. 分布性指标}：

\begin{itemize}
  % \item Response length distribution (shouldn’t shift dramatically)
  \item 回答长度分布（不应发生剧烈漂移）。
  % \item Vocabulary diversity (unique tokens per response)
  \item 词汇多样性（每条回答的 unique tokens）。
  % \item Format compliance (if trained for specific format)
  \item 格式合规（若训练目标包含特定格式）。
\end{itemize}

% \textbf{5. Human evaluation} (gold standard, expensive):
\textbf{5. 人工评测}（金标准，昂贵）：

\begin{itemize}
  % \item Blind A/B with 3+ skilled annotators per example. Inter-annotator agreement $>$ 70\%.
  \item 每个样本由 3+ 名熟练标注者做盲 A/B。标注者间一致性 $>$ 70\%。
  % \item Only for final model selection, not during training (too slow/expensive).
  \item 仅用于最终模型选择，训练过程中不用（太慢/太贵）。
\end{itemize}

% \textbf{Red flags}: RM score up + win-rate flat = reward hacking, not real improvement. Win-rate up + benchmarks down = alignment tax too high, reduce RL strength.
\textbf{红旗信号}：RM 分数上升 + 胜率持平 = reward hacking，不是真改进。胜率上升 + 基准下降 = 对齐税过高，需降低 RL 强度。

% \emph{\textbf{Review:} Chapter~14 (LLM Evaluation).}
\emph{\textbf{复习：} 第~14 章（LLM 评估）。}
\end{questionbox}

% \begin{questionbox}[Q16: Describe the reward model training pipeline end-to-end.]
\begin{questionbox}[Q16：端到端描述 reward 模型训练流水线。]
\textbf{答}：

% \textbf{Phase 1 --- Data Generation}:
\textbf{阶段 1——数据生成}：

\begin{itemize}
  % \item Collect 50K--100K diverse prompts (real traffic + synthetic)
  \item 收集 50K--100K 多样 prompt（真实流量 + 合成）。
  % \item Generate 4--8 responses per prompt at varying temperatures (0.3, 0.7, 1.0) and from multiple models (diversity is key --- if all responses are similar, preferences are uninformative)
  \item 对每个 prompt 用不同温度（0.3、0.7、1.0）和多个模型生成 4--8 个回答（多样性是关键——如果所有回答都相似，偏好就毫无信息量）。
  % \item Total: 200K--800K candidate responses
  \item 总计：200K--800K 个候选回答。
\end{itemize}

% \textbf{Phase 2 --- Preference Collection}:
\textbf{阶段 2——偏好收集}：

\begin{itemize}
  % \item Option A (expensive, best quality): Human annotators compare pairs. 3 annotators per pair. Cost: \$2--5 per comparison.
  \item 方案 A（昂贵，质量最佳）：人工标注者两两比较。每对 3 名标注者。成本：每次比较 \$2--5。
  % \item Option B (cheap, 85--90\% agreement with humans): LLM judge (GPT-4/Claude). 10$\times$ cheaper. Good for scale.
  \item 方案 B（便宜，与人类一致率 85--90\%）：LLM judge（GPT-4/Claude）。便宜 10$\times$。适合规模化。
  % \item Format: (prompt, chosen response, rejected response). Pairs with annotator disagreement ($<$70\% agreement) are discarded.
  \item 格式：(prompt, chosen response, rejected response)。标注者不一致（一致率 $<$70\%）的对被丢弃。
  % \item Final dataset: 100K--500K pairs.
  \item 最终数据集：100K--500K 对。
\end{itemize}

% \textbf{Phase 3 --- Training}:
\textbf{阶段 3——训练}：

\begin{itemize}
  % \item Architecture: Same as base LLM + scalar head (one regression output per sequence).
  \item 架构：与 base LLM 相同 + 标量头（每个序列一个回归输出）。
  % \item Loss: $\mathcal{L} = -\mathbb{E}[\log\sigma(r(x,y_w) - r(x,y_l))]$ (Bradley-Terry).
  \item Loss：$\mathcal{L} = -\mathbb{E}[\log\sigma(r(x,y_w) - r(x,y_l))]$（Bradley-Terry）。
  % \item Training: \textbf{1 epoch only!} RMs overfit extremely fast. Validation accuracy 68--75\% is good (higher often means overfitting to annotation artifacts).
  \item 训练：\textbf{仅 1 个 epoch！} RM 过拟合极快。验证准确率 68--75\% 即为良好（更高往往意味着过拟合到标注 artifact）。
  % \item Tricks: Center rewards around 0 (subtract running mean). Check for length bias (if correlation between length and score $>$ 0.3, add length penalty to training).
  \item 技巧：将 reward 中心化到 0（减去滑动均值）。检查长度偏差（若长度与分数相关 $>$ 0.3，在训练中加长度惩罚）。
\end{itemize}

% \textbf{Phase 4 --- Validation}:
\textbf{阶段 4——验证}：

\begin{itemize}
  % \item Hold-out preference pairs: accuracy should be 68--75\%.
  \item 留出偏好对：准确率应在 68--75\%。
  % \item Agreement with humans on new data: $>$ 80\%.
  \item 在新数据上与人类的一致率：$>$ 80\%。
  % \item Length bias check: correlation between response length and RM score should be $<$ 0.2.
  \item 长度偏差检查：回答长度与 RM 分数的相关 $<$ 0.2。
  % \item Consistency check: same prompt, paraphrased responses should get similar scores.
  \item 一致性检查：同一 prompt 下，改写后的回答应得到相近分数。
\end{itemize}

% \emph{\textbf{Review:} Chapter~9 (Reward Model Training).}
\emph{\textbf{复习：} 第~9 章（Reward 模型训练）。}
\end{questionbox}

% \begin{questionbox}[Q17: What happens when KL divergence explodes? Root cause and fix.]
\begin{questionbox}[Q17：KL 散度爆炸时会发生什么？根因与修复。]
\textbf{答}：

% \textbf{What KL measures}: Average log-ratio between current policy and reference: $D_\text{KL} = \mathbb{E}_{y\sim\pi_\theta}[\log(\pi_\theta(y|x)/\pi_\text{ref}(y|x))]$. KL=0 means identical to reference. KL=10 means the policy puts 10 nats more probability on its preferred outputs.
\textbf{KL 衡量什么}：当前 policy 与参考之间的平均对数比：$D_\text{KL} = \mathbb{E}_{y\sim\pi_\theta}[\log(\pi_\theta(y|x)/\pi_\text{ref}(y|x))]$。KL=0 意味着与参考完全相同。KL=10 意味着 policy 对其偏好的输出多投放了 10 nats 的概率。

% \textbf{Healthy range}: 3--10 during training. Slowly increasing is fine. Sudden spike = problem.
\textbf{健康区间}：训练中 3--10。缓慢增长无妨。突发飙升 = 出问题。

% \textbf{Root causes of KL explosion}:
\textbf{KL 爆炸的根因}：

\begin{enumerate}
  % \item \textbf{Learning rate too high}: Policy takes giant step, diverges from reference. Fix: reduce LR by 2--5$\times$.
  \item \textbf{学习率过高}：policy 步幅巨大、偏离参考。修复：将 LR 降低 2--5$\times$。
  % \item \textbf{Reward hacking}: Found an exploit that gets high reward far from reference behavior. Fix: increase $\beta$, add RM ensemble.
  \item \textbf{reward hacking}：找到了远离参考行为的高 reward 漏洞。修复：增大 $\beta$、加入 RM 集成。
  % \item \textbf{Mode collapse}: Policy concentrates on one response template. KL is high at that template, low everywhere else. Fix: increase entropy bonus, increase temperature.
  \item \textbf{Mode collapse}：policy 集中到一个回答模板上。在该模板处 KL 很高，其他地方都低。修复：加大熵奖励、提高温度。
  % \item \textbf{Bad batch}: One unlucky batch with extreme advantages pushed the policy. Fix: gradient clipping, reduce mini-batch size.
  \item \textbf{坏 batch}：一个具有极端 advantage 的不幸 batch 推动了 policy。修复：梯度裁剪、减小 mini-batch 大小。
  % \item \textbf{Value function diverged}: Wrong advantage estimates cause wrong updates. Fix: reduce value function LR, or switch to GRPO (no value function).
  \item \textbf{Value function 偏离}：错误的 advantage 估计导致错误的更新。修复：降低 value function 的 LR，或切换到 GRPO（无 value function）。
\end{enumerate}

% \textbf{Recovery protocol}:
\textbf{恢复流程}：

\begin{enumerate}
  % \item Detect: KL $>$ 15 for 50+ steps, or KL jumps $>$5 in one step.
  \item 检测：KL $>$ 15 持续 50+ 步，或单步内 KL 跳 $>$5。
  % \item Immediate: Load last clean checkpoint (KL $<$ 10).
  \item 立即：加载最近一个干净的 checkpoint（KL $<$ 10）。
  % \item Adjust: Reduce LR by 50\%. Increase $\beta$ by 2$\times$. Lower cliprange to 0.1.
  \item 调整：LR 降低 50\%。$\beta$ 增大 2$\times$。cliprange 降至 0.1。
  % \item Resume: Monitor closely for first 200 steps.
  \item 恢复：在头 200 步密切监控。
\end{enumerate}

% \emph{\textbf{Review:} Chapters~5 and~7 (PPO; GRPO).}
\emph{\textbf{复习：} 第~5 与~7 章（PPO；GRPO）。}
\end{questionbox}

% \begin{questionbox}[Q18: Compare monolithic vs decoupled RLHF architectures. When does each make sense?]
\begin{questionbox}[Q18：比较单体 vs 解耦的 RLHF 架构。各自在什么场景合理？]
\textbf{答}：

% \textbf{Monolithic} (TRL default: single process, all models on same GPUs):
\textbf{单体}（TRL 默认：单进程，所有模型同 GPU）：

\begin{itemize}
  % \item Pros: Simple code. No distributed systems complexity. Easy debugging.
  \item 优点：代码简单。无分布式系统复杂度。易于调试。
  % \item Cons: GPUs idle 60\% of time (compute idle during gen, bandwidth idle during training). Doesn’t scale past $\sim$16 GPUs efficiently. All models compete for same memory.
  \item 缺点：GPU 闲置 60\% 时间（生成时计算闲、训练时带宽闲）。无法高效扩展到 $\sim$16 GPU 以上。所有模型争抢同一显存。
  % \item Use when: Model $\leq$ 13B, single node, research/prototyping.
  \item 适用：模型 $\leq$ 13B，单节点，研究/原型。
\end{itemize}

% \textbf{Semi-decoupled} (vLLM gen + FSDP training, same cluster):
\textbf{半解耦}（vLLM 生成 + FSDP 训练，同一集群）：

\begin{itemize}
  % \item Pros: Better utilization (gen and training can partially overlap). Scales to 64 GPUs.
  \item 优点：利用率更好（生成与训练可部分重叠）。可扩展到 64 GPU。
  % \item Cons: Still shares hardware, can’t optimize independently. More complex than monolithic.
  \item 缺点：仍共享硬件，无法独立优化。比单体更复杂。
  % \item Use when: 13B--70B, 2--8 nodes, production experiments.
  \item 适用：13B--70B，2--8 节点，生产实验。
\end{itemize}

% \textbf{Fully decoupled} (separate clusters connected by queues):
\textbf{完全解耦}（独立集群通过队列连接）：

\begin{itemize}
  % \item Pros: Each cluster optimized for its workload. Gen scales independently from training. Gen cluster is stateless (trivial fault tolerance). Scales to hundreds of GPUs.
  \item 优点：每个集群针对其工作负载优化。生成与训练独立扩展。生成集群无状态（容错极简）。可扩展到数百 GPU。
  % \item Cons: Distributed systems complexity. Weight staleness. Queue management. Network overhead.
  \item 缺点：分布式系统复杂度。权重陈旧度。队列管理。网络开销。
  % \item Use when: $\geq$ 70B production training. Need scale, fault tolerance, high utilization.
  \item 适用：$\geq$ 70B 生产训练。需要规模、容错、高利用率。
\end{itemize}

% \textbf{The key insight}: Generation is bandwidth-bound, training is compute-bound. Same hardware can’t optimize for both. Decoupling lets gen nodes have: more memory bandwidth, INT8 weights, large batch. Training nodes have: full BF16 precision, FlashAttention, FSDP sharding.
\textbf{关键洞见}：生成受带宽约束，训练受算力约束。同一硬件无法同时为二者优化。解耦让生成节点拥有：更多显存带宽、INT8 权重、大 batch。训练节点拥有：完整 BF16 精度、FlashAttention、FSDP 分片。

% \emph{\textbf{Review:} Chapter~11 (System Architecture \& Infrastructure at Scale).}
\emph{\textbf{复习：} 第~11 章（系统架构与大规模基础设施）。}
\end{questionbox}

% \begin{questionbox}[Q19: How do you set up curriculum learning for RL training?]
\begin{questionbox}[Q19：如何为 RL 训练设置课程学习？]
% \textbf{Answer}: Curriculum = gradually increasing difficulty so the model learns progressively.
\textbf{答}：课程学习 = 逐渐提升难度，让模型循序渐进地学习。

% \textbf{Why it matters}: If you throw the hardest prompts at a weak model, it gets all-negative rewards $\rightarrow$ no learning signal (everything is equally bad). If you start easy, the model develops basic capabilities, then builds on them.
\textbf{为什么重要}：若把最难的 prompt 丢给弱模型，它会得到全负的 reward $\rightarrow$ 无学习信号（一切都同样糟）。若从易入手，模型先发展基础能力，再在其上构建。

% \textbf{Implementation}:
\textbf{实现}：

\begin{enumerate}
  % \item \textbf{Difficulty scoring}: Rate each prompt by current model’s pass rate (from Goldilocks filtering). Easy = $>$80\% pass, Medium = 30--80\%, Hard = $<$30\%.
  \item \textbf{难度评分}：根据当前模型的通过率（来自 Goldilocks 过滤）对每个 prompt 打分。易 = 通过 $>$80\%，中 = 30--80\%，难 = $<$30\%。
  % \item \textbf{Schedule}: Steps 0--1000: 70\% easy, 20\% medium, 10\% hard. Steps 1000--5000: 30\% easy, 50\% medium, 20\% hard. Steps 5000+: 10\% easy, 40\% medium, 50\% hard.
  \item \textbf{排程}：步 0--1000：70\% 易，20\% 中，10\% 难。步 1000--5000：30\% 易，50\% 中，20\% 难。步 5000+：10\% 易，40\% 中，50\% 难。
  % \item \textbf{Dynamic adjustment}: Every 500 steps, re-evaluate difficulty distribution. Prompts the model has ``mastered'' (pass rate $>$ 95\%) get retired. New harder prompts are introduced.
  \item \textbf{动态调整}：每 500 步重新评估难度分布。模型已“掌握”（通过率 $>$ 95\%）的 prompt 退役。引入更难的新 prompt。
  % \item \textbf{For GRPO specifically}: Curriculum ensures groups always have a mix of successes and failures. Without curriculum, hard prompts give all-zero groups (useless).
  \item \textbf{对 GRPO}：课程保证 group 总是混合成功与失败。没有课程，难 prompt 会产生全零 group（无用）。
\end{enumerate}

% \textbf{Evidence}: DeepSeek-R1 used implicit curriculum --- starting with easy math/code problems, the model developed basic reasoning, then solved progressively harder problems without explicit scheduling.
\textbf{证据}：DeepSeek-R1 使用隐式课程——从简单的数学/代码问题起步，模型发展出基本推理，再逐步解出更难的问题，无需显式排程。

% \emph{\textbf{Review:} Chapters~7 and~12 (GRPO; LLM Agentic Training).}
\emph{\textbf{复习：} 第~7 与~12 章（GRPO；LLM 智能体训练）。}
\end{questionbox}

% \begin{questionbox}[Q20: You have a budget of 64 A100-80GB GPUs and need to RL-train a 70B model. Design the allocation.]
\begin{questionbox}[Q20：你有 64 张 A100-80GB GPU 的预算，需要 RL 训练一个 70B 模型。设计资源分配。]
% \textbf{Answer}: 8 nodes $\times$ 8 GPUs = 64 total. Need to split across generation, scoring, and training.
\textbf{答}：8 节点 $\times$ 8 GPU = 64 张总。需在生成、打分、训练之间分配。

% \textbf{My allocation}:
\textbf{我的分配}：

\begin{itemize}
  % \item \textbf{Generation}: 24 GPUs (3 nodes). 6 vLLM instances with TP=4. INT8 weights = 70GB/model, leaves room for KV cache. Continuous batching, batch $\approx$ 128 total.
  \item \textbf{生成}：24 GPU（3 节点）。6 个 vLLM 实例，TP=4。INT8 权重 = 70GB/模型，留下空间给 KV cache。连续批处理，总 batch $\approx$ 128。
  % \item \textbf{Scoring}: 8 GPUs (1 node). RM (INT8, TP=4) + Reference model (INT8, TP=4) on same node. Or share 4 GPUs with TP=4 alternating between RM and ref.
  \item \textbf{打分}：8 GPU（1 节点）。RM（INT8，TP=4）+ 参考模型（INT8，TP=4）位于同一节点。或共享 4 GPU，TP=4 在 RM 与 ref 间轮换。
  % \item \textbf{Training}: 32 GPUs (4 nodes). FSDP across all 32. Each GPU holds $\sim$70B/32 = 2.2GB params + optimizer fraction. Plenty of headroom for activations. Gradient checkpointing for safety.
  \item \textbf{训练}：32 GPU（4 节点）。在全部 32 张上做 FSDP。每张 GPU 持有 $\sim$70B/32 = 2.2GB 参数 + 优化器片段。激活留有充足余量。开启梯度检查点以保险。
\end{itemize}

% \textbf{Expected throughput}:
\textbf{预期吞吐}：

\begin{itemize}
  % \item Generation: 6 instances $\times$ $\sim$80 responses/min = 480 responses/min
  \item 生成：6 实例 $\times$ $\sim$80 回答/分钟 = 480 回答/分钟。
  % \item Training: Batch of 128, one step every $\sim$15s (training only, no gen wait)
  \item 训练：batch=128，每 $\sim$15s 一步（仅训练时间，不含等待生成）。
  % \item Overlap: While training step $N$ (15s), generation produces $\sim$120 responses for step $N+1$. Perfect pipeline.
  \item 重叠：当训练第 $N$ 步进行（15s）时，生成已为第 $N+1$ 步产出 $\sim$120 个回答。完美流水。
\end{itemize}

% \textbf{Bottleneck analysis}: Generation takes $\sim$45s for 128 responses. Training takes $\sim$15s. Scoring takes $\sim$5s. Generation is bottleneck. Could move 8 GPUs from training to gen (40 gen, 24 training) to balance, but then training is bottleneck. Current allocation is near-optimal.
\textbf{瓶颈分析}：128 个回答的生成需 $\sim$45s。训练 $\sim$15s。打分 $\sim$5s。生成是瓶颈。可将 8 GPU 从训练移到生成（40 生成、24 训练）以平衡，但那样训练变成瓶颈。当前分配接近最优。

% \textbf{Alternative if memory-tight}: Move scoring onto training nodes (time-share: score during gen, train during training). Saves 8 GPUs. Slightly worse pipelining but works.
\textbf{显存吃紧时的替代方案}：把打分放到训练节点（分时复用：生成时打分，训练时训练）。省下 8 GPU。流水稍差但可用。

% \emph{\textbf{Review:} Chapters~2 and~11 (Systems Foundations; System Architecture).}
\emph{\textbf{复习：} 第~2 与~11 章（系统基础；系统架构）。}
\end{questionbox}

\section{GRPO 变体与高级 RL 题}
\label{grpo-variants-and-advanced-rl-questions}

% \begin{questionbox}[Q21: What is DAPO and how does it improve over standard GRPO?]
\begin{questionbox}[Q21：什么是 DAPO？它相对标准 GRPO 有哪些改进？]
% \textbf{Answer}: DAPO (Dynamic Adaptive Policy Optimization) introduces 5 key modifications:
\textbf{答}：DAPO（动态自适应策略优化，Dynamic Adaptive Policy Optimization）引入了 5 项关键改动：

% \textbf{1. Clip-Higher (asymmetric clipping)}: Standard PPO/GRPO clips both directions equally at $\epsilon=0.2$. DAPO uses $\epsilon_\text{low}=0.2$ but $\epsilon_\text{high}=0.28$. This allows the model to \emph{increase} good action probabilities more aggressively while still restricting how much it suppresses bad ones. Intuition: exploration needs more room than exploitation.
\textbf{1. Clip-Higher（非对称裁剪）}：标准 PPO/GRPO 两侧都按 $\epsilon=0.2$ 等量裁剪。DAPO 使用 $\epsilon_\text{low}=0.2$ 但 $\epsilon_\text{high}=0.28$。这允许模型\emph{更激进地提升}好动作的概率，同时仍限制对坏动作的抑制幅度。直觉：探索比利用需要更大的空间。

% \textbf{2. Overlong Filtering}: If a response hits the max length limit (truncated, no EOS token), it’s masked entirely from the loss. Rationale: truncated responses contain no natural stopping signal --- training on them teaches the model that ``stopping mid-sentence'' is acceptable.
\textbf{2. 超长过滤（Overlong Filtering）}：若回答达到最大长度上限（被截断，没有 EOS Token），就将其完全从 loss 中屏蔽。理由：截断的回答不包含自然停止信号——在其上训练会教会模型“句子中间停下”是可接受的。

% \textbf{3. Token-level Loss}: Loss is normalized by total token count across all sequences, not by number of sequences. This prevents longer sequences from dominating the gradient.
\textbf{3. Token 级 Loss}：loss 按所有序列的总 Token 数归一，而非按序列数。这防止更长的序列主导梯度。

% \textbf{4. Soft Overlong Punishment}: Instead of binary truncation filtering, apply a gradual penalty as responses approach max length. $r_\text{soft} = -c \cdot \max(0, \text{len} - L_\text{soft})/(L_\text{max} - L_\text{soft})$.
\textbf{4. 软超长惩罚（Soft Overlong Punishment）}：不再是二元截断过滤，而是在回答接近最大长度时施加渐进惩罚。$r_\text{soft} = -c \cdot \max(0, \text{len} - L_\text{soft})/(L_\text{max} - L_\text{soft})$。

% \textbf{5. Dynamic Sampling}: Resample prompts during training to ensure each batch has a mix of success/failure (not yet in TRL).
\textbf{5. 动态采样（Dynamic Sampling）}：训练中重采样 prompt，确保每个 batch 都有成功/失败的混合（TRL 中尚未实现）。

% \textbf{When to use}: Large-scale reasoning RL where you need maximum exploration and long completions (32K+ tokens). The asymmetric clipping is particularly valuable.
\textbf{适用场景}：需要极大探索与超长补全（32K+ Token）的大规模推理 RL。非对称裁剪尤其有价值。

% \emph{\textbf{Review:} Chapters~7 and~8 (GRPO; Preference Optimization Variants).}
\emph{\textbf{复习：} 第~7 与~8 章（GRPO；偏好优化变体）。}
\end{questionbox}

% \begin{questionbox}[Q22: Explain the vLLM train-inference mismatch. Why does it happen and how do TIS/MIS fix it?]
\begin{questionbox}[Q22：解释 vLLM 的训练-推理不一致。为何发生，TIS/MIS 如何修复？]
% \textbf{Answer}: \textbf{The problem}: When using vLLM for generation and a training framework (DeepSpeed/FSDP) for updates, the same model with the same weights produces \emph{different} token probabilities. This happens because:
\textbf{答}：\textbf{问题}：当用 vLLM 做生成、用训练框架（DeepSpeed/FSDP）做更新时，同一模型、同一权重会产生\emph{不同}的 Token 概率。原因：

\begin{itemize}
  % \item Different numerical kernels (vLLM uses custom CUDA kernels optimized for throughput)
  \item 数值 kernel 不同（vLLM 使用针对吞吐优化的自定义 CUDA kernel）。
  % \item Different attention implementations (FlashAttention in training vs PagedAttention in vLLM)
  \item attention 实现不同（训练用 FlashAttention，vLLM 用 PagedAttention）。
  % \item Different precision handling (FP8/INT8 in vLLM vs BF16 in training)
  \item 精度处理不同（vLLM 用 FP8/INT8，训练用 BF16）。
  % \item Batching differences affecting layer normalization numerics
  \item 批处理差异影响 layer normalization 的数值。
\end{itemize}

% This silently breaks PPO’s on-policy assumption: we compute the ratio $\pi_\theta/\pi_\text{old}$ using $\pi_\text{old}$ from vLLM but $\pi_\theta$ from the training framework. The ratio is wrong from step zero!
这会悄无声息地破坏 PPO 的 on-policy 假设：我们计算 $\pi_\theta/\pi_\text{old}$ 时，$\pi_\text{old}$ 来自 vLLM 而 $\pi_\theta$ 来自训练框架。从第零步起比率就是错的！

% \textbf{TIS (Truncated Importance Sampling)}: Correct the gradient by multiplying by $\min(\pi_\text{train}/\pi_\text{inference}, C)$. The $\min$ with cap $C$ prevents extreme corrections from destabilizing training. Typical $C=2.0$.
\textbf{TIS（截断重要性采样，Truncated Importance Sampling）}：通过乘以 $\min(\pi_\text{train}/\pi_\text{inference}, C)$ 修正梯度。$\min$ 与上限 $C$ 防止极端修正破坏训练稳定性。典型 $C=2.0$。

% \textbf{MIS (Masked Importance Sampling)}: More aggressive --- simply discard any token where $\pi_\text{train}/\pi_\text{inference} > C$. Zero contribution to gradient. Prevents any badly-estimated token from affecting the update.
\textbf{MIS（掩码重要性采样，Masked Importance Sampling）}：更激进——直接丢弃任何 $\pi_\text{train}/\pi_\text{inference} > C$ 的 Token，对梯度贡献为零。防止任何估计糟糕的 Token 影响更新。

% \textbf{Sequence vs Token level}: Sequence-level IS is theoretically correct (unbiased); token-level IS is biased but lower variance. In practice, sequence-level with truncation works best.
\textbf{序列级 vs Token 级}：序列级 IS 理论上正确（无偏）；Token 级 IS 有偏但方差更低。实践中，序列级带截断效果最好。

% \emph{\textbf{Review:} Chapters~7 and~11 (GRPO; System Architecture).}
\emph{\textbf{复习：} 第~7 与~11 章（GRPO；系统架构）。}
\end{questionbox}

% \begin{questionbox}[Q23: GSPO vs GRPO — what’s the fundamental difference and when does it matter?]
\begin{questionbox}[Q23：GSPO vs GRPO——根本区别在哪？什么时候重要？]
% \textbf{Answer}: \textbf{GRPO}: Computes importance ratio \emph{per token}: $w_{i,t} = \pi_\theta(o_{i,t}|q, o_{i,<t}) / \pi_\text{old}(o_{i,t}|q, o_{i,<t})$, then clips each token independently.
\textbf{答}：\textbf{GRPO}：\emph{逐 Token}计算重要性比 $w_{i,t} = \pi_\theta(o_{i,t}|q, o_{i,<t}) / \pi_\text{old}(o_{i,t}|q, o_{i,<t})$，然后独立裁剪每个 Token。

% \textbf{GSPO}: Computes importance ratio at the \emph{sequence level}: $s_i(\theta) = (\pi_\theta(o_i|q)/\pi_\text{old}(o_i|q))^{1/|o_i|}$ --- the geometric mean of token probabilities. Clips this single sequence-level ratio.
\textbf{GSPO}：在\emph{序列级}计算重要性比：$s_i(\theta) = (\pi_\theta(o_i|q)/\pi_\text{old}(o_i|q))^{1/|o_i|}$——Token 概率的几何均值。裁剪这单一的序列级比率。

% \textbf{Why it matters}: GRPO’s per-token clipping treats each token as independent, but in language they’re deeply correlated. A small per-token change early in the sequence compounds exponentially over many tokens. GSPO captures this by looking at the full sequence probability.
\textbf{为何重要}：GRPO 的逐 Token 裁剪把每个 Token 视为独立，但语言中它们高度相关。序列前段的微小逐 Token 变化在多个 Token 上指数级放大。GSPO 通过审视完整序列概率来捕捉这一点。

% \textbf{Length normalization}: The $1/|o_i|$ exponent ensures fair comparison across different-length sequences. Without it, longer sequences would always have lower probability ratios.
\textbf{长度归一化}：$1/|o_i|$ 指数保证不同长度序列间的公平比较。否则更长的序列总是有更低的概率比。

% \textbf{When to use GSPO}: When training goes off-policy (\texttt{steps\_per\_generation > 1} or \texttt{num\_iterations > 1}). If fully on-policy (ratio $\approx 1$), GRPO and GSPO are equivalent.
\textbf{何时用 GSPO}：当训练变为 off-policy 时（\texttt{steps\_per\_generation > 1} 或 \texttt{num\_iterations > 1}）。如果完全 on-policy（比率 $\approx 1$），GRPO 与 GSPO 等价。

% \emph{\textbf{Review:} Chapters~7 and~8 (GRPO; Preference Optimization Variants).}
\emph{\textbf{复习：} 第~7 与~8 章（GRPO；偏好优化变体）。}
\end{questionbox}

% \begin{questionbox}[Q24: The paper ``It Takes Two'' shows G=2 matches G=16. How is that possible?]
\begin{questionbox}[Q24：论文 ``It Takes Two'' 表明 G=2 与 G=16 效果相当。这怎么可能？]
% \textbf{Answer}: The key insight is that GRPO’s effectiveness doesn’t come from accurate advantage estimation (which would need large $G$), but from an \textbf{implicit contrastive objective}.
\textbf{答}：关键洞见是 GRPO 的有效性并非来自精确的 advantage 估计（那需要大 $G$），而来自一个\textbf{隐式的对比目标}。

% With $G=2$ and binary rewards (one correct, one wrong): $\hat{A}_\text{correct} = +1$, $\hat{A}_\text{wrong} = -1$ (after normalization). The loss becomes: increase probability of correct response, decrease probability of wrong response. This is essentially a DPO-style contrastive loss!
$G=2$ 加二元 reward（一个对一个错）时：归一化后 $\hat{A}_\text{correct} = +1$，$\hat{A}_\text{wrong} = -1$。loss 变成：提升正确回答的概率、降低错误回答的概率。这本质上就是一个 DPO 风格的对比 loss！

% \textbf{Why large $G$ doesn’t help much}: The normalized advantage $\hat{A}_i = (r_i - \mu)/\sigma$ already creates a contrast between good and bad. More samples give a better $\mu$ estimate, but the gradient direction is dominated by the \emph{contrast} between best and worst, not the mean accuracy.
\textbf{为什么大 $G$ 帮助不大}：归一化 advantage $\hat{A}_i = (r_i - \mu)/\sigma$ 本身已制造了好坏之间的对比。更多样本能更准确估计 $\mu$，但梯度方向由最好与最差之间的\emph{对比}主导，而非 $\mu$ 的精度。

% \textbf{Compute savings}: $G=2$ means 8$\times$ less generation compute than $G=16$. Since generation is 60\% of training time, this translates to $\sim$4$\times$ faster training.
\textbf{算力节省}：$G=2$ 意味着生成算力比 $G=16$ 少 8$\times$。由于生成占训练时间 60\%，整体训练加速 $\sim$4$\times$。

% \textbf{Caveat}: Works best when pass rate is 30--70\%. If pass rate is very low ($<$10\%), $G=2$ often gives two failures (no signal). Need larger $G$ for hard problems.
\textbf{注意}：通过率在 30--70\% 时效果最佳。如果通过率极低（$<$10\%），$G=2$ 常常给出两个失败（无信号）。难题需要更大的 $G$。

% \emph{\textbf{Review:} Chapter~7 (GRPO).}
\emph{\textbf{复习：} 第~7 章（GRPO）。}
\end{questionbox}

% \begin{questionbox}[Q25: What is SAPO and how does its soft gating differ from hard clipping?]
\begin{questionbox}[Q25：什么是 SAPO？它的软门控与硬裁剪有何不同？]
% \textbf{Answer}: Standard PPO/GRPO uses hard clipping: $\text{clip}(r, 1-\epsilon, 1+\epsilon)$. At the boundary, the gradient suddenly drops to zero. This creates a ``dead zone'' where the model receives no learning signal.
\textbf{答}：标准 PPO/GRPO 使用硬裁剪：$\text{clip}(r, 1-\epsilon, 1+\epsilon)$。在边界处梯度突然降为零。这制造了一个“死区”，模型在其中收不到任何学习信号。

% \textbf{SAPO} replaces this with a smooth sigmoid gate: the gradient is gradually attenuated as the ratio moves away from 1, never suddenly zeroed out. It uses asymmetric temperatures:
\textbf{SAPO} 用平滑的 sigmoid 门控取代它：随着比率偏离 1，梯度被逐渐衰减，绝不会突兀归零。它使用非对称温度：

\begin{itemize}
  % \item $\tau_+ = 1.0$ for positive advantages (standard attenuation)
  \item 对正 advantage 用 $\tau_+ = 1.0$（标准衰减）。
  % \item $\tau_- = 1.05$ for negative advantages (slightly more aggressive attenuation for suppression)
  \item 对负 advantage 用 $\tau_- = 1.05$（对抑制略激进一些的衰减）。
\end{itemize}

% \textbf{Benefits}: (1) No ``cliff'' in gradient landscape. (2) Tokens slightly outside the clip range still contribute (attenuated, not zeroed). (3) More stable optimization trajectory. (4) Sequence-coherent --- considers the full sequence context.
\textbf{收益}：(1) 梯度地形上没有“悬崖”。(2) 略超出裁剪范围的 Token 仍有贡献（衰减而非归零）。(3) 优化轨迹更稳定。(4) 序列连贯——考虑了完整序列上下文。

% \textbf{Trade-off}: Slightly less restrictive trust region than hard clipping, so requires careful temperature tuning. But more robust to hyperparameter choices overall.
\textbf{权衡}：信任域比硬裁剪略宽松，因此需要仔细调温度。但整体对超参更鲁棒。

% \emph{\textbf{Review:} Chapters~7 and~8 (GRPO; Preference Optimization Variants).}
\emph{\textbf{复习：} 第~7 与~8 章（GRPO；偏好优化变体）。}
\end{questionbox}

\section{DPO 扩展题}
\label{dpo-extensions-questions}

% \begin{questionbox}[Q26: Compare f-DPO divergence choices. When would you use forward KL vs JS vs reverse KL?]
\begin{questionbox}[Q26：比较 f-DPO 的散度选择。前向 KL、JS 与反向 KL 各自何时使用？]
% \textbf{Answer}: Standard DPO uses reverse KL implicitly ($D_\text{KL}[\pi_\theta \| \pi_\text{ref}]$):
\textbf{答}：标准 DPO 隐式使用反向 KL（$D_\text{KL}[\pi_\theta \| \pi_\text{ref}]$）：

\begin{itemize}
  % \item \textbf{Reverse KL} (default): Mode-seeking. The policy concentrates probability where the reference is high. Avoids generating text the reference wouldn’t. Good for safety (conservative).
  \item \textbf{反向 KL}（默认）：寻峰（mode-seeking）。policy 把概率集中在参考概率高的位置。避免生成参考不会生成的文本。利于安全（保守）。
  % \item \textbf{Forward KL}: Mass-covering. The policy tries to cover all modes of the reference, even low-probability ones. Good for diversity but can generate low-quality outputs.
  \item \textbf{前向 KL}：覆盖（mass-covering）。policy 试图覆盖参考的所有峰，甚至低概率的峰。利于多样性但可能生成低质量输出。
  % \item \textbf{Jensen-Shannon}: Symmetric compromise between forward and reverse. Balanced mode-coverage and mode-seeking. Often best for general alignment.
  \item \textbf{Jensen-Shannon}：前向与反向之间的对称折中。在覆盖与寻峰间取得平衡。对通用对齐通常最佳。
  % \item \textbf{Alpha-divergence} ($\alpha=0.5$): Interpolates between forward ($\alpha=0$) and reverse ($\alpha=1$). Tunable.
  \item \textbf{Alpha 散度}（$\alpha=0.5$）：在前向（$\alpha=0$）与反向（$\alpha=1$）之间插值。可调。
\end{itemize}

% \textbf{Practical recommendation}: Start with reverse KL (standard DPO). If the model is too conservative (won’t try creative solutions), switch to JS divergence. If diversity is critical (creative writing, brainstorming), try forward KL.
\textbf{实践建议}：先用反向 KL（标准 DPO）。若模型过于保守（不愿尝试创造性方案），切换到 JS 散度。若多样性至关重要（创意写作、头脑风暴），试试前向 KL。

% \emph{\textbf{Review:} Chapters~6 and~8 (DPO; Preference Optimization Variants).}
\emph{\textbf{复习：} 第~6 与~8 章（DPO；偏好优化变体）。}
\end{questionbox}

% \begin{questionbox}[Q27: Your DPO preference data has 15\% label noise. What do you do?]
\begin{questionbox}[Q27：你的 DPO 偏好数据有 15\% 的标签噪声。怎么办？]
% \textbf{Answer}: Three options in order of sophistication:
\textbf{答}：按复杂度递增的三个方案：

% \textbf{1. Robust DPO} (best for known noise rate): Analytically debiases the loss: $\mathcal{L}_\text{robust} = \frac{(1-\varepsilon)\mathcal{L}_\text{DPO}(y_w, y_l) - \varepsilon \mathcal{L}_\text{DPO}(y_l, y_w)}{1 - 2\varepsilon}$. Set $\varepsilon = 0.15$. This provably recovers the clean DPO objective in expectation. TRL: \texttt{loss\_type="robust", label\_smoothing=0.15}.
\textbf{1. Robust DPO}（已知噪声率时最佳）：解析地去偏 loss：$\mathcal{L}_\text{robust} = \frac{(1-\varepsilon)\mathcal{L}_\text{DPO}(y_w, y_l) - \varepsilon \mathcal{L}_\text{DPO}(y_l, y_w)}{1 - 2\varepsilon}$。设 $\varepsilon = 0.15$。在期望意义上可证明恢复干净的 DPO 目标。TRL：\texttt{loss\_type="robust", label\_smoothing=0.15}。

% \textbf{2. IPO} (best when noise rate unknown): Squared loss with target margin. Mislabeled pairs have bounded influence (the squared loss doesn’t diverge). More robust to arbitrary noise patterns without needing to know $\varepsilon$. TRL: \texttt{loss\_type="ipo"}.
\textbf{2. IPO（Identity Preference Optimization, IPO）}（噪声率未知时最佳）：带目标 margin 的平方 loss。被误标的对影响有界（平方 loss 不发散）。对任意噪声模式更鲁棒，且无需知道 $\varepsilon$。TRL：\texttt{loss\_type="ipo"}。

% \textbf{3. TR-DPO} (best for distribution shift): Updates the reference model via EMA during training. Even if early data is noisy, the evolving reference helps the model self-correct. TRL: \texttt{sync\_ref\_model=True, ref\_model\_mixup\_alpha=0.6}.
\textbf{3. TR-DPO}（分布漂移时最佳）：训练中通过 EMA 更新参考模型。即便早期数据有噪声，演化的参考也能帮助模型自我修正。TRL：\texttt{sync\_ref\_model=True, ref\_model\_mixup\_alpha=0.6}。

% \textbf{Data-side fixes}: (1) Filter pairs with $<$70\% inter-annotator agreement. (2) Use RM to score pairs; discard those where RM disagrees with label. (3) Active learning: re-label the most uncertain pairs.
\textbf{数据侧修复}：(1) 过滤标注者一致率 $<$70\% 的对。(2) 用 RM 给对打分；丢弃 RM 与标签不一致的对。(3) 主动学习：重新标注最不确定的对。

% \emph{\textbf{Review:} Chapters~6 and~8 (DPO; Preference Optimization Variants).}
\emph{\textbf{复习：} 第~6 与~8 章（DPO；偏好优化变体）。}
\end{questionbox}

% \begin{questionbox}[Q28: What is SimPO and why is being reference-free advantageous?]
\begin{questionbox}[Q28：什么是 SimPO？为什么“无参考模型”是优势？]
% \textbf{Answer}: SimPO uses the average log-probability of a response as an implicit reward signal: $r(x,y) = \frac{1}{|y|}\sum_t \log \pi_\theta(y_t|x, y_{<t})$ --- no reference model needed.
\textbf{答}：SimPO 用回答的平均对数概率作为隐式 reward 信号：$r(x,y) = \frac{1}{|y|}\sum_t \log \pi_\theta(y_t|x, y_{<t})$——无需参考模型。

% The loss adds a target margin $\gamma$: the chosen response should have average log-prob at least $\gamma$ higher than rejected.
loss 中加入目标 margin $\gamma$：chosen 回答的平均 log-prob 应至少比 rejected 高 $\gamma$。

% \textbf{Why reference-free matters}:
\textbf{为什么“无参考”重要}：

\begin{enumerate}
  % \item \textbf{Memory}: No reference model = 70--140GB saved for 70B models. Can train larger models on same hardware.
  \item \textbf{显存}：无参考模型 = 70B 节省 70--140GB。可在同样硬件上训练更大的模型。
  % \item \textbf{Simplicity}: No need to manage/load/serve a second model copy.
  \item \textbf{简洁性}：无需管理/加载/服务第二份模型副本。
  % \item \textbf{No stale reference}: DPO’s reference becomes increasingly irrelevant as training progresses. SimPO doesn’t have this problem.
  \item \textbf{无陈旧参考}：DPO 的参考随着训练推进越来越不相关。SimPO 没有这个问题。
  % \item \textbf{Length normalization built in}: The $1/|y|$ naturally prevents length bias (DPO needs explicit handling).
  \item \textbf{内置长度归一化}：$1/|y|$ 天然防止长度偏差（DPO 需显式处理）。
\end{enumerate}

% \textbf{Trade-off}: Without a reference anchor, the model has more freedom to collapse or drift. The $\gamma$ margin and length normalization partially mitigate this, but SimPO can be less stable than DPO for aggressive training.
\textbf{权衡}：没有参考锚点，模型有更多自由去塌缩或漂移。$\gamma$ margin 和长度归一化部分缓解了这一点，但在激进训练时 SimPO 可能不如 DPO 稳定。

% \emph{\textbf{Review:} Chapter~8 (Preference Optimization Variants).}
\emph{\textbf{复习：} 第~8 章（偏好优化变体）。}
\end{questionbox}

% \begin{questionbox}[Q29: Explain Iterative RPO. Why combine DPO with NLL loss for reasoning?]
\begin{questionbox}[Q29：解释 Iterative RPO。为何在推理任务中把 DPO 与 NLL loss 结合？]
% \textbf{Answer}: Standard DPO for reasoning has a subtle failure mode: it learns to \emph{discriminate} (assign higher implicit reward to correct traces) but doesn’t necessarily learn to \emph{generate} them.
\textbf{答}：用于推理的标准 DPO 有一种微妙的失效模式：它学会\emph{判别}（给正确轨迹更高的隐式 reward），但不一定学会\emph{生成}它们。

% \textbf{Why}: DPO’s gradient pushes chosen probability up and rejected probability down. But the chosen response might be so different from what the model would generate that increasing its probability doesn’t teach the model how to produce similar reasoning patterns.
\textbf{为什么}：DPO 的梯度推高 chosen 的概率、压低 rejected 的概率。但 chosen 回答可能与模型自身会生成的东西差异巨大，以至于提升其概率并不能教会模型产出类似的推理模式。

% \textbf{RPO’s fix}: Add a negative log-likelihood (NLL/SFT) loss on the chosen response: $\mathcal{L} = \mathcal{L}_\text{DPO} + \alpha \cdot \mathcal{L}_\text{NLL}(y_w)$.
\textbf{RPO 的修复}：在 chosen 回答上添加负对数似然（NLL/SFT）loss：$\mathcal{L} = \mathcal{L}_\text{DPO} + \alpha \cdot \mathcal{L}_\text{NLL}(y_w)$。

% The NLL term explicitly trains the model to generate the winning response step by step. The DPO term ensures it also learns to avoid the losing response. Combined: the model learns both ``how to reason correctly'' (NLL) and ``what to avoid'' (DPO).
NLL 项显式训练模型逐步生成获胜回答。DPO 项确保模型同时学会避免落败回答。结合起来：模型同时学到“如何正确推理”（NLL）和“该避免什么”（DPO）。

% \textbf{Iterative}: Generate responses $\rightarrow$ check correctness $\rightarrow$ create pairs $\rightarrow$ train with RPO $\rightarrow$ repeat. Each iteration the model gets better at generating correct reasoning, creating higher-quality training data for the next round.
\textbf{迭代}：生成回答 $\rightarrow$ 检查正确性 $\rightarrow$ 构造对 $\rightarrow$ 用 RPO 训练 $\rightarrow$ 重复。每次迭代，模型更擅长生成正确推理，为下一轮提供更高质量的训练数据。

% TRL: \texttt{loss\_type=["sigmoid", "sft"], loss\_weights=[1.0, 1.0]}
TRL：\texttt{loss\_type=["sigmoid", "sft"], loss\_weights=[1.0, 1.0]}

% \emph{\textbf{Review:} Chapters~8 and~13 (Preference Optimization Variants; RL for Large Reasoning Models).}
\emph{\textbf{复习：} 第~8 与~13 章（偏好优化变体；大型推理模型的 RL）。}
\end{questionbox}

\section{GPU 架构与硬件题}
\label{gpu-architecture-and-hardware-questions}

% \begin{questionbox}[Q30: Explain the GPU memory hierarchy. Why does it matter for LLM inference?]
\begin{questionbox}[Q30：解释 GPU 内存层次结构。它对 LLM 推理为何重要？]
% \textbf{Answer}: From fastest to slowest:
\textbf{答}：从最快到最慢：

\begin{enumerate}
  % \item \textbf{Registers}: Per-thread, $\sim$256 KB/SM. Instant access (0 cycles latency).
  \item \textbf{寄存器（Registers）}：线程私有，$\sim$256 KB/SM。即时访问（0 周期延迟）。
  % \item \textbf{SRAM (Shared Memory)}: Per-SM, $\sim$192--228 KB/SM (A100: 164 KB configurable). Bandwidth: $\sim$19 TB/s aggregate. Latency: $\sim$20 cycles.
  \item \textbf{SRAM（共享内存）}：每 SM 私有，$\sim$192--228 KB/SM（A100：164 KB 可配置）。聚合带宽：$\sim$19 TB/s。延迟：$\sim$20 周期。
  % \item \textbf{L2 Cache}: Shared across GPU, 40--60 MB (H100: 50 MB). Bandwidth: $\sim$5 TB/s. Latency: $\sim$200 cycles.
  \item \textbf{L2 缓存}：跨整张 GPU 共享，40--60 MB（H100：50 MB）。带宽：$\sim$5 TB/s。延迟：$\sim$200 周期。
  % \item \textbf{HBM}: Main GPU memory, 80 GB (A100). Bandwidth: 2--3.35 TB/s. Latency: $\sim$400 cycles.
  \item \textbf{HBM}：GPU 主显存，80 GB（A100）。带宽：2--3.35 TB/s。延迟：$\sim$400 周期。
  % \item \textbf{CPU DRAM}: Via PCIe, 512 GB+. Bandwidth: 32--64 GB/s. Latency: $\sim$10K cycles.
  \item \textbf{CPU DRAM}：通过 PCIe，512 GB+。带宽：32--64 GB/s。延迟：$\sim$10K 周期。
\end{enumerate}

% \textbf{Why it matters for LLMs}: Autoregressive generation reads the full model weights ($\sim$140 GB for 70B) for every single token. At 2 TB/s HBM bandwidth, that’s 70ms just to stream the weights. The actual computation (one matrix-vector multiply) takes only 0.5ms. The GPU is 99\% waiting for data.
\textbf{对 LLM 为何重要}：自回归生成在每个 Token 都要读取完整模型权重（70B 约 $\sim$140 GB）。以 2 TB/s 的 HBM 带宽，仅流式读权重就要 70ms。实际计算（一次矩阵-向量乘法）只要 0.5ms。GPU 有 99\% 的时间在等数据。

% \textbf{FlashAttention exploits this}: By keeping intermediate results (QK scores, softmax) in SRAM (19 TB/s) instead of writing them to HBM (2 TB/s), it eliminates 90\% of the memory traffic for attention. The compute is the same, but HBM reads/writes drop 10$\times$.
\textbf{FlashAttention 利用了这一点}：通过把中间结果（QK 分数、softmax）保留在 SRAM（19 TB/s）而非写到 HBM（2 TB/s），消除了 attention 的 90\% 内存流量。计算量未变，但 HBM 的读写降低 10$\times$。

% \emph{\textbf{Review:} Chapter~2 (Systems Foundations for LLMs).}
\emph{\textbf{复习：} 第~2 章（LLM 的系统基础）。}
\end{questionbox}

% \begin{questionbox}[Q31: How does FlashAttention work? What is the online softmax trick?]
\begin{questionbox}[Q31：FlashAttention 如何工作？什么是 online softmax 技巧？]
% \textbf{Answer}: \textbf{Problem}: Standard attention materializes the $n \times n$ attention matrix in HBM. For $n=8192$: $8192^2 \times 2 = 134$ MB per head, 4.3 GB per layer with 32 heads. Must write to HBM then read back for softmax and multiply --- 3 full HBM round-trips.
\textbf{答}：\textbf{问题}：标准 attention 在 HBM 中物化 $n \times n$ 的 attention 矩阵。$n=8192$ 时：每个 head $8192^2 \times 2 = 134$ MB，32 头每层 4.3 GB。必须写到 HBM 再读回做 softmax 与乘法——3 次完整的 HBM 往返。

% \textbf{FlashAttention solution}: Never store the full $n \times n$ matrix. Process in tiles that fit in SRAM.
\textbf{FlashAttention 方案}：从不存储完整的 $n \times n$ 矩阵。以可放进 SRAM 的分块（tile）处理。

% \textbf{Algorithm}:
\textbf{算法}：

\begin{enumerate}
  % \item Split $Q$ into blocks of size $B_r$ rows, $K/V$ into blocks of $B_c$ rows.
  \item 将 $Q$ 拆为 $B_r$ 行的块，$K/V$ 拆为 $B_c$ 行的块。
  % \item For each $Q$ block: iterate over all $K$ blocks, computing partial attention scores.
  \item 对每个 $Q$ 块：遍历所有 $K$ 块，计算部分 attention 分数。
  % \item \textbf{Online softmax trick}: Maintain running max $m$ and running sum $\ell$ for softmax normalization. When processing a new $K$ block, update: $m_\text{new} = \max(m_\text{old}, \max(\text{scores}))$, rescale previous accumulator by $e^{m_\text{old} - m_\text{new}}$, add new contribution.
  \item \textbf{Online softmax 技巧}：维护滑动 max $m$ 与滑动 sum $\ell$ 以做 softmax 归一。处理新的 $K$ 块时更新：$m_\text{new} = \max(m_\text{old}, \max(\text{scores}))$，将上一累加器按 $e^{m_\text{old} - m_\text{new}}$ 重新缩放，再加上新的贡献。
  % \item Output is accumulated incrementally --- never needs the full $n \times n$ matrix.
  \item 输出以增量方式累加——从不需要完整的 $n \times n$ 矩阵。
\end{enumerate}

% \textbf{Key insight}: Softmax is normally a global operation ($\max$ and $\sum$ over all elements). The online trick decomposes it into local updates with a correction factor. Mathematically exact --- no approximation.
\textbf{关键洞见}：softmax 本来是全局操作（对所有元素取 $\max$ 和 $\sum$）。online 技巧将其分解为带修正因子的局部更新。数学上精确——并非近似。

% \textbf{Result}: Memory $O(n)$ instead of $O(n^2)$. Speed 2--4$\times$ faster (fewer HBM accesses, more time in SRAM).
\textbf{结果}：显存 $O(n)$ 而非 $O(n^2)$。速度提升 2--4$\times$（更少 HBM 访问，更多时间在 SRAM）。

% \textbf{FlashAttention 2}: Better work partitioning across warps, reduces non-matmul FLOPs by 2$\times$.
\textbf{FlashAttention 2}：在 warp 之间做更好的工作划分，将非矩阵乘的 FLOPs 减半。

% \textbf{FlashAttention 3} (H100/Hopper): Uses Tensor Memory Accelerator (TMA) for async loads, warp specialization (producer/consumer warps), FP8 support.
\textbf{FlashAttention 3}（H100/Hopper）：使用张量内存加速器（Tensor Memory Accelerator, TMA）做异步加载、warp 专门化（生产者/消费者 warp）、支持 FP8。

% \emph{\textbf{Review:} Chapters~1 and~2 (LLM Architecture; Systems Foundations).}
\emph{\textbf{复习：} 第~1 与~2 章（LLM 架构；系统基础）。}
\end{questionbox}

% \begin{questionbox}[Q32: Explain PagedAttention. How does it solve the KV cache problem?]
\begin{questionbox}[Q32：解释 PagedAttention。它如何解决 KV cache 问题？]
% \textbf{Answer}: \textbf{The problem}: During generation, each sequence needs a KV cache (stores K and V tensors for all previous tokens). For a 70B model: each token needs $2 \times n_\text{layers} \times d_\text{model} \times 2$ bytes = $2 \times 80 \times 8192 \times 2 \approx 2.5$ MB per token. A 2048-token sequence: $\sim$5 GB of KV cache.
\textbf{答}：\textbf{问题}：生成期间，每条序列都需要 KV cache（存储所有先前 Token 的 K、V 张量）。对 70B 模型：每个 Token 需要 $2 \times n_\text{layers} \times d_\text{model} \times 2$ 字节 = $2 \times 80 \times 8192 \times 2 \approx 2.5$ MB。2048 Token 的序列：$\sim$5 GB 的 KV cache。

% \textbf{Traditional allocation}: Pre-allocate max\_sequence\_length for each active sequence. If max=2048 but average=500, you waste 75\% of allocated memory. With 50 concurrent sequences, that’s hundreds of GB wasted.
\textbf{传统分配}：为每条活跃序列预分配 max\_sequence\_length。若 max=2048 而平均=500，就浪费 75\% 已分配显存。50 条并发序列就浪费数百 GB。

% \textbf{PagedAttention}: Inspired by OS virtual memory:
\textbf{PagedAttention}：受操作系统虚拟内存启发：

\begin{enumerate}
  % \item KV cache is split into fixed-size \emph{blocks} (pages), each holding KV for 16 tokens.
  \item KV cache 被切分为固定大小的\emph{块}（页），每块保存 16 个 Token 的 KV。
  % \item A \emph{block table} (like a page table) maps logical token positions to physical memory blocks.
  \item 一张\emph{块表}（类似页表）将逻辑 Token 位置映射到物理显存块。
  % \item Blocks are allocated on demand as the sequence grows. No pre-allocation waste.
  \item 块按需随序列增长分配。无预分配浪费。
  % \item Freed blocks return to a pool immediately when sequences finish.
  \item 序列结束时，释放的块立即归还到池中。
\end{enumerate}

% \textbf{Extra benefits}:
\textbf{额外收益}：

\begin{itemize}
  % \item \textbf{Prefix sharing}: Multiple sequences with the same system prompt share KV cache blocks (copy-on-write). Saves 30--50\% memory for chat applications.
  \item \textbf{前缀共享}：拥有相同系统 prompt 的多条序列共享 KV cache 块（写时复制）。聊天应用可省 30--50\% 显存。
  % \item \textbf{Preemption}: Can ``swap out'' a low-priority sequence’s blocks to CPU, freeing GPU memory for higher-priority requests.
  \item \textbf{抢占}：可将低优先级序列的块“换出”到 CPU，为高优先级请求腾出 GPU 显存。
  % \item \textbf{Near-zero fragmentation}: Internal fragmentation limited to last block ($<$16 tokens). External fragmentation eliminated (any free block can be used anywhere).
  \item \textbf{接近零碎片}：内部碎片仅限末块（$<$16 Token）。外部碎片被消除（任意空闲块可用于任意位置）。
\end{itemize}

% \textbf{Result}: 3--5$\times$ more concurrent sequences in the same memory $\rightarrow$ 3--5$\times$ better throughput for serving.
\textbf{结果}：同等显存下并发序列数提升 3--5$\times$ $\rightarrow$ 服务吞吐提升 3--5$\times$。

% \emph{\textbf{Review:} Chapter~2 (Systems Foundations for LLMs).}
\emph{\textbf{复习：} 第~2 章（LLM 的系统基础）。}
\end{questionbox}

% \begin{questionbox}[Q33: Compare NVLink vs InfiniBand. When do you use each in RLHF training?]
\begin{questionbox}[Q33：比较 NVLink 与 InfiniBand。在 RLHF 训练中各用于何处？]
\textbf{答}：

% \textbf{NVLink} (intra-node, GPU-to-GPU):
\textbf{NVLink}（节点内，GPU 到 GPU）：

\begin{itemize}
  % \item Bandwidth: 600 GB/s (A100), 900 GB/s (H100) --- total bidirectional
  \item 带宽：600 GB/s（A100）、900 GB/s（H100）——双向总和。
  % \item Latency: $\sim$1 $\mu$s
  \item 延迟：$\sim$1 $\mu$s。
  % \item Scope: Within one physical node (8 GPUs connected via NVSwitch)
  \item 范围：单个物理节点内（8 GPU 通过 NVSwitch 互联）。
  % \item Use case: \textbf{Tensor Parallelism} (TP=8). Each layer’s matrix multiply is split across GPUs, requiring AllReduce after every layer. Needs ultra-high bandwidth + low latency.
  \item 用例：\textbf{张量并行}（TP=8）。每层矩阵乘法在 GPU 间切分，每层后都要 AllReduce。需要超高带宽 + 低延迟。
\end{itemize}

% \textbf{InfiniBand NDR} (inter-node, node-to-node):
\textbf{InfiniBand NDR}（跨节点，节点到节点）：

\begin{itemize}
  % \item Bandwidth: 400 Gb/s = 50 GB/s per port. With 8 ports (GPUDirect RDMA): 400 GB/s aggregate per node.
  \item 带宽：每端口 400 Gb/s = 50 GB/s。8 端口（GPUDirect RDMA）：每节点聚合 400 GB/s。
  % \item Latency: $\sim$1--5 $\mu$s (RDMA)
  \item 延迟：$\sim$1--5 $\mu$s（RDMA）。
  % \item Scope: Between nodes in a cluster. Requires switches (fat-tree topology).
  \item 范围：集群中跨节点。需要交换机（fat-tree 拓扑）。
  % \item Use case: \textbf{Data Parallelism / FSDP} gradient synchronization. AllReduce of gradients happens once per training step (not per layer), so latency tolerance is higher.
  \item 用例：\textbf{数据并行 / FSDP} 梯度同步。梯度 AllReduce 每训练步发生一次（而非每层），因此延迟容忍度更高。
\end{itemize}

% \textbf{In RLHF specifically}:
\textbf{在 RLHF 中具体而言}：

\begin{itemize}
  % \item \emph{Generation}: TP=8 over NVLink within node. Multiple vLLM instances across nodes don’t communicate (embarrassingly parallel).
  \item \emph{生成}：节点内 NVLink 上 TP=8。跨节点的多个 vLLM 实例无需通信（embarrassingly parallel）。
  % \item \emph{Training}: TP=8 over NVLink intra-node + FSDP over InfiniBand inter-node. Gradients synced after full backward pass.
  \item \emph{训练}：节点内 NVLink 上 TP=8 + 节点间 InfiniBand 上 FSDP。梯度在完整反向传播后同步。
  % \item \emph{Weight sync}: Training $\to$ Generation uses InfiniBand (140 GB transfer, async, takes $\sim$3s at 50 GB/s).
  \item \emph{权重同步}：训练 $\to$ 生成使用 InfiniBand（140 GB 传输，异步，50 GB/s 下需 $\sim$3s）。
\end{itemize}

% \emph{\textbf{Review:} Chapters~2 and~11 (Systems Foundations; System Architecture).}
\emph{\textbf{复习：} 第~2 与~11 章（系统基础；系统架构）。}
\end{questionbox}

\section{优化与训练题}
\label{optimization-and-training-questions}

% \begin{questionbox}[Q34: Explain Adam vs AdamW. Why does the difference matter for LLMs?]
\begin{questionbox}[Q34：解释 Adam 与 AdamW。这一差异对 LLM 为何重要？]
% \textbf{Answer}: \textbf{Adam with L2 regularization}: $\theta_{t+1} = \theta_t - \alpha \cdot (\hat{m}_t / (\sqrt{\hat{v}_t} + \epsilon) + \lambda\theta_t)$. The weight decay term $\lambda\theta_t$ is \emph{inside} the adaptive scaling. Parameters with large gradients (large $v_t$) get \emph{less} weight decay (divided by $\sqrt{v_t}$). This is not true weight decay --- it’s scale-dependent.
\textbf{答}：\textbf{Adam 加 L2 正则化}：$\theta_{t+1} = \theta_t - \alpha \cdot (\hat{m}_t / (\sqrt{\hat{v}_t} + \epsilon) + \lambda\theta_t)$。权重衰减项 $\lambda\theta_t$ \emph{位于}自适应缩放\emph{内部}。梯度大（$v_t$ 大）的参数\emph{衰减更少}（除以 $\sqrt{v_t}$）。这不是真正的权重衰减——它与尺度相关。

% \textbf{AdamW (decoupled weight decay)}: $\theta_{t+1} = (1 - \alpha\lambda)\theta_t - \alpha \cdot \hat{m}_t / (\sqrt{\hat{v}_t} + \epsilon)$. Weight decay is applied \emph{outside} and \emph{before} the adaptive update. Every parameter gets the same proportional decay regardless of gradient history.
\textbf{AdamW（解耦权重衰减）}：$\theta_{t+1} = (1 - \alpha\lambda)\theta_t - \alpha \cdot \hat{m}_t / (\sqrt{\hat{v}_t} + \epsilon)$。权重衰减在自适应更新\emph{之外}且\emph{之前}施加。无论梯度历史如何，每个参数都获得相同比例的衰减。

% \textbf{Why it matters for LLMs}:
\textbf{对 LLM 为何重要}：

\begin{enumerate}
  % \item LLMs have parameters spanning many orders of magnitude (embedding layers vs attention vs FFN). Adam’s coupled L2 effectively penalizes small-gradient params more than large-gradient ones --- wrong behavior.
  \item LLM 的参数跨越多个数量级（embedding 层 vs attention vs FFN）。Adam 的耦合 L2 实际上对小梯度参数惩罚得比大梯度参数更多——错误行为。
  % \item Decoupled WD provides uniform regularization across all layers, preventing some layers from growing unbounded while others over-shrink.
  \item 解耦 WD 在所有层提供统一正则化，防止某些层无界增长而另一些层过度收缩。
  % \item Empirically: AdamW gives 2--5\% better perplexity on long pretraining runs compared to Adam+L2 with the same effective regularization.
  \item 经验：在长预训练上，AdamW 相比相同有效正则强度的 Adam+L2 给出 2--5\% 更好的 perplexity。
\end{enumerate}

% \textbf{For RL specifically}: Often use $\lambda = 0$ (no weight decay). The KL penalty provides regularization. But for SFT: $\lambda = 0.01$--$0.1$ with AdamW is standard.
\textbf{对 RL 而言}：通常使用 $\lambda = 0$（无权重衰减）。由 KL 惩罚提供正则化。但对 SFT：AdamW 配合 $\lambda = 0.01$--$0.1$ 是标准做法。

% \emph{\textbf{Review:} Chapter~1 (LLM Architecture and Optimization Methods).}
\emph{\textbf{复习：} 第~1 章（LLM 架构与优化方法）。}
\end{questionbox}

% \begin{questionbox}[Q35: Why is learning rate warmup necessary? What happens without it?]
\begin{questionbox}[Q35：为什么学习率 warmup 是必要的？不做 warmup 会发生什么？]
% \textbf{Answer}: \textbf{The problem}: Adam’s second moment estimate $v_t = \beta_2 v_{t-1} + (1-\beta_2)g_t^2$ starts at $v_0 = 0$. The bias correction $\hat{v}_t = v_t/(1-\beta_2^t)$ compensates mathematically, but in practice:
\textbf{答}：\textbf{问题}：Adam 的二阶矩估计 $v_t = \beta_2 v_{t-1} + (1-\beta_2)g_t^2$ 从 $v_0 = 0$ 起步。偏置修正 $\hat{v}_t = v_t/(1-\beta_2^t)$ 在数学上做了补偿，但在实践中：

\begin{itemize}
  % \item First few steps: $v_t$ is based on 1--5 gradient samples. Highly inaccurate estimate of true variance.
  \item 头几步：$v_t$ 基于 1--5 个梯度样本。对真实方差的估计极不准确。
  % \item If a parameter happens to get a small gradient initially, $v_t$ is tiny $\rightarrow$ effective LR is huge $\rightarrow$ catastrophic update.
  \item 若某参数初期恰好得到小梯度，$v_t$ 微小 $\rightarrow$ 有效 LR 巨大 $\rightarrow$ 灾难性更新。
  % \item The bias correction amplifies early updates: at step 1, $\hat{v}_1 = v_1/(1-0.999) = 1000 \cdot v_1$.
  \item 偏置修正放大早期更新：第 1 步时 $\hat{v}_1 = v_1/(1-0.999) = 1000 \cdot v_1$。
\end{itemize}

% \textbf{Without warmup}: First 10--100 steps often have gradient spikes that permanently damage the model. Early representations get scrambled before the optimizer stabilizes.
\textbf{无 warmup 时}：头 10--100 步经常出现梯度尖峰，永久性损坏模型。在优化器稳定之前，早期表征就被打乱。

% \textbf{Warmup fix}: Start with LR $\approx 0$ and linearly increase to target over $W$ steps (typically 3--10\% of training). By the time LR reaches full value, $v_t$ has accumulated enough samples to be accurate.
\textbf{Warmup 的修复}：从 LR $\approx 0$ 起步，在 $W$ 步（通常为训练量的 3--10\%）内线性升至目标值。当 LR 到达满值时，$v_t$ 已积累了足够样本以变得准确。

% \textbf{Typical settings}:
\textbf{典型设置}：

\begin{itemize}
  % \item Pretraining: 2000 steps warmup ($\sim$1\% of 200K steps)
  \item 预训练：2000 步 warmup（占 200K 步的 $\sim$1\%）。
  % \item SFT: 100 steps warmup ($\sim$5\% of 2000 steps)
  \item SFT：100 步 warmup（占 2000 步的 $\sim$5\%）。
  % \item RL (PPO/GRPO): 20--50 steps warmup (short, model already stable from SFT)
  \item RL（PPO/GRPO）：20--50 步 warmup（短，模型经 SFT 已稳定）。
\end{itemize}

% \emph{\textbf{Review:} Chapters~1 and~10 (LLM Architecture; SFT Best Practices).}
\emph{\textbf{复习：} 第~1 与~10 章（LLM 架构；SFT 最佳实践）。}
\end{questionbox}

% \begin{questionbox}[Q36: Compare learning rate schedules. Which would you choose for RL fine-tuning?]
\begin{questionbox}[Q36：比较学习率调度策略。RL 微调你会选哪种？]
\textbf{答}：

% \textbf{Cosine decay}: $\eta_t = \eta_\text{min} + \frac{1}{2}(\eta_\text{max} - \eta_\text{min})(1 + \cos(\pi t/T))$. Standard for pretraining and SFT. Smooth decay, most time at moderate LR.
\textbf{余弦衰减（Cosine decay）}：$\eta_t = \eta_\text{min} + \frac{1}{2}(\eta_\text{max} - \eta_\text{min})(1 + \cos(\pi t/T))$。预训练与 SFT 的标准选择。衰减平滑，大部分时间停留在中等 LR。

% \textbf{Linear decay}: $\eta_t = \eta_\text{max}(1 - t/T)$. Simpler, similar results to cosine for short runs.
\textbf{线性衰减}：$\eta_t = \eta_\text{max}(1 - t/T)$。更简单，短训练下与余弦衰减效果相近。

% \textbf{WSD (Warmup-Stable-Decay)}: Warmup $\rightarrow$ constant LR for 80\% $\rightarrow$ rapid decay in last 20\%. New standard for pretraining. The ``stable'' phase gives consistent learning; the final decay squeezes out remaining gains.
\textbf{预热-稳定-衰减（Warmup-Stable-Decay, WSD）}：warmup $\rightarrow$ 80\% 时间常数 LR $\rightarrow$ 最后 20\% 快速衰减。预训练的新标准。“稳定”阶段提供一致学习；最终衰减挤出剩余收益。

% \textbf{Constant}: No decay. $\eta_t = \eta_\text{max}$ after warmup.
\textbf{常数}：无衰减。warmup 之后 $\eta_t = \eta_\text{max}$。

% \textbf{For RL fine-tuning (PPO/GRPO), I’d choose: Constant with short warmup}. Reasons:
\textbf{对 RL 微调（PPO/GRPO），我会选择：短 warmup + 常数}。理由：

\begin{enumerate}
  % \item RL training length is highly unpredictable (you stop based on win-rate, not epochs).
  \item RL 训练长度高度不可预测（按胜率停训，而非 epoch）。
  % \item Cosine/linear decay assumes you know the total steps in advance.
  \item 余弦/线性衰减假设你提前知道总步数。
  % \item The LR is already very low ($10^{-6}$), further decay makes updates invisible.
  \item LR 已经非常低（$10^{-6}$），再衰减会让更新几乎不可见。
  % \item PPO’s adaptive KL controller already modulates the effective step size.
  \item PPO 的自适应 KL 控制器已在调控有效步长。
  % \item If you must decay: use linear decay over a generous budget, stop early when metrics plateau.
  \item 如必须衰减：在宽裕预算上做线性衰减，指标停滞时早停。
\end{enumerate}

% \emph{\textbf{Review:} Chapters~1 and~5 (LLM Architecture; PPO).}
\emph{\textbf{复习：} 第~1 与~5 章（LLM 架构；PPO）。}
\end{questionbox}

% \begin{questionbox}[Q37: Why is gradient clipping critical for RL training but less important for SFT?]
\begin{questionbox}[Q37：为何梯度裁剪（Gradient Clipping）对 RL 训练至关重要、对 SFT 却不那么重要？]
% \textbf{Answer}: \textbf{SFT}: Supervised loss is smooth and well-behaved. Gradient norms are consistent across batches (typically 0.1--1.0). Clipping at 1.0 rarely activates --- it’s a safety net.
\textbf{答}：\textbf{SFT}：监督 loss 平滑且行为良好。梯度范数在 batch 间一致（通常 0.1--1.0）。在 1.0 处裁剪很少触发——只是安全网。

% \textbf{RL (PPO/GRPO)}: Gradient norms are highly variable because:
\textbf{RL（PPO/GRPO）}：梯度范数高度多变，因为：

\begin{enumerate}
  % \item \textbf{Reward variance}: One batch might have all high-reward responses, next might have all low. The advantage $\hat{A}$ swings wildly.
  \item \textbf{Reward 方差}：一个 batch 可能全是高 reward 回答，下一个全是低的。advantage $\hat{A}$ 剧烈摆动。
  % \item \textbf{Ratio explosion}: If a rare token’s probability changed a lot, $r_t = \pi_\text{new}/\pi_\text{old}$ can be very large $\rightarrow$ large gradient before clipping kicks in.
  \item \textbf{比率爆炸}：若某个稀有 Token 的概率变化很大，$r_t = \pi_\text{new}/\pi_\text{old}$ 可能极大 $\rightarrow$ 在裁剪生效前就产生大梯度。
  % \item \textbf{Sparse reward}: In GRPO with binary rewards, some prompts give all-correct (advantage $\approx 0$) then suddenly a hard prompt gives extreme advantages.
  \item \textbf{稀疏 reward}：在二元 reward 的 GRPO 中，有些 prompt 全部正确（advantage $\approx 0$），突然一道难题给出极端 advantage。
  % \item \textbf{KL term}: The KL penalty gradient can spike when policy diverges.
  \item \textbf{KL 项}：当 policy 偏离时，KL 惩罚的梯度可能尖峰。
\end{enumerate}

% \textbf{Without clipping}: A single bad batch can produce a gradient 100$\times$ normal magnitude $\rightarrow$ destroys the model in one step. Recovery is impossible (catastrophic forgetting of all pretraining).
\textbf{无裁剪时}：一个坏 batch 就能产生比正常大 100$\times$ 的梯度 $\rightarrow$ 一步毁掉模型。无法恢复（所有预训练遭灾难性遗忘）。

% \textbf{Typical setting}: \texttt{max\_grad\_norm=1.0}. Some use 0.5 for extra safety in early RL training. The norm is computed globally across all parameters (not per-layer).
\textbf{典型设置}：\texttt{max\_grad\_norm=1.0}。一些人在 RL 训练早期用 0.5 以更安全。范数在所有参数上全局计算（非逐层）。

% \textbf{Monitoring}: If clipping activates more than 20\% of steps, your LR is probably too high or your batch size too small.
\textbf{监控}：若裁剪在超过 20\% 步触发，你的 LR 可能过高，或 batch size 过小。

% \emph{\textbf{Review:} Chapters~5 and~7 (PPO; GRPO).}
\emph{\textbf{复习：} 第~5 与~7 章（PPO；GRPO）。}
\end{questionbox}

% \begin{questionbox}[Q38: BF16 vs FP16 for training. When does the choice matter?]
\begin{questionbox}[Q38：训练中 BF16 vs FP16。这个选择何时重要？]
\textbf{答}：

% \textbf{FP16}: 1 sign + 5 exponent + 10 mantissa bits. Range: $\pm 65504$. Precision: $\sim$3.3 decimal digits.
\textbf{FP16}：1 符号 + 5 指数 + 10 尾数位。范围：$\pm 65504$。精度：$\sim$3.3 十进制位。

% \textbf{BF16}: 1 sign + 8 exponent + 7 mantissa bits. Range: $\pm 3.4 \times 10^{38}$ (same as FP32!). Precision: $\sim$2.4 decimal digits.
\textbf{BF16}：1 符号 + 8 指数 + 7 尾数位。范围：$\pm 3.4 \times 10^{38}$（与 FP32 相同！）。精度：$\sim$2.4 十进制位。

% \textbf{Why BF16 wins for LLMs}:
\textbf{为什么 BF16 在 LLM 上胜出}：

\begin{enumerate}
  % \item \textbf{No loss scaling needed}: FP16’s small range ($\pm$65K) means gradients and activations frequently overflow/underflow. Requires dynamic loss scaling (multiply loss by 1024, divide gradients back). BF16 has FP32’s range --- overflow is essentially impossible.
  \item \textbf{无需 loss scaling}：FP16 范围窄（$\pm$65K）导致梯度和激活频繁溢出/下溢。需要动态 loss scaling（loss 乘 1024，梯度再除回去）。BF16 拥有 FP32 的范围——溢出基本不可能。
  % \item \textbf{Simpler code}: No loss scaler, no inf/nan checking, no dynamic scaling adjustment.
  \item \textbf{代码更简单}：无 loss scaler，无 inf/nan 检查，无动态 scaling 调整。
  % \item \textbf{Critical for RL}: RL gradients are noisier and spikier than SFT. FP16 loss scaling often fails (picks wrong scale, causes NaN). BF16 ``just works.''
  \item \textbf{对 RL 关键}：RL 梯度比 SFT 更嘈杂、更易尖峰。FP16 的 loss scaling 经常失败（选错尺度，产生 NaN）。BF16 “开箱即用”。
\end{enumerate}

% \textbf{When FP16 might be better}: If you need maximum precision (some scientific computing tasks) and can manage the loss scaling. FP16 has 3 more mantissa bits = slightly more accurate results.
\textbf{FP16 可能更好的场景}：若你需要极致精度（某些科学计算任务）并能管理好 loss scaling。FP16 多 3 个尾数位 = 结果略更精确。

% \textbf{FP32 master weights}: Even with BF16 forward/backward, accumulate gradient updates in FP32 to prevent rounding errors from compounding over thousands of small steps. Standard practice for all LLM training.
\textbf{FP32 主权重}：即便前向/反向使用 BF16，也要在 FP32 中累加梯度更新，以防舍入误差在数千次小步中累积。所有 LLM 训练的标准做法。

% \emph{\textbf{Review:} Chapter~2 (Systems Foundations for LLMs).}
\emph{\textbf{复习：} 第~2 章（LLM 的系统基础）。}
\end{questionbox}

\section{Reward 模型与 SFT 题}
\label{reward-model-and-sft-questions}

% \begin{questionbox}[Q39: Derive the Bradley-Terry reward model loss. What are its limitations?]
\begin{questionbox}[Q39：推导 Bradley-Terry reward 模型的 loss。它有哪些局限？]
% \textbf{Answer}: \textbf{Bradley-Terry model}: Given two responses, the probability the better one ($y_w$) is preferred: $P(y_w \succ y_l | x) = \sigma(r(x, y_w) - r(x, y_l))$ where $\sigma$ is the sigmoid.
\textbf{答}：\textbf{Bradley-Terry 模型（Bradley-Terry Model）}：给定两个回答，更好的那个（$y_w$）被偏好的概率：$P(y_w \succ y_l | x) = \sigma(r(x, y_w) - r(x, y_l))$，其中 $\sigma$ 是 sigmoid。

% \textbf{MLE derivation}: Given $N$ preference pairs, maximize likelihood: $\prod_i P(y_w^i \succ y_l^i)$. Take negative log: $\mathcal{L} = -\sum_i \log\sigma(r(x_i, y_w^i) - r(x_i, y_l^i))$.
\textbf{MLE 推导}：给定 $N$ 个偏好对，最大化似然：$\prod_i P(y_w^i \succ y_l^i)$。取负对数：$\mathcal{L} = -\sum_i \log\sigma(r(x_i, y_w^i) - r(x_i, y_l^i))$。

% \textbf{Limitations}:
\textbf{局限}：

\begin{enumerate}
  % \item \textbf{No ties}: BT can’t model ``equally good'' --- forces a strict preference.
  \item \textbf{无平局}：BT 无法建模“同样好”——强制严格偏好。
  % \item \textbf{Transitivity}: Assumes if A$>$B and B$>$C then A$>$C. Humans aren’t transitive.
  \item \textbf{可传递性}：假设若 A$>$B 且 B$>$C 则 A$>$C。人类并不可传递。
  % \item \textbf{Context-free}: Same reward regardless of what alternatives were available.
  \item \textbf{上下文无关}：无论备选项是什么，reward 都相同。
  % \item \textbf{Scalar collapse}: Compresses all quality dimensions into one number. A response can be safe but unhelpful --- RM must trade off.
  \item \textbf{标量塌缩}：把所有质量维度压成一个数。一个回答可能既安全又无用——RM 必须权衡。
  % \item \textbf{Length bias}: Longer responses get higher scores (more information = more likely to contain what annotator wanted). Must explicitly decorrelate.
  \item \textbf{长度偏差}：更长的回答得到更高分（信息更多 = 更可能包含标注者想要的内容）。必须显式去相关。
\end{enumerate}

% \textbf{Mitigations}: Margin loss (require minimum gap $\delta$), reward centering (subtract running mean), length penalty during training, multi-head RM (separate scores for helpfulness/safety/accuracy).
\textbf{缓解措施}：margin loss（要求最小间隔 $\delta$）、reward 中心化（减去滑动均值）、训练时长度惩罚、多头 RM（为有用性/安全性/准确性提供独立分数）。

% \emph{\textbf{Review:} Chapter~9 (Reward Model Training).}
\emph{\textbf{复习：} 第~9 章（Reward 模型训练）。}
\end{questionbox}

% \begin{questionbox}[Q40: What is sequence packing in SFT and why does it matter?]
\begin{questionbox}[Q40：SFT 中的序列打包（Sequence Packing）是什么？为何重要？]
% \textbf{Answer}: \textbf{Problem}: Training examples have variable length. Standard batching pads all examples to max\_length. If max=4096 but average=500, you waste 88\% of compute on padding tokens (which contribute zero gradient).
\textbf{答}：\textbf{问题}：训练样本长度可变。标准批处理把所有样本填充到 max\_length。若 max=4096 而平均=500，就有 88\% 的算力浪费在填充 Token 上（梯度贡献为零）。

% \textbf{Packing solution}: Concatenate multiple short examples into a single max\_length sequence. Separate with EOS tokens. Train on all examples simultaneously.
\textbf{打包方案}：把多个短样本拼接成一条 max\_length 序列，用 EOS Token 分隔。同时对所有样本训练。

% Example: Instead of 4 sequences padded to 4096 (16K tokens, 14K padding), pack into 1 sequence of 4096 with 4 examples end-to-end (4096 real tokens, 0 padding). \textbf{4$\times$ more efficient.}
例子：不再用 4 条填充到 4096 的序列（16K Token，14K 填充），而是打包成 1 条 4096 的序列，首尾相连地容纳 4 个样本（4096 真实 Token，0 填充）。\textbf{效率提升 4$\times$。}

% \textbf{Critical detail --- block-diagonal attention mask}: Without special handling, example 2 attends to example 1’s tokens (cross-contamination). Must use a block-diagonal attention mask that restricts each example to only attend to its own tokens.
\textbf{关键细节——块对角 attention 掩码}：若不做特殊处理，样本 2 会关注样本 1 的 Token（交叉污染）。必须使用块对角 attention 掩码，限制每个样本只能关注自身的 Token。

% \textbf{In TRL}: \texttt{SFTConfig(packing=True, max\_seq\_length=4096)}. Handles mask automatically.
\textbf{在 TRL 中}：\texttt{SFTConfig(packing=True, max\_seq\_length=4096)}。自动处理掩码。

% \textbf{Caveats}: (1) Longer examples still need their own batch entries (can’t split mid-sequence). (2) Slight implementation complexity for position embeddings (reset per example). (3) Some argue packing changes the effective batch size (more examples per step) --- adjust LR accordingly.
\textbf{注意}：(1) 更长的样本仍需自己的 batch 条目（不能在序列中间切分）。(2) 位置 Embedding 实现略复杂（每个样本重置）。(3) 有人认为打包改变了有效 batch size（每步样本更多）——相应调整 LR。

% \emph{\textbf{Review:} Chapter~10 (SFT Best Practices and Techniques).}
\emph{\textbf{复习：} 第~10 章（SFT 最佳实践与技巧）。}
\end{questionbox}

% \begin{questionbox}[Q41: Explain completion-only masking for SFT. What happens if you don’t use it?]
\begin{questionbox}[Q41：解释 SFT 的 completion-only masking。不使用它会怎样？]
% \textbf{Answer}: In chat-format SFT data: \texttt{[system] + [user message] + [assistant response]}. Standard NLL loss computes loss on all tokens including the system prompt and user message.
\textbf{答}：在 chat 格式的 SFT 数据中：\texttt{[system] + [user message] + [assistant response]}。标准 NLL loss 会在所有 Token 上计算 loss，包括 system prompt 和用户消息。

% \textbf{Problem without masking}: The model wastes capacity learning to predict the user’s message (which it will never need to generate at inference). Worse: if the training data has diverse user messages, the model gets confused about ``whose turn is it?''
\textbf{不掩码的问题}：模型浪费容量去预测用户消息（推理时永远不需要生成）。更糟：若训练数据中用户消息多样，模型会困惑“现在轮到谁说话？”

% \textbf{Completion-only masking}: Set loss weight to 0 for all tokens in the prompt (system + user). Only compute loss on assistant response tokens.
\textbf{Completion-only masking}：将 prompt（system + user）中所有 Token 的 loss 权重设为 0。仅在 assistant 回答 Token 上计算 loss。

% TRL: \texttt{DataCollatorForCompletionOnlyLM(response\_template="<|assistant|>")}
TRL：\texttt{DataCollatorForCompletionOnlyLM(response\_template="<|assistant|>")}

% \textbf{Impact}: Typically 5--15\% better on instruction-following benchmarks. Faster convergence (gradient signal is concentrated on useful tokens). No change to compute cost.
\textbf{影响}：在指令跟随基准上通常提升 5--15\%。收敛更快（梯度信号集中于有用 Token）。算力成本不变。

% \textbf{Subtlety}: Must include the response template token in the loss (teaches the model to start responding). But exclude everything before it.
\textbf{细节}：response 模板 Token 必须包含在 loss 内（教模型开始回答），但之前的所有内容都要排除。

% \emph{\textbf{Review:} Chapter~10 (SFT Best Practices and Techniques).}
\emph{\textbf{复习：} 第~10 章（SFT 最佳实践与技巧）。}
\end{questionbox}

% \begin{questionbox}[Q42: How does SFT quality affect the RL ceiling? What is the pass@k diagnostic?]
\begin{questionbox}[Q42：SFT 质量如何影响 RL 的天花板？pass@k 诊断是什么？]
% \textbf{Answer}: \textbf{Ceiling theorem (informal)}: RL can only reinforce behaviors the model can already produce with non-negligible probability. If the SFT model has 0\% chance of generating a correct solution, RL will never find it.
\textbf{答}：\textbf{天花板定理（非正式表述）}：RL 只能强化模型已能以非可忽略概率产出的行为。若 SFT 模型生成正确解的概率为 0\%，RL 永远找不到它。

% \textbf{Why}: GRPO/PPO sample from the current policy and reinforce good samples. If good samples don’t exist in the distribution, there’s nothing to reinforce. RL is exploration-limited by the base policy’s support.
\textbf{为什么}：GRPO/PPO 从当前 policy 采样并强化好样本。若好样本不在分布内，就没什么可强化。RL 的探索受 base policy 支撑集（support）的限制。

% \textbf{pass@k diagnostic}: Generate $k$ responses per prompt, check if \emph{any} is correct:
\textbf{pass@k 诊断}：对每个 prompt 生成 $k$ 个回答，检查\emph{是否有任一个}正确：

\begin{itemize}
  % \item pass@1: Model’s typical performance (greedy/low temp).
  \item pass@1：模型的典型表现（greedy/低温）。
  % \item pass@8: Upper bound of what GRPO with $G=8$ can achieve.
  \item pass@8：$G=8$ 的 GRPO 所能达到的上界。
  % \item pass@64: Upper bound for aggressive Best-of-N.
  \item pass@64：激进 Best-of-N 采样（Rejection Sampling）的上界。
  % \item pass@256: Approximate ceiling for RL improvement.
  \item pass@256：RL 改进的近似天花板。
\end{itemize}

% \textbf{Interpretation}:
\textbf{解读}：

\begin{itemize}
  % \item pass@1=20\%, pass@64=80\%: Great! 4$\times$ headroom for RL. Strong gains expected.
  \item pass@1=20\%、pass@64=80\%：好！RL 有 4$\times$ 上行空间。预期强劲收益。
  % \item pass@1=20\%, pass@64=25\%: Almost no headroom. RL won’t help much. Need better SFT first.
  \item pass@1=20\%、pass@64=25\%：几乎没有上行空间。RL 帮不上忙。需先做更好的 SFT。
  % \item pass@1=5\%, pass@64=60\%: Model \emph{can} solve it but rarely does. Perfect case for RL (reinforce the rare successes).
  \item pass@1=5\%、pass@64=60\%：模型\emph{能}解但很少解出来。RL 的完美场景（强化稀有的成功）。
\end{itemize}

% \textbf{Rule}: If pass@64 $<$ 1.5$\times$ pass@1, invest in better SFT data before starting RL.
\textbf{规则}：若 pass@64 $<$ 1.5$\times$ pass@1，在开始 RL 之前先投入更好的 SFT 数据。

% \emph{\textbf{Review:} Chapters~7 and~10 (GRPO; SFT Best Practices).}
\emph{\textbf{复习：} 第~7 与~10 章（GRPO；SFT 最佳实践）。}
\end{questionbox}

% \begin{questionbox}[Q43: Design a multi-objective reward system for a chat model. How do you balance helpfulness vs safety?]
\begin{questionbox}[Q43：为聊天模型设计一个多目标 reward 系统。如何平衡有用性与安全性？]
% \textbf{Answer}: \textbf{Architecture}: Separate reward models for each objective:
\textbf{答}：\textbf{架构}：每个目标使用独立 reward 模型：

\begin{itemize}
  % \item $r_\text{helpful}$: Trained on helpfulness preferences (quality, accuracy, completeness)
  \item $r_\text{helpful}$：在有用性偏好上训练（质量、准确性、完整性）。
  % \item $r_\text{safe}$: Trained on safety preferences (refusals, harmlessness, no hallucination)
  \item $r_\text{safe}$：在安全偏好上训练（拒答、无害、无幻觉）。
  % \item $r_\text{format}$: Rule-based (follows instructions, proper formatting, appropriate length)
  \item $r_\text{format}$：基于规则（遵循指令、合适格式、合理长度）。
\end{itemize}

% \textbf{Combination strategies}:
\textbf{组合策略}：

\begin{enumerate}
  % \item \textbf{Weighted sum} (simplest): $r = w_1 r_\text{helpful} + w_2 r_\text{safe} + w_3 r_\text{format}$. Problem: safety can be outweighed by helpfulness.
  \item \textbf{加权求和}（最简）：$r = w_1 r_\text{helpful} + w_2 r_\text{safe} + w_3 r_\text{format}$。问题：安全性可能被有用性压过。
  % \item \textbf{Constrained} (safer): Maximize $r_\text{helpful}$ subject to $r_\text{safe} > \tau$. Implemented via: $r = r_\text{helpful} - \lambda \cdot \max(0, \tau - r_\text{safe})$ with large $\lambda$.
  \item \textbf{约束式}（更安全）：在 $r_\text{safe} > \tau$ 约束下最大化 $r_\text{helpful}$。通过 $r = r_\text{helpful} - \lambda \cdot \max(0, \tau - r_\text{safe})$ 实现，$\lambda$ 取较大值。
  % \item \textbf{GDPO normalization} (best for GRPO): Normalize each reward independently within group, then combine: $\hat{A} = w_1 \hat{A}_\text{helpful} + w_2 \hat{A}_\text{safe}$. Prevents one reward from dominating due to scale differences.
  \item \textbf{GDPO 归一化}（对 GRPO 最佳）：在 group 内独立归一化每个 reward，再组合：$\hat{A} = w_1 \hat{A}_\text{helpful} + w_2 \hat{A}_\text{safe}$。防止某个 reward 因尺度差异而主导。
  % \item \textbf{Lexicographic}: Safety is hard constraint (must pass), then optimize helpfulness. Train in stages: safety alignment first, then helpfulness.
  \item \textbf{字典序（Lexicographic）}：安全是硬约束（必须通过），再优化有用性。分阶段训练：先做安全对齐，再做有用性。
\end{enumerate}

% \textbf{Practical weights}: Start with $w_\text{safe}=2.0, w_\text{helpful}=1.0, w_\text{format}=0.5$. Safety gets 2$\times$ weight because its failure mode (harmful content) is much worse than helpfulness failure (mediocre answer).
\textbf{实用权重}：从 $w_\text{safe}=2.0, w_\text{helpful}=1.0, w_\text{format}=0.5$ 起步。安全权重 2$\times$，因为它的失败模式（有害内容）远比有用性失败（平庸答案）更糟。

% \emph{\textbf{Review:} Chapters~9 and~12 (Reward Model Training; LLM Agentic Training).}
\emph{\textbf{复习：} 第~9 与~12 章（Reward 模型训练；LLM 智能体训练）。}
\end{questionbox}

\section{系统架构扩展题}
\label{system-architecture-extension-questions}

% \begin{questionbox}[Q44: How does speculative decoding work? When does it help for RLHF?]
\begin{questionbox}[Q44：投机解码（Speculative Decoding）如何工作？它在 RLHF 中何时有帮助？]
% \textbf{Answer}: \textbf{Problem}: Large model generates one token per forward pass ($\sim$70ms for 70B). Slow.
\textbf{答}：\textbf{问题}：大模型每次前向只产出一个 Token（70B 约 $\sim$70ms）。慢。

% \textbf{Speculative decoding}:
\textbf{投机解码}：

\begin{enumerate}
  % \item \textbf{Draft}: Small model (1--7B) generates $k$ candidate tokens quickly ($\sim$5ms for all $k$).
  \item \textbf{草稿}：小模型（1--7B）快速生成 $k$ 个候选 Token（全部 $k$ 个约 $\sim$5ms）。
  % \item \textbf{Verify}: Large model does one forward pass scoring all $k$ tokens in parallel. Accepts tokens where $p_\text{large}(t_i) \geq p_\text{draft}(t_i)$ (always). Probabilistically accepts others.
  \item \textbf{验证}：大模型做一次前向，并行给所有 $k$ 个 Token 打分。$p_\text{large}(t_i) \geq p_\text{draft}(t_i)$ 的 Token 总是接受。其他按概率接受。
  % \item \textbf{Result}: On average, 3--4 tokens accepted per verification step. Speedup: 2--3$\times$.
  \item \textbf{结果}：每次验证步平均接受 3--4 个 Token。加速：2--3$\times$。
\end{enumerate}

% \textbf{Key property}: The output distribution is \emph{identical} to sampling from the large model alone. No quality loss. The draft model only affects speed, not output.
\textbf{关键性质}：输出分布与单独从大模型采样\emph{完全相同}。无质量损失。草稿模型只影响速度，不影响输出。

% \textbf{For RLHF specifically}: Generation is 60\% of compute. 2--3$\times$ speedup on generation = 1.5--2$\times$ end-to-end speedup. Combined with vLLM + INT8: generation goes from the bottleneck to parity with training.
\textbf{对 RLHF 而言}：生成占算力 60\%。生成加速 2--3$\times$ = 端到端加速 1.5--2$\times$。结合 vLLM + INT8：生成从瓶颈变为与训练持平。

% \textbf{Limitations}: (1) Draft model must share tokenizer. (2) Less effective at high temperature (draft model less accurate). (3) Needs additional GPU memory for draft model. (4) Diminishing returns beyond $k=5$ (acceptance rate drops).
\textbf{局限}：(1) 草稿模型必须共享 tokenizer。(2) 高温下效果较差（草稿模型准确率下降）。(3) 草稿模型需要额外 GPU 显存。(4) 超过 $k=5$ 后边际收益递减（接受率下降）。

% \emph{\textbf{Review:} Chapters~2 and~11 (Systems Foundations; System Architecture).}
\emph{\textbf{复习：} 第~2 与~11 章（系统基础；系统架构）。}
\end{questionbox}

% \begin{questionbox}[Q45: Explain the roofline model. How do you determine if a kernel is compute-bound or memory-bound?]
\begin{questionbox}[Q45：解释 roofline 模型。如何判断一个 kernel 是计算受限还是内存受限？]
% \textbf{Answer}: The roofline model plots achievable performance (FLOPS) as a function of arithmetic intensity (FLOPS per byte of memory traffic).
\textbf{答}：roofline 模型将可达性能（FLOPS）刻画为算术强度（每字节内存流量的 FLOPS）的函数。

% \textbf{Two regimes}:
\textbf{两个区域}：

\begin{itemize}
  % \item \textbf{Memory-bound} (left of crossover): Performance limited by how fast you can feed data to compute units. Actual FLOPS = bandwidth $\times$ arithmetic intensity. GPU utilization $<$ 100\%.
  \item \textbf{内存受限}（拐点左侧）：性能受限于向计算单元喂数据的速度。实际 FLOPS = 带宽 $\times$ 算术强度。GPU 利用率 $<$ 100\%。
  % \item \textbf{Compute-bound} (right of crossover): Performance limited by peak FLOPS. Memory is fast enough. GPU at max utilization.
  \item \textbf{计算受限}（拐点右侧）：性能受限于峰值 FLOPS。内存够快。GPU 达到最大利用率。
\end{itemize}

% \textbf{Crossover point}: Peak FLOPS / Peak Bandwidth. For A100: 312 TF / 2 TB/s = 156 FLOP/byte.
\textbf{拐点}：峰值 FLOPS / 峰值带宽。A100：312 TF / 2 TB/s = 156 FLOP/byte。

% \textbf{LLM operations}:
\textbf{LLM 操作}：

\begin{itemize}
  % \item \textbf{Autoregressive generation} (batch=1): Read 140GB weights, do 140G FLOPs = 1 FLOP/byte. \emph{Extremely} memory-bound (156$\times$ below crossover). Only 0.6\% GPU utilization.
  \item \textbf{自回归生成}（batch=1）：读取 140GB 权重，做 140G FLOPs = 1 FLOP/byte。\emph{极度}内存受限（比拐点低 156$\times$）。GPU 利用率仅 0.6\%。
  % \item \textbf{Training forward pass} (batch=128, seq=2048): Arithmetic intensity $\approx 200$+ FLOP/byte. Compute-bound. Near peak utilization.
  \item \textbf{训练前向}（batch=128，seq=2048）：算术强度 $\approx 200$+ FLOP/byte。计算受限。接近峰值利用率。
  % \item \textbf{Attention} (long sequence): $O(n^2 d)$ FLOPs / $O(n^2 + nd)$ bytes. For long $n$: compute-bound. For short $n$: memory-bound. FlashAttention keeps it in SRAM regardless.
  \item \textbf{Attention}（长序列）：$O(n^2 d)$ FLOPs / $O(n^2 + nd)$ 字节。$n$ 长：计算受限；$n$ 短：内存受限。FlashAttention 无论如何都让其留在 SRAM。
\end{itemize}

% \textbf{Practical use}: If your kernel is memory-bound, reduce memory traffic (quantization, caching, tiling). If compute-bound, reduce FLOPs (pruning, distillation, lower precision).
\textbf{实用法}：若 kernel 内存受限，减少内存流量（量化、缓存、分块）。若计算受限，减少 FLOPs（剪枝、蒸馏、降低精度）。

% \emph{\textbf{Review:} Chapter~2 (Systems Foundations for LLMs).}
\emph{\textbf{复习：} 第~2 章（LLM 的系统基础）。}
\end{questionbox}

% \begin{questionbox}[Q46: How does continuous batching work and why is it essential for RLHF generation?]
\begin{questionbox}[Q46：连续批处理（Continuous Batching）如何工作？为何它对 RLHF 生成至关重要？]
% \textbf{Answer}: \textbf{Static batching}: Start $B$ sequences. Wait for ALL to finish. If one sequence generates 500 tokens and another generates 50 tokens, the 50-token sequence’s GPU slot sits idle for 450 tokens.
\textbf{答}：\textbf{静态批处理}：启动 $B$ 条序列。等\emph{所有}序列结束。若一条生成 500 Token、另一条只生成 50 Token，那 50 Token 序列的 GPU 槽位会闲置 450 Token 的时间。

% \textbf{Continuous batching} (iteration-level scheduling): After each generation step, check which sequences are done. Immediately insert new sequences into freed slots. GPU slots are never idle.
\textbf{连续批处理}（迭代级调度）：每个生成步后检查哪些序列已完成。立即在腾出的槽位插入新序列。GPU 槽位从不闲置。

% \textbf{Why essential for RLHF}:
\textbf{为何对 RLHF 至关重要}：

\begin{enumerate}
  % \item RLHF generates diverse outputs (high temperature). Length variance is huge --- some responses are 50 tokens, others 2000+.
  \item RLHF 生成多样输出（高温度）。长度方差巨大——有些回答 50 Token，有些 2000+。
  % \item Without continuous batching: average utilization $\sim$40--50\% (waiting for slowest sequence).
  \item 无连续批处理：平均利用率 $\sim$40--50\%（等最慢的序列）。
  % \item With continuous batching: utilization $>$90\%. Throughput 2--3$\times$ higher.
  \item 有连续批处理：利用率 $>$90\%。吞吐高 2--3$\times$。
  % \item RLHF needs large batches (128+ responses per step). Generating 128 responses with static batching requires max\_tokens $\times$ 128 sequential steps. Continuous batching amortizes this.
  \item RLHF 需要大 batch（每步 128+ 回答）。用静态批处理生成 128 个回答需要 max\_tokens $\times$ 128 个串行步。连续批处理摊销了这一开销。
\end{enumerate}

% \textbf{Implementation}: vLLM’s scheduler checks after every decode step. Preemption: if a new high-priority request arrives and memory is full, can swap out a low-priority sequence’s KV cache to CPU and resume later.
\textbf{实现}：vLLM 的调度器在每个解码步后检查。抢占：若有新的高优先级请求到达且显存已满，可将低优先级序列的 KV cache 换出到 CPU，稍后恢复。

% \emph{\textbf{Review:} Chapters~2 and~11 (Systems Foundations; System Architecture).}
\emph{\textbf{复习：} 第~2 与~11 章（系统基础；系统架构）。}
\end{questionbox}

\section{Transformer 架构问题}
\label{transformer-architecture-questions}

% \begin{questionbox}[Q: Why does RoPE dominate over learned absolute positional embeddings in modern LLMs?]
\begin{questionbox}[Q：为什么在现代 LLM 中 RoPE 主导地位胜过可学习的绝对位置嵌入？]
% \textbf{Answer}: RoPE encodes \emph{relative} position directly into the Q/K dot product via rotation matrices. Key advantages:
\textbf{答}：RoPE 通过旋转矩阵将\emph{相对}位置直接编码到 Q/K 点积中。关键优势：

\begin{enumerate}
%   \item Attention scores depend only on relative distance $i-j$, not absolute position --- this generalizes better to unseen sequence lengths.
  \item Attention 分数只依赖于相对距离 $i-j$，而不依赖绝对位置——这能更好地泛化到未见过的序列长度。
%   \item Can be extended beyond training length via frequency scaling (NTK-aware, YaRN) without retraining.
  \item 可通过频率缩放（NTK-aware、YaRN）扩展到超出训练长度，无需重新训练。
%   \item No additional parameters (rotations are deterministic from position index).
  \item 无需额外参数（旋转由位置索引确定性地给出）。
%   \item Learned absolute embeddings are fixed to training length and don’t extrapolate --- a model trained at 4K context fails at 8K.
  \item 可学习的绝对嵌入固定到训练长度，无法外推——在 4K 上下文训练的模型在 8K 时会失败。
\end{enumerate}

% \emph{\textbf{Review:} Chapter~1 (LLM Architecture and Optimization Methods).}
\emph{\textbf{复习：}第 1 章（LLM 架构与优化方法）。}
\end{questionbox}

% \begin{questionbox}[Q: Explain SwiGLU and why it replaced ReLU in modern transformers.]
\begin{questionbox}[Q：解释 SwiGLU，以及它为何在现代 Transformer 中取代了 ReLU。]
% \textbf{Answer}: SwiGLU: $\text{FFN}(x) = W_2 (\text{Swish}(W_1 x) \odot W_3 x)$, where $\text{Swish}(x) = x \cdot \sigma(x)$.
\textbf{答}：SwiGLU：$\text{FFN}(x) = W_2 (\text{Swish}(W_1 x) \odot W_3 x)$，其中 $\text{Swish}(x) = x \cdot \sigma(x)$。

% \textbf{Why it’s better}:
\textbf{为什么它更好}：

\begin{itemize}
%   \item The \emph{gating} mechanism ($\odot W_3 x$) allows the network to selectively suppress or amplify dimensions --- more expressive than pointwise ReLU.
  \item \emph{门控}机制（$\odot W_3 x$）让网络可以有选择地抑制或放大维度——比逐点 ReLU 更具表达力。
%   \item Swish is smooth (no dead neurons like ReLU’s zero-gradient region).
  \item Swish 是平滑的（不存在 ReLU 零梯度区那样的死神经元）。
%   \item Empirically: 1--2\% improvement on language modeling benchmarks at same FLOP count.
  \item 经验上：在相同 FLOP 预算下，语言建模基准上提升 1\%--2\%。
%   \item Tradeoff: requires 3 weight matrices instead of 2 (solved by reducing hidden dim from $4d$ to $8d/3$).
  \item 权衡：需要 3 个权重矩阵而非 2 个（通过将隐藏维度从 $4d$ 降到 $8d/3$ 来解决）。
\end{itemize}

\emph{\textbf{复习：}第 1 章（LLM 架构与优化方法）。}
\end{questionbox}

% \begin{questionbox}[Q: What is Grouped Query Attention (GQA) and why does Llama-3 use it?]
\begin{questionbox}[Q：什么是分组查询注意力（Grouped Query Attention, GQA），Llama-3 为何采用它？]
% \textbf{Answer}: Standard MHA: $H$ query heads, $H$ key heads, $H$ value heads. GQA: $H$ query heads but only $G < H$ key/value heads (shared across query groups).
\textbf{答}：标准 MHA：$H$ 个 query 头、$H$ 个 key 头、$H$ 个 value 头。GQA：$H$ 个 query 头但只有 $G < H$ 个 key/value 头（在 query 组间共享）。

% Llama-3 70B: 64 query heads, 8 KV heads (each KV head shared by 8 query heads).
Llama-3 70B：64 个 query 头、8 个 KV 头（每个 KV 头被 8 个 query 头共享）。

% \textbf{Benefits}:
\textbf{优势}：

\begin{itemize}
%   \item KV cache size reduced by $H/G = 8\times$ --- critical for inference (KV cache is the dominant memory cost at long sequences).
  \item KV 缓存大小缩小 $H/G = 8\times$——对推理至关重要（长序列下 KV 缓存是主导的显存开销）。
%   \item Minimal quality loss ($<$0.5\% on benchmarks) because KV patterns are highly correlated across heads.
  \item 质量损失极小（基准上 $<$0.5\%），因为 KV 模式在不同头之间高度相关。
%   \item Inference throughput increases proportionally to KV cache reduction (more sequences fit in memory = higher batch size).
  \item 推理吞吐与 KV 缓存的减少成比例增长（更多序列能装进显存 = 更大的 Batch）。
\end{itemize}

\emph{\textbf{复习：}第 1 章（LLM 架构与优化方法）。}
\end{questionbox}

% \begin{questionbox}[Q: Why did decoder-only architectures win over encoder-decoder for LLMs?]
\begin{questionbox}[Q：为何 decoder-only 架构在 LLM 中胜过了 encoder-decoder？]
% \textbf{Answer}:
\textbf{答}：

\begin{enumerate}
%   \item \textbf{Unified objective}: Pretraining = fine-tuning = inference all use next-token prediction. No architectural mismatch.
  \item \textbf{统一目标}：预训练 = 微调 = 推理，都使用下一个 token 预测。不存在架构不匹配。
%   \item \textbf{Parameter efficiency}: All parameters contribute to generation. In encoder-decoder, encoder params are ``wasted'' during pure generation tasks.
  \item \textbf{参数效率}：所有参数都对生成有贡献。在 encoder-decoder 中，纯生成任务下 encoder 参数被“浪费”。
%   \item \textbf{Simpler scaling}: One model, one loss function, one set of hyperparameters to tune.
  \item \textbf{扩展更简单}：一个模型、一个损失函数、一套超参数需要调。
%   \item \textbf{KV cache efficiency}: Decoder-only has one KV cache; encoder-decoder has two (encoder + decoder cross-attention).
  \item \textbf{KV 缓存效率}：decoder-only 只有一份 KV 缓存；encoder-decoder 有两份（encoder + decoder 的交叉注意力）。
%   \item \textbf{Emergent few-shot}: Decoder-only naturally supports in-context learning (prepend examples to the prompt).
  \item \textbf{涌现的少样本能力}：decoder-only 天然支持上下文学习（把示例前置到 Prompt 中）。
\end{enumerate}

% Encoder-decoder still wins for seq2seq tasks with fixed input length (translation), but these are a shrinking fraction of LLM use cases.
对于固定输入长度的 seq2seq 任务（如翻译），encoder-decoder 仍然占优，但这部分在 LLM 用例中占比正在缩小。

\emph{\textbf{复习：}第 1 章（LLM 架构与优化方法）。}
\end{questionbox}

\section{FlashAttention 问题}
\label{flash-attention-questions}

% \begin{questionbox}[Q: FlashAttention computes the same result as standard attention but is 2--4$\times$ faster. How is this possible if it does the same number of FLOPs?]
\begin{questionbox}[Q：FlashAttention 计算的结果与标准 attention 相同，但快 2--4 倍。如果 FLOPs 数相同，这怎么可能？]
% \textbf{Answer}: FlashAttention is faster because it reduces \emph{HBM memory traffic}, not FLOPs. Standard attention materializes the $n \times n$ attention matrix in HBM (slow), reads it back for softmax, reads again for $PV$ multiply --- 4 HBM round-trips over $O(n^2)$ data.
\textbf{答}：FlashAttention 之所以更快，是因为它减少了\emph{HBM 显存流量}，而非 FLOPs。标准 attention 在 HBM（慢）中物化 $n \times n$ 的 attention 矩阵，再读回来做 Softmax，再读一次做 $PV$ 乘法——总共在 $O(n^2)$ 数据上有 4 次 HBM 往返。

% FlashAttention tiles the computation so the $n \times n$ matrix is computed and consumed entirely in SRAM (fast, 19 TB/s) without ever writing it to HBM (2 TB/s). The ``online softmax'' trick enables this by maintaining running statistics.
FlashAttention 对计算进行分块，使 $n \times n$ 矩阵完全在 SRAM（快，19 TB/s）中计算并消费，从不写入 HBM（2 TB/s）。“在线 Softmax”技巧通过维护滚动统计量来实现这一点。

% Result: HBM traffic drops from $O(n^2 d)$ to $O(n^2 d / M)$ where $M$ is SRAM size. Same FLOPs, 10--50$\times$ less memory traffic $\to$ 2--4$\times$ wall-clock speedup.
结果：HBM 流量从 $O(n^2 d)$ 降到 $O(n^2 d / M)$，其中 $M$ 为 SRAM 大小。相同 FLOPs，显存流量减少 10--50$\times$ $\to$ 实际墙钟加速 2--4$\times$。

% \emph{\textbf{Review:} Chapters~1 and~2 (LLM Architecture; Systems Foundations).}
\emph{\textbf{复习：}第 1 章与第 2 章（LLM 架构；系统基础）。}
\end{questionbox}

% \begin{questionbox}[Q: Why doesn’t FlashAttention help the FFN layers?]
\begin{questionbox}[Q：为什么 FlashAttention 对 FFN 层没帮助？]
% \textbf{Answer}: FFN layers are \emph{compute-bound}, not memory-bound. Their arithmetic intensity ($I \approx 300$ FLOP/byte for large batch GEMMs) is already above the roofline ridge point (156 FLOP/byte on A100).
\textbf{答}：FFN 层是\emph{计算瓶颈}，而非显存瓶颈。其算术强度（大 Batch GEMM 下 $I \approx 300$ FLOP/byte）已经高于 roofline 拐点（A100 上为 156 FLOP/byte）。

% FlashAttention helps attention because attention is deeply memory-bound ($I \approx 1$--$60$ FLOP/byte). By keeping data in SRAM, it removes the memory bottleneck.
FlashAttention 对 attention 有帮助是因为 attention 严重受限于显存（$I \approx 1$--$60$ FLOP/byte）。通过将数据保留在 SRAM 中，它消除了显存瓶颈。

% For FFN: the bottleneck is already the Tensor Cores (not memory bandwidth), so reducing memory traffic doesn’t help. Instead, FFN benefits from quantization (reduces weight size $\to$ higher arithmetic intensity) and larger batch sizes.
对 FFN 而言：瓶颈已经是 Tensor Core（而非显存带宽），所以减少显存流量没用。FFN 的收益来自量化（减小权重大小 $\to$ 提高算术强度）和更大的 Batch。

\emph{\textbf{复习：}第 1 章与第 2 章（LLM 架构；系统基础）。}
\end{questionbox}

% \begin{questionbox}[Q: Explain the online softmax trick and why it’s essential for FlashAttention.]
\begin{questionbox}[Q：解释在线 Softmax 技巧，以及它对 FlashAttention 为何不可或缺。]
% \textbf{Answer}: Standard softmax needs the global maximum $m = \max_j x_j$ before computing any output --- this requires seeing all $n$ attention scores first, forcing materialization of the full $n \times n$ matrix.
\textbf{答}：标准 Softmax 在计算任何输出之前都需要全局最大值 $m = \max_j x_j$——这要求先看到所有 $n$ 个 attention 分数，迫使物化完整的 $n \times n$ 矩阵。

% The online softmax trick processes blocks sequentially, maintaining a running $(m, \ell, O)$ state:
在线 Softmax 技巧按顺序处理块，维护一个滚动的 $(m, \ell, O)$ 状态：

\begin{enumerate}
%   \item Process new block $\to$ update running max: $m_{\text{new}} = \max(m_{\text{old}}, \max(s_{\text{new}}))$
  \item 处理新块 $\to$ 更新滚动最大值：$m_{\text{new}} = \max(m_{\text{old}}, \max(s_{\text{new}}))$
%   \item Rescale old sum: $\ell_{\text{new}} = e^{m_{\text{old}} - m_{\text{new}}} \cdot \ell_{\text{old}} + \text{new terms}$
  \item 重新缩放旧的求和：$\ell_{\text{new}} = e^{m_{\text{old}} - m_{\text{new}}} \cdot \ell_{\text{old}} + \text{new terms}$
%   \item Rescale output: $O_{\text{new}} = \text{rescaled}(O_{\text{old}}) + \text{new contribution}$
  \item 重新缩放输出：$O_{\text{new}} = \text{rescaled}(O_{\text{old}}) + \text{new contribution}$
\end{enumerate}

% This is mathematically exact --- no approximation. It enables block-by-block processing where each block fits in SRAM, never needing the full $n \times n$ matrix in memory.
这在数学上是完全精确的——没有任何近似。它使得逐块处理成为可能，每块都能放进 SRAM，永远不需要在显存中保留完整的 $n \times n$ 矩阵。

\emph{\textbf{复习：}第 1 章与第 2 章（LLM 架构；系统基础）。}
\end{questionbox}

\section{LoRA 与 PEFT 问题}
\label{lora-and-peft-questions}

% \begin{questionbox}[Q: Why does LoRA work? What theoretical insight justifies low-rank updates?]
\begin{questionbox}[Q：LoRA 为什么有效？什么理论洞见支持低秩更新？]
% \textbf{Answer}: Aghajanyan et al.~\cite{aghajanyan2020intrinsic} showed that fine-tuning operates in a very low \emph{intrinsic dimensionality} --- the effective parameter space for a fine-tuning task is far smaller than the model’s total parameter count. A 175B model may have intrinsic dimensionality $<$10,000 for a given task.
\textbf{答}：Aghajanyan 等人~\cite{aghajanyan2020intrinsic} 表明，微调是在一个非常低的\emph{内在维度（intrinsic dimensionality）}上进行的——某个微调任务的有效参数空间远远小于模型的总参数量。175B 模型在给定任务上的内在维度可能 $<$10,000。

% LoRA directly exploits this: by constraining updates to rank $r$ ($W' = W + BA$, $B \in \mathbb{R}^{d \times r}$), it restricts learning to an $r$-dimensional subspace per weight matrix. Since the true task subspace is low-dimensional, this loses almost nothing while reducing trainable parameters by 100--1000$\times$.
LoRA 直接利用了这一点：通过将更新约束到秩 $r$（$W' = W + BA$，$B \in \mathbb{R}^{d \times r}$），它将每个权重矩阵的学习限制在一个 $r$ 维子空间内。由于真实的任务子空间是低维的，这几乎没有任何损失，同时将可训练参数减少 100--1000$\times$。

% \textbf{Intuition}: Fine-tuning doesn’t change what the model ``knows'' (the full-rank $W$ stays frozen); it only adjusts \emph{how} existing knowledge is combined for the new task --- a low-rank perturbation.
\textbf{直觉}：微调并不改变模型“知道什么”（满秩 $W$ 保持冻结）；它只调整\emph{如何}为新任务组合现有知识——一个低秩扰动。

\emph{\textbf{复习：}第 1 章（LLM 架构与优化方法）。}
\end{questionbox}

% \begin{questionbox}[Q: Compare QLoRA vs. full LoRA vs. full fine-tuning for a 70B model. When would you choose each?]
\begin{questionbox}[Q：对 70B 模型，比较 QLoRA、全 LoRA、全量微调。何时选择每种？]
\textbf{答}：

{\small
% \begin{tabular}{@{}llllp{4.5cm}@{}}
\begin{tabularx}{\textwidth}{@{}llllX@{}}
\toprule
% \textbf{Method} & \textbf{Memory} & \textbf{GPUs} & \textbf{Quality} & \textbf{Use When} \\
\textbf{方法} & \textbf{显存} & \textbf{GPU 数} & \textbf{质量} & \textbf{适用场景} \\
\midrule
% Full fine-tune & 560+ GB & 8+ A100 & Best & Unlimited budget; pre-training continuation \\
全量微调 & 560+ GB & 8+ A100 & 最佳 & 预算无上限；持续预训练 \\
% LoRA ($r=16$) & 145 GB & 2 A100 & 95--98\% & Good budget; general fine-tuning \\
LoRA ($r=16$) & 145 GB & 2 A100 & 95--98\% & 预算充足；一般微调 \\
% QLoRA ($r=16$) & 44 GB & 1$\times$48GB & 93--96\% & Single-GPU; prototyping; constrained resources \\
QLoRA ($r=16$) & 44 GB & 1$\times$48GB & 93--96\% & 单 GPU；原型；资源受限 \\
\bottomrule
\end{tabularx}
}

% \textbf{Decision tree}: (1) If task requires deep knowledge change $\to$ full fine-tune. (2) If adapting to new style/format $\to$ LoRA. (3) If memory-constrained or rapid iteration $\to$ QLoRA. (4) If rank matters: start $r=16$; increase if training loss plateaus above full fine-tune level.
\textbf{决策树}：(1) 若任务需要深层知识变化 $\to$ 全量微调。(2) 若适配新风格/格式 $\to$ LoRA。(3) 若显存受限或快速迭代 $\to$ QLoRA。(4) 若秩很关键：从 $r=16$ 开始；如果训练 Loss 在全量微调水平之上停滞则增大。

% \emph{\textbf{Review:} Chapters~1 and~10 (LLM Architecture; SFT Best Practices).}
\emph{\textbf{复习：}第 1 章与第 10 章（LLM 架构；SFT 最佳实践）。}
\end{questionbox}

% \begin{questionbox}[Q: What is DoRA and why does it outperform standard LoRA?]
\begin{questionbox}[Q：什么是 DoRA，为何它优于标准 LoRA？]
% \textbf{Answer}: DoRA (Weight-Decomposed Low-Rank Adaptation) decomposes $W$ into magnitude $\|W\|$ and direction $W/\|W\|$, then applies LoRA only to the direction component:
\textbf{答}：DoRA（权重分解低秩自适应，Weight-Decomposed Low-Rank Adaptation）将 $W$ 分解为幅度 $\|W\|$ 和方向 $W/\|W\|$，然后仅对方向分量应用 LoRA：
\[
W' = m \odot \frac{W + BA}{\|W + BA\|}
\]
% where $m$ (magnitude) is also trainable but as a simple scalar per output neuron.
其中 $m$（幅度）也是可训练的，但仅作为每个输出神经元上的简单标量。

% \textbf{Why it helps}: Full fine-tuning naturally updates both magnitude and direction independently. Standard LoRA couples them (the low-rank update changes both simultaneously in a constrained way). DoRA decouples them, giving LoRA the same ``degrees of freedom'' structure as full fine-tuning. Result: 1--3\% improvement on reasoning tasks with no extra compute at inference (merge adapters).
\textbf{为何有效}：全量微调自然地独立更新幅度和方向。标准 LoRA 将二者耦合（低秩更新以受约束的方式同时改变两者）。DoRA 将它们解耦，让 LoRA 拥有与全量微调相同的“自由度”结构。结果：推理任务上提升 1\%--3\%，推理阶段无额外计算开销（合并适配器即可）。

\emph{\textbf{复习：}第 1 章（LLM 架构与优化方法）。}
\end{questionbox}

\section{模型压缩问题}
\label{model-compression-questions}

% \begin{questionbox}[Q: Explain AWQ. Why does protecting 1\% of weights preserve 99\% of quality?]
\begin{questionbox}[Q：解释 AWQ。为什么保护 1\% 的权重就能保留 99\% 的质量？]
% \textbf{Answer}: AWQ (Activation-Aware Weight Quantization) observes that weight importance is highly non-uniform: weights that multiply large activations contribute disproportionately to the output.
\textbf{答}：AWQ（激活感知权重量化，Activation-Aware Weight Quantization）观察到权重的重要性高度不均匀：乘以大激活值的权重对输出的贡献不成比例地大。

% The key insight: $\|W \cdot X\|$ depends on both $W$ and $X$. A small weight multiplied by a large activation matters more than a large weight multiplied by a near-zero activation.
关键洞见：$\|W \cdot X\|$ 同时取决于 $W$ 和 $X$。一个小权重乘以一个大激活，比一个大权重乘以一个接近零的激活更重要。

% AWQ identifies the top 1\% of ``salient'' channels (those with consistently large activation magnitudes across calibration data) and protects them by scaling: multiply salient channels by a factor $s > 1$ before quantization (then divide by $s$ in the activation). This reduces relative quantization error for important channels.
AWQ 识别出 top 1\% 的“显著”通道（在校准数据上始终具有较大激活幅度的通道），并通过缩放对其加以保护：量化前将显著通道乘以因子 $s > 1$（然后在激活中除以 $s$）。这降低了重要通道的相对量化误差。

% Result: 4-bit quantization with $<$1\% quality loss on 70B models, because the 99\% of non-salient weights can tolerate aggressive quantization.
结果：在 70B 模型上 4-bit 量化的质量损失 $<$1\%，因为 99\% 的非显著权重能容忍激进的量化。

\emph{\textbf{复习：}第 1 章与第 2 章（LLM 架构；系统基础）。}
\end{questionbox}

% \begin{questionbox}[Q: When should you use FP8 vs. 4-bit quantization vs. BF16?]
\begin{questionbox}[Q：什么时候用 FP8、4-bit 量化、BF16？]
\textbf{答}：

\begin{itemize}
%   \item \textbf{BF16}: Training (policy model in RLHF), when precision matters. Default for any model being updated by gradients.
  \item \textbf{BF16}：训练（RLHF 中的 Policy 模型），当精度重要时。任何被 Gradient 更新的模型的默认选择。
%   \item \textbf{FP8 (E4M3)}: H100 training with Transformer Engine (2$\times$ throughput, $<$0.5\% quality loss). Also for inference on H100 when you need maximum throughput.
  \item \textbf{FP8 (E4M3)}：H100 上使用 Transformer Engine 训练（2$\times$ 吞吐，$<$0.5\% 质量损失）。也用于 H100 上需要最大吞吐的推理。
%   \item \textbf{INT8/FP8 inference}: Frozen models in RLHF (reference model, reward model) --- not being trained, so reduced precision is safe.
  \item \textbf{INT8/FP8 推理}：RLHF 中的冻结模型（reference 模型、Reward 模型）——不被训练，所以降低精度是安全的。
%   \item \textbf{4-bit (AWQ/GPTQ)}: Inference serving at scale. Best memory/quality tradeoff for deployment. Also for QLoRA base model.
  \item \textbf{4-bit (AWQ/GPTQ)}：大规模推理服务。部署中显存/质量权衡最佳。也用于 QLoRA 的基模。
%   \item \textbf{2-bit}: Experimental; edge deployment where memory is extreme constraint. Quality loss 5--10\%.
  \item \textbf{2-bit}：实验性；显存极度受限的边缘部署。质量损失 5\%--10\%。
\end{itemize}

% \textbf{Rule}: quantize as aggressively as possible for inference, keep BF16 (or FP8 on H100) for training.
\textbf{规则}：推理时尽可能激进地量化，训练时保持 BF16（H100 上可用 FP8）。

% \emph{\textbf{Review:} Chapter~2 (Systems Foundations for LLMs).}
\emph{\textbf{复习：}第 2 章（LLM 的系统基础）。}
\end{questionbox}

% \begin{questionbox}[Q: Explain NVIDIA 2:4 structured sparsity. What’s the speedup and constraint?]
\begin{questionbox}[Q：解释 NVIDIA 2:4 结构化稀疏。加速比和约束是什么？]
% \textbf{Answer}: 2:4 sparsity means: in every group of 4 consecutive elements, exactly 2 must be zero. This is enforced at the weight level.
\textbf{答}：2:4 稀疏意味着：每 4 个连续元素一组中，恰好有 2 个必须为零。这在权重层面强制执行。

% \textbf{Hardware support}: A100/H100 Tensor Cores have dedicated 2:4 sparse GEMM instructions that skip the zero elements, achieving exactly \textbf{2$\times$ throughput} with no software overhead.
\textbf{硬件支持}：A100/H100 的 Tensor Core 有专门的 2:4 稀疏 GEMM 指令，可以跳过零元素，无软件开销地实现恰好 \textbf{2$\times$ 吞吐}。

% \textbf{Constraint}: You must achieve \emph{exactly} 50\% sparsity in this specific pattern. You can’t have 30\% or 70\% sparsity; you can’t have arbitrary sparsity patterns. The pruning must respect the 4-element group structure.
\textbf{约束}：你必须在这一特定模式下实现\emph{恰好}50\% 的稀疏度。不能是 30\% 或 70\% 的稀疏度；不能是任意稀疏模式。剪枝必须遵守 4 元素组结构。

% \textbf{How to achieve it}: After training (or during fine-tuning), for each group of 4 weights, zero out the 2 smallest by magnitude. Then fine-tune for a few hundred steps to recover quality. Quality loss: typically $<$1\% for large models (70B+).
\textbf{如何实现}：训练之后（或微调期间），对每组 4 个权重，将幅度最小的 2 个置零。然后微调数百步以恢复质量。质量损失：大模型（70B+）通常 $<$1\%。

\emph{\textbf{复习：}第 2 章（LLM 的系统基础）。}
\end{questionbox}

\section{专家混合（Mixture of Experts）问题}
\label{mixture-of-experts-questions}

% \begin{questionbox}[Q: Mixtral 8x7B has 47B total parameters but only 13B active per token. Explain how this works and why it’s efficient.]
\begin{questionbox}[Q：Mixtral 8x7B 总共有 47B 参数但每个 token 只激活 13B。请解释它如何工作，以及为何高效。]
% \textbf{Answer}: Mixtral replaces each FFN layer with 8 parallel expert FFNs (each $\sim$7B params for the FFN portion). A router network selects the Top-2 experts per token.
\textbf{答}：Mixtral 将每个 FFN 层替换为 8 个并行的专家 FFN（每个专家的 FFN 部分约 $\sim$7B 参数）。一个 router 网络为每个 token 选择 Top-2 专家。

% \textbf{Why 47B total}: Attention layers are shared (not replicated) = $\sim$5B. FFN experts: 8 $\times$ $\sim$5.25B = 42B. Total: $\sim$47B.
\textbf{为什么总共 47B}：Attention 层是共享的（不复制）= $\sim$5B。FFN 专家：8 $\times$ $\sim$5.25B = 42B。总计：$\sim$47B。

% \textbf{Why 13B active}: Per token, only 2 experts fire. Active params = attention ($\sim$5B) + 2 FFN experts ($\sim$2 $\times$ 5.25B) $\approx$ 13B.
\textbf{为什么 13B 激活}：每个 token 只有 2 个专家被激活。激活参数 = attention（$\sim$5B）+ 2 个 FFN 专家（$\sim$2 $\times$ 5.25B）$\approx$ 13B。

% \textbf{Why efficient}: Compute cost scales with \emph{active} params (13B), matching a 13B dense model. But capacity (knowledge stored) scales with \emph{total} params (47B), matching much larger models. Result: Mixtral matches Llama-2 70B quality at 13B compute cost.
\textbf{为什么高效}：计算成本随\emph{激活}参数（13B）扩展，与 13B 稠密模型相当。但容量（存储的知识）随\emph{总}参数（47B）扩展，与更大的模型相当。结果：Mixtral 以 13B 的计算成本匹配 Llama-2 70B 的质量。

% \textbf{Memory cost}: Still need all 47B params in memory (all experts loaded), so memory = 47B model, but compute = 13B model.
\textbf{显存成本}：仍需把全部 47B 参数装进显存（所有专家都加载），所以显存 = 47B 模型，但计算 = 13B 模型。

\emph{\textbf{复习：}第 1 章（LLM 架构与优化方法）。}
\end{questionbox}

% \begin{questionbox}[Q: What is the load balancing problem in MoE and how is it solved?]
\begin{questionbox}[Q：MoE 中的负载均衡问题是什么，如何解决？]
% \textbf{Answer}: Without constraints, the router tends to send most tokens to 1--2 ``popular'' experts (rich-get-richer dynamics). This causes:
\textbf{答}：不加约束时，router 倾向于将大部分 token 路由到 1--2 个“热门”专家（强者愈强的正反馈）。这导致：

\begin{itemize}
%   \item Capacity waste: 6 of 8 experts are unused, model effectively shrinks to 2-expert size.
  \item 容量浪费：8 个专家中有 6 个未被使用，模型实际上缩小到 2 个专家的规模。
%   \item Compute imbalance: If each expert is on a different GPU, popular experts become bottlenecks while others idle.
  \item 计算不均衡：如果每个专家在不同 GPU 上，热门专家会成为瓶颈，其他专家闲置。
\end{itemize}

% \textbf{Solution}: Auxiliary load-balancing loss: $\mathcal{L}_{\text{bal}} = \alpha \cdot N \sum_{i=1}^N f_i \cdot p_i$, where $f_i$ = fraction of tokens routed to expert $i$, $p_i$ = mean router probability for expert $i$. This penalizes uneven distributions.
\textbf{解决方案}：辅助负载均衡 Loss：$\mathcal{L}_{\text{bal}} = \alpha \cdot N \sum_{i=1}^N f_i \cdot p_i$，其中 $f_i$ = 路由到专家 $i$ 的 token 比例，$p_i$ = 专家 $i$ 的平均 router 概率。这会惩罚不均匀的分布。

% \textbf{Alternative}: Expert capacity factor --- hard cap on max tokens per expert per batch. Overflow tokens are dropped or re-routed.
\textbf{替代方案}：专家容量因子——对每个 Batch 中每个专家的最大 token 数做硬上限。溢出的 token 被丢弃或重新路由。

% Typical $\alpha$: 0.01--0.1 (small enough not to hurt main loss, large enough to prevent collapse).
典型 $\alpha$：0.01--0.1（足够小以不损害主 Loss，足够大以防止坍缩）。

\emph{\textbf{复习：}第 1 章（LLM 架构与优化方法）。}
\end{questionbox}

\section{训练中的多样性问题}
\label{diversity-in-training-questions}

% \begin{questionbox}[Q: What happens if all N responses in a GRPO group are identical?]
\begin{questionbox}[Q：如果一个 GRPO 组中 N 个响应全部相同会怎样？]
% \textbf{Answer}: If all $N$ responses are identical: all rewards $r_i$ are equal, so $\sigma_G = 0$ and advantages $\hat{A}_i = (r_i - \mu_G)/\sigma_G$ are undefined (division by zero). In practice, implementations set $\hat{A}_i = 0$ for all, meaning \textbf{zero learning signal} --- the step is wasted.
\textbf{答}：如果全部 $N$ 个响应相同：所有 Reward $r_i$ 相等，因此 $\sigma_G = 0$，优势 $\hat{A}_i = (r_i - \mu_G)/\sigma_G$ 未定义（除以零）。实践中，实现把所有 $\hat{A}_i = 0$，意味着\textbf{零学习信号}——这一步被浪费。

% \textbf{Prevention}:
\textbf{预防}：

\begin{enumerate}
%   \item \textbf{Temperature}: Use $\tau = 0.7$--$1.0$ during generation (not greedy).
  \item \textbf{温度}：生成时使用 $\tau = 0.7$--$1.0$（不要贪心解码）。
%   \item \textbf{Large $N$}: $N=8$--$16$ increases probability of diverse responses.
  \item \textbf{大 $N$}：$N=8$--$16$ 增加多样化响应的概率。
%   \item \textbf{Duplicate rejection}: DAPO’s approach --- reject duplicate responses and resample.
  \item \textbf{重复拒绝}：DAPO 的做法——拒绝重复响应并重采样。
%   \item \textbf{Frequency penalty}: Penalize repeated n-grams during generation.
  \item \textbf{频率惩罚}：生成时对重复的 n-gram 施加惩罚。
%   \item \textbf{Monitor}: Track unique-response ratio per group. If $<$50\%, increase temperature.
  \item \textbf{监控}：跟踪每组唯一响应比例。若 $<$50\%，调高温度。
\end{enumerate}

% \emph{\textbf{Review:} Chapter~7 (GRPO).}
\emph{\textbf{复习：}第 7 章（GRPO）。}
\end{questionbox}

% \begin{questionbox}[Q: Explain the diversity-quality tradeoff in RLHF. How do you detect mode collapse?]
\begin{questionbox}[Q：解释 RLHF 中的多样性—质量权衡。如何检测模式坍缩？]
% \textbf{Answer}: \textbf{Tradeoff}: High diversity (high entropy/temperature) = varied but potentially random/low-quality responses. Low diversity = consistent but repetitive, reward-hacked responses.
\textbf{答}：\textbf{权衡}：高多样性（高熵/高温度）= 多样但可能随机/低质量的响应。低多样性 = 一致但重复、被 Reward hack 的响应。

% \textbf{Detecting mode collapse} (all should be monitored during training):
\textbf{检测模式坍缩}（训练期间均应监控）：

\begin{enumerate}
%   \item \textbf{Response entropy}: Compute per-token entropy $H = -\sum p_i \log p_i$. If dropping rapidly $\to$ collapse.
  \item \textbf{响应熵}：计算每个 token 的熵 $H = -\sum p_i \log p_i$。若快速下降 $\to$ 坍缩。
%   \item \textbf{Unique n-gram ratio}: Fraction of unique 4-grams across responses to same prompt. Healthy: $>$0.6.
  \item \textbf{唯一 n-gram 比例}：同一 Prompt 不同响应间唯一 4-gram 的比例。健康：$>$0.6。
%   \item \textbf{Reward distribution width}: If $\sigma(\text{rewards})$ shrinks to near-zero $\to$ all responses are the same quality $\to$ likely identical.
  \item \textbf{Reward 分布宽度}：若 $\sigma(\text{rewards})$ 收缩到接近零 $\to$ 所有响应质量相同 $\to$ 可能完全一致。
%   \item \textbf{KL divergence}: If $D_\text{KL}[\pi_\theta \| \pi_\text{ref}]$ is growing rapidly, the policy is moving far from reference $\to$ often toward a narrow mode.
  \item \textbf{KL 散度}：若 $D_\text{KL}[\pi_\theta \| \pi_\text{ref}]$ 快速增长，Policy 正远离参考 $\to$ 通常朝向某个狭窄模式。
%   \item \textbf{Length histogram}: If all responses converge to same length $\to$ template behavior.
  \item \textbf{长度直方图}：若所有响应收敛到相同长度 $\to$ 模板化行为。
\end{enumerate}

% \textbf{Fix}: Increase KL coefficient $\beta$, increase entropy bonus, increase sampling temperature, or rollback to earlier checkpoint.
\textbf{修复}：提高 KL 系数 $\beta$、增大熵奖励、提高采样温度，或回滚到更早的 checkpoint。

% \emph{\textbf{Review:} Chapters~7 and~9 (GRPO; Reward Model Training).}
\emph{\textbf{复习：}第 7 章与第 9 章（GRPO；Reward 模型训练）。}
\end{questionbox}

\section{推测解码（Speculative Decoding）问题}
\label{speculative-decoding-questions}

% \begin{questionbox}[Q: Speculative decoding claims ``no quality loss.'' How can generating tokens differently produce identical output distribution?]
\begin{questionbox}[Q：推测解码声称“无质量损失”。以不同方式生成 token 怎么能产生完全相同的输出分布？]
% \textbf{Answer}: The acceptance/rejection scheme guarantees distributional equivalence:
\textbf{答}：接受/拒绝方案保证了分布等价性：

% For each draft token $\hat{x}$ with draft probability $q(\hat{x})$ and target probability $p(\hat{x})$:
对每个 draft token $\hat{x}$，给定 draft 概率 $q(\hat{x})$ 与目标概率 $p(\hat{x})$：

\begin{itemize}
%   \item Accept with probability $\min(1, p(\hat{x})/q(\hat{x}))$
  \item 以概率 $\min(1, p(\hat{x})/q(\hat{x}))$ 接受
%   \item On rejection: sample from the \emph{residual distribution} $\propto \max(0, p(x) - q(x))$
  \item 拒绝时：从\emph{残差分布} $\propto \max(0, p(x) - q(x))$ 中采样
\end{itemize}

% This is mathematically equivalent to sampling directly from $p$ (the target). Proof sketch: the probability of outputting token $x$ is $q(x) \cdot \min(1, p(x)/q(x)) + P(\text{reject}) \cdot \frac{\max(0, p(x)-q(x))}{\sum_y \max(0, p(y)-q(y))} = p(x)$.
这在数学上等价于直接从 $p$（目标）采样。证明草图：输出 token $x$ 的概率为 $q(x) \cdot \min(1, p(x)/q(x)) + P(\text{reject}) \cdot \frac{\max(0, p(x)-q(x))}{\sum_y \max(0, p(y)-q(y))} = p(x)$。

% The speedup comes from amortizing: when the draft is good (high acceptance), multiple tokens are confirmed in one target forward pass. The guarantee holds regardless of draft quality --- bad drafts just give lower speedup (more rejections), not worse quality.
加速来自摊销：当 draft 质量好（接受率高）时，一次目标模型前向就能确认多个 token。无论 draft 质量如何，这一保证始终成立——糟糕的 draft 只是带来更低的加速比（更多拒绝），而不是更差的质量。

\emph{\textbf{复习：}第 2 章（LLM 的系统基础）。}
\end{questionbox}

% \begin{questionbox}[Q: Compare Medusa vs. Eagle for speculative decoding. When would you choose each?]
\begin{questionbox}[Q：比较 Medusa 和 Eagle 在推测解码上的差异。何时选择哪一个？]
\textbf{答}：

% \textbf{Medusa}: Adds $k$ parallel prediction heads to the target model. Each head independently predicts token at position $t+i$. \emph{Pro}: No separate model, $<$1\% memory overhead. \emph{Con}: Heads predict independently --- cannot condition position $t+2$ on what was predicted at $t+1$. Acceptance rate: 60--80\%.
\textbf{Medusa}：在目标模型上增加 $k$ 个并行预测头。每个头独立预测位置 $t+i$ 的 token。\emph{优点}：无需独立模型，$<$1\% 显存开销。\emph{缺点}：各头独立预测——无法让位置 $t+2$ 以 $t+1$ 的预测为条件。接受率：60\%--80\%。

% \textbf{Eagle}: Lightweight autoregressive decoder on target model’s hidden states. Draft tokens are generated autoregressively (each conditioned on previous). \emph{Pro}: Captures inter-token dependencies $\to$ 85--95\% acceptance rate. \emph{Con}: Slightly more memory (small decoder) and sequential draft generation.
\textbf{Eagle}：在目标模型隐状态上的轻量自回归解码器。Draft token 自回归生成（每个以前一个为条件）。\emph{优点}：捕捉 token 间依赖 $\to$ 85\%--95\% 接受率。\emph{缺点}：稍多显存（小型解码器）以及顺序的 draft 生成。

% \textbf{Choose Medusa when}: Memory is extremely tight; simple integration; moderate speedup is sufficient (2--2.5$\times$).
\textbf{选 Medusa 当}：显存极度紧张；集成简单；中等加速即可（2--2.5$\times$）。

% \textbf{Choose Eagle when}: Maximum speedup needed (3--4$\times$); can afford small extra model; latency-critical single-stream generation.
\textbf{选 Eagle 当}：需要最大加速（3--4$\times$）；可承受额外的小模型；延迟敏感的单流生成。

% \textbf{Choose N-gram when}: Repetitive outputs (code, structured data); zero cost; no training needed.
\textbf{选 N-gram 当}：输出重复性强（代码、结构化数据）；零成本；无需训练。

\emph{\textbf{复习：}第 2 章（LLM 的系统基础）。}
\end{questionbox}

% \begin{questionbox}[Q: Why does speculative decoding NOT help at high batch sizes?]
\begin{questionbox}[Q：为什么推测解码在大 Batch 下\emph{帮不上忙}？]
% \textbf{Answer}: At high batch sizes ($\geq$64), autoregressive generation is already \emph{compute-efficient}: the weight-read cost is amortized across many sequences. The arithmetic intensity approaches the roofline ridge point.
\textbf{答}：在大 Batch（$\geq$64）下，自回归生成已经\emph{计算高效}：权重读取成本在多个序列间摊销。算术强度接近 roofline 拐点。

% Speculative decoding adds overhead:
推测解码会带来开销：

\begin{enumerate}
%   \item Draft generation cost (even if small model, it’s not free at high batch)
  \item Draft 生成成本（即便是小模型，大 Batch 下也不是免费的）
%   \item Verification forward pass processes $k$ extra tokens per sequence (batch $\times$ $k$ tokens)
  \item 验证前向每个序列处理 $k$ 个额外 token（Batch $\times$ $k$ 个 token）
%   \item Memory for draft model or Medusa heads
  \item Draft 模型或 Medusa 头的显存
%   \item Rejected tokens waste compute
  \item 被拒绝的 token 浪费算力
\end{enumerate}

% At batch=1 (latency-bound, memory-bound): speculation turns 1 token/step into 3--4 tokens/step --- huge win.
在 batch=1（延迟瓶颈、显存瓶颈）：推测把 1 token/步变成 3--4 token/步——巨大胜利。

% At batch=128 (already compute-efficient): the extra tokens from speculation barely help throughput because the GPU is already near-saturated. The overhead may even \emph{reduce} throughput.
在 batch=128（已经计算高效）：推测带来的额外 token 几乎不能提升吞吐，因为 GPU 已接近饱和。开销甚至可能\emph{降低}吞吐。

% \textbf{Rule}: Speculative decoding is for latency (small batch); batching is for throughput (large batch). Don’t combine them.
\textbf{规则}：推测解码用于延迟（小 Batch）；批处理用于吞吐（大 Batch）。别把它们组合在一起。

\emph{\textbf{复习：}第 2 章（LLM 的系统基础）。}
\end{questionbox}

\section{Agent 强化学习问题}
\label{agentic-rl-questions}

% \begin{questionbox}[Q: Why does standard RLHF (single-turn PPO/DPO) fail for multi-step agents?]
\begin{questionbox}[Q：为什么标准 RLHF（单轮 PPO/DPO）对多步 Agent 失败？]
% \textbf{Answer}: Standard RLHF optimizes for single-turn quality: given a prompt, produce one good response. Multi-step agents face fundamentally different challenges:
\textbf{答}：标准 RLHF 优化单轮质量：给定 Prompt，生成一个好响应。多步 Agent 面临根本不同的挑战：

\begin{enumerate}
%   \item \textbf{Credit assignment}: In a 50-step trajectory, which step caused the failure? Single-turn reward assigns credit to the entire response uniformly; multi-step needs \emph{per-step} credit.
  \item \textbf{信用分配}：在 50 步轨迹中，哪一步导致失败？单轮 Reward 把信用均匀地分配给整个响应；多步需要\emph{每一步}的信用。
%   \item \textbf{Sparse rewards}: Success/failure only at trajectory end. PPO’s GAE assumes intermediate rewards; without them, advantage estimates are noisy.
  \item \textbf{稀疏 Reward}：仅在轨迹末尾才有成功/失败。PPO 的 GAE 假设有中间 Reward；缺失时，优势估计会有噪声。
%   \item \textbf{Action space}: Actions are structured tool calls (JSON), not just token sequences. The model must learn syntax + semantics + strategy simultaneously.
  \item \textbf{动作空间}：动作是结构化的工具调用（JSON），不仅是 token 序列。模型必须同时学习语法 + 语义 + 策略。
%   \item \textbf{Non-stationarity}: The environment changes with each action (tool outputs modify state). Each step has a different ``prompt'' unlike single-turn where input is fixed.
  \item \textbf{非平稳性}：每个动作都会改变环境（工具输出修改状态）。每一步都有不同的"Prompt"，不像单轮那样输入是固定的。
%   \item \textbf{Exploration}: Agent must discover novel tool-use strategies, not just rephrase text.
  \item \textbf{探索}：Agent 必须发现新颖的工具使用策略，而不仅仅是改写文本。
\end{enumerate}

% \textbf{Solution}: Trajectory-level GRPO (rank complete trajectories), process reward models (per-step feedback), or filtered SFT on successful trajectories.
\textbf{解决方案}：轨迹级 GRPO（对完整轨迹排序）、过程奖励模型（Process Reward Model, PRM）（每一步反馈）、或对成功轨迹做过滤后的 SFT。

% \emph{\textbf{Review:} Chapter~12 (LLM Agentic Training).}
\emph{\textbf{复习：}第 12 章（LLM Agent 训练）。}
\end{questionbox}

% \begin{questionbox}[Q: Explain how GRPO is adapted for agentic training. What are the key differences from single-turn GRPO?]
\begin{questionbox}[Q：解释 GRPO 如何适配到 Agent 训练。与单轮 GRPO 的关键区别是什么？]
% \textbf{Answer}: Single-turn GRPO: generate $N$ responses to a prompt, rank by reward, compute advantages.
\textbf{答}：单轮 GRPO：对一个 Prompt 生成 $N$ 个响应，按 Reward 排序，计算优势。

% \textbf{Agentic GRPO} differences:
\textbf{Agent GRPO} 的区别：

\begin{enumerate}
%   \item \textbf{Unit of generation}: A full \emph{trajectory} (10--100 steps) instead of a single response. Each trajectory is one ``sample'' in the group.
  \item \textbf{生成单位}：完整的\emph{轨迹}（10--100 步），而非单个响应。每条轨迹是组内的一个“样本”。
%   \item \textbf{Reward}: Terminal (task success/failure) or trajectory-level (sum of step rewards). NOT per-token.
  \item \textbf{Reward}：终止 Reward（任务成功/失败）或轨迹级 Reward（步 Reward 求和）。\emph{不是}逐 token 的。
%   \item \textbf{Masking}: Only compute policy loss on the agent’s outputs (reasoning + tool calls). Mask tool \emph{outputs} (environment responses) from gradient computation.
  \item \textbf{掩码}：仅在 Agent 的输出（推理 + 工具调用）上计算 Policy Loss。把工具\emph{输出}（环境响应）从 Gradient 计算中掩掉。
%   \item \textbf{Group size}: Typically smaller ($N=4$--8) because trajectories are expensive (many forward passes per trajectory).
  \item \textbf{组大小}：通常较小（$N=4$--8），因为轨迹昂贵（每条轨迹很多次前向）。
%   \item \textbf{KL penalty}: Applied per-step to prevent drift from SFT policy at each decision point.
  \item \textbf{KL 惩罚}：在每一步应用，防止每个决策点偏离 SFT Policy。
%   \item \textbf{Length normalization}: Normalize by number of agent \emph{actions} (not tokens) to avoid penalizing thorough reasoning.
  \item \textbf{长度归一化}：按 Agent \emph{动作}数（不是 token 数）归一化，避免惩罚充分的推理。
\end{enumerate}

% \emph{\textbf{Review:} Chapters~7 and~12 (GRPO; LLM Agentic Training).}
\emph{\textbf{复习：}第 7 章与第 12 章（GRPO；LLM Agent 训练）。}
\end{questionbox}

% \begin{questionbox}[Q: Compare STaR and Reflexion and ReAct for agents.]
\begin{questionbox}[Q：比较 Agent 场景下的 STaR、Reflexion 和 ReAct。]
\textbf{答}：

% \textbf{STaR (Self-Taught Reasoner)}: Generate reasoning chains $\to$ filter by correctness $\to$ fine-tune on correct ones. \emph{Use when}: You have verifiable tasks (math, code) and want to bootstrap reasoning from a base model without RL.
\textbf{STaR（自学推理器，Self-Taught Reasoner）}：生成推理链 $\to$ 按正确性筛选 $\to$ 在正确链上微调。\emph{适用}：你有可验证任务（数学、代码），希望在不使用 RL 的情况下从基模引导出推理能力。

% \textbf{Reflexion}: After failure, generate verbal feedback (``What went wrong?'') $\to$ retry with reflection in context. No weight updates. \emph{Use when}: Inference-time improvement; limited compute for training; tasks where self-diagnosis is possible.
\textbf{Reflexion}：失败后生成自然语言反馈（“哪里出错了？”）$\to$ 把反思放进上下文重试。不更新权重。\emph{适用}：推理时改进；训练算力有限；任务允许自我诊断。

% \textbf{ReAct}: Interleave Reasoning (think) + Acting (tool use) in a structured loop. \emph{Use when}: Multi-step tool use tasks; you need transparency (reasoning traces are interpretable); the agent must decide between thinking and acting.
\textbf{ReAct}：在结构化循环中交替进行 Reasoning（思考）+ Acting（工具调用）。\emph{适用}：多步工具调用任务；需要透明性（推理轨迹可解释）；Agent 必须在思考与行动之间抉择。

% \textbf{Key differences}:
\textbf{关键差异}：

\resizebox{\textwidth}{!}{%
\begin{tabular}{@{}llll@{}}
\toprule
 & \textbf{STaR} & \textbf{Reflexion} & \textbf{ReAct} \\
\midrule
% Updates weights? & Yes (SFT) & No (in-context) & No (prompting) \\
是否更新权重？ & 是 (SFT) & 否（上下文内） & 否（Prompt） \\
% Multi-step? & No (single reasoning chain) & Yes (retry loops) & Yes (think-act cycles) \\
多步？ & 否（单条推理链） & 是（重试循环） & 是（思考—行动循环） \\
% Tools? & No & Optional & Yes (required) \\
工具？ & 否 & 可选 & 是（必需） \\
% Best for & Reasoning improvement & Error recovery & Tool-augmented tasks \\
最适合 & 推理能力提升 & 错误恢复 & 工具增强任务 \\
\bottomrule
\end{tabular}}

% \emph{\textbf{Review:} Chapters~12 and~18 (LLM Agentic Training; Agent Design Patterns).}
\emph{\textbf{复习：}第 12 章与第 18 章（LLM Agent 训练；Agent 设计模式）。}
\end{questionbox}

% \begin{questionbox}[Q: Why is GRPO preferred over PPO for a research agent?]
\begin{questionbox}[Q：为什么研究型 Agent 偏好 GRPO 而非 PPO？]
% \textbf{Answer}: For a research agent with 20--100 step trajectories:
\textbf{答}：对于具有 20--100 步轨迹的研究型 Agent：

% \textbf{PPO requires a value model}: $V(s_t)$ must predict expected total reward from the current state. For research (where state = 128K tokens of context including papers, code, and results), training an accurate value function is extremely difficult --- the value of ``having read 3 papers and written partial code'' is hard to predict.
\textbf{PPO 需要一个 Value 模型}：$V(s_t)$ 必须预测从当前状态出发的期望总 Reward。对于研究场景（其中状态 = 128K token 的上下文，包括论文、代码和结果），训练一个精确的 Value 函数极其困难——“读过 3 篇论文并写了一部分代码”的价值难以预测。

% \textbf{GRPO avoids value estimation entirely}: It generates $N$ complete trajectories per research question and uses within-group ranking as the advantage. No need to predict intermediate value --- just compare outcomes.
\textbf{GRPO 完全避免了 Value 估计}：它对每个研究问题生成 $N$ 条完整轨迹，并把组内排名作为优势。无需预测中间价值——只需比较结果。

% \textbf{Additional reasons}:
\textbf{其他原因}：

\begin{itemize}
%   \item Research quality is binary-ish (good report vs.~bad report) --- ranking is natural.
  \item 研究质量近似二元（好报告 vs.~差报告）——排序很自然。
%   \item Trajectories are long and expensive; GRPO’s $N=4$ is manageable; PPO would need many rollouts for stable value estimates.
  \item 轨迹长且昂贵；GRPO 的 $N=4$ 可承受；PPO 需要很多 Rollout 才能得到稳定的 Value 估计。
%   \item Terminal reward is sparse; GAE with sparse rewards gives noisy per-step advantages anyway.
  \item 终止 Reward 稀疏；稀疏 Reward 下 GAE 给出的每步优势本来就很有噪声。
\end{itemize}

\emph{\textbf{复习：}第 7 章与第 12 章（GRPO；LLM Agent 训练）。}
\end{questionbox}

% \begin{questionbox}[Q: Design a reward function for a coding agent. What reward hacking risks exist?]
\begin{questionbox}[Q：为编程 Agent 设计一个 Reward 函数。存在哪些 Reward hacking 风险？]
% \textbf{Answer}: \textbf{Reward design}:
\textbf{答}：\textbf{Reward 设计}：
\[
R = 0.5 \cdot R_{\text{tests}} + 0.2 \cdot R_{\text{quality}} + 0.2 \cdot R_{\text{efficiency}} + 0.1 \cdot R_{\text{safety}}
\]

\begin{itemize}
%   \item $R_{\text{tests}}$: Fraction of unit tests passing (0--1). Ground-truth verifiable.
  \item $R_{\text{tests}}$：单元测试通过比例（0--1）。可基于真值验证。
%   \item $R_{\text{quality}}$: LLM judge on code style, documentation, maintainability.
  \item $R_{\text{quality}}$：LLM 评判代码风格、文档、可维护性。
%   \item $R_{\text{efficiency}}$: $\max(0, 1 - \text{steps}/30)$ --- bonus for finishing quickly.
  \item $R_{\text{efficiency}}$：$\max(0, 1 - \text{steps}/30)$——快速完成给奖励。
%   \item $R_{\text{safety}}$: No dangerous operations (rm -rf, network access outside sandbox).
  \item $R_{\text{safety}}$：无危险操作（rm -rf、沙箱外网络访问）。
\end{itemize}

% \textbf{Reward hacking risks}:
\textbf{Reward hacking 风险}：

\begin{enumerate}
%   \item \textbf{Hardcoded outputs}: Agent learns to print expected test outputs directly without computing them. \emph{Fix}: Randomize test inputs; test on held-out cases.
  \item \textbf{硬编码输出}：Agent 学会直接打印预期测试输出而不真正计算。\emph{修复}：随机化测试输入；在保留测试用例上测试。
%   \item \textbf{Test deletion}: Agent modifies/deletes failing tests. \emph{Fix}: Sandbox tests as read-only.
  \item \textbf{删除测试}：Agent 修改/删除失败的测试。\emph{修复}：测试沙箱设为只读。
%   \item \textbf{Trivial solutions}: Agent writes minimal code that passes tests but doesn’t generalize. \emph{Fix}: Large, diverse test suites; property-based testing.
  \item \textbf{平凡解}：Agent 写出能通过测试但不泛化的最小代码。\emph{修复}：大型、多样化的测试集；基于属性的测试。
%   \item \textbf{Efficiency gaming}: Agent skips reasoning steps to maximize efficiency bonus. \emph{Fix}: Minimum quality threshold before efficiency bonus applies.
  \item \textbf{效率作弊}：Agent 跳过推理步骤以最大化效率奖励。\emph{修复}：先满足最低质量阈值，再给效率奖励。
\end{enumerate}

% \emph{\textbf{Review:} Chapters~9, 12, and~19 (Reward Model Training; LLM Agentic Training; Agentic Environments).}
\emph{\textbf{复习：}第 9、12、19 章（Reward 模型训练；LLM Agent 训练；Agent 环境）。}
\end{questionbox}

\section{列表式 Reward 与高级 RM 问题}
\label{listwise-rewards-and-advanced-rm-questions}

% \begin{questionbox}[Q: Explain the Plackett-Luce model. How does it generalize Bradley-Terry?]
\begin{questionbox}[Q：解释 Plackett-Luce 模型。它如何推广 Bradley-Terry？]
% \textbf{Answer}: Bradley-Terry models \emph{pairwise} preferences: $P(y_1 \succ y_2) = \sigma(r(y_1) - r(y_2))$.
\textbf{答}：Bradley-Terry 建模\emph{成对}偏好：$P(y_1 \succ y_2) = \sigma(r(y_1) - r(y_2))$。

% Plackett-Luce models \emph{full rankings} of $K$ items as sequential selection:
Plackett-Luce 把 $K$ 项的\emph{完整排序}建模为顺序选择：
\[
P(\pi) = \prod_{i=1}^K \frac{e^{r(y_{\pi(i)})}}{\sum_{j=i}^K e^{r(y_{\pi(j)})}}
\]
% Interpretation: Sequentially pick the best remaining item. Position 1 = softmax over all $K$; position 2 = softmax over remaining $K-1$; etc.
解释：依次挑选剩余项中最好的。位置 1 = 对全部 $K$ 做 Softmax；位置 2 = 对剩余 $K-1$ 做 Softmax；以此类推。

% \textbf{Generalization}: For $K=2$, PL reduces exactly to BT: $P(y_1 \succ y_2) = \frac{e^{r(y_1)}}{e^{r(y_1)} + e^{r(y_2)}} = \sigma(r(y_1) - r(y_2))$.
\textbf{推广关系}：当 $K=2$ 时，PL 恰好退化为 BT：$P(y_1 \succ y_2) = \frac{e^{r(y_1)}}{e^{r(y_1)} + e^{r(y_2)}} = \sigma(r(y_1) - r(y_2))$。

% \textbf{Advantage}: A ranking of $K=8$ items provides $\binom{8}{2} = 28$ implicit pairwise comparisons plus relative margin information --- much richer training signal than a single pair.
\textbf{优势}：$K=8$ 项的排序提供 $\binom{8}{2} = 28$ 个隐式成对比较，外加相对差距信息——比单一成对样本丰富得多。

% \emph{\textbf{Review:} Chapter~9 (Reward Model Training).}
\emph{\textbf{复习：}第 9 章（Reward 模型训练）。}
\end{questionbox}

% \begin{questionbox}[Q: What is a Process Reward Model (PRM) and when is it better than an Outcome Reward Model (ORM)?]
\begin{questionbox}[Q：什么是过程奖励模型（PRM），它在什么情况下优于结果奖励模型（Outcome Reward Model, ORM）？]
% \textbf{Answer}: \textbf{ORM}: Scores the final output only. $r(x, y_{\text{final}})$ = one scalar for the complete response.
\textbf{答}：\textbf{ORM}：仅对最终输出打分。$r(x, y_{\text{final}})$ = 整个响应一个标量。

% \textbf{PRM}: Scores each \emph{step} of reasoning. $r(x, y_{\text{step } t})$ = scalar per intermediate step.
\textbf{PRM}：对每一\emph{步}推理打分。$r(x, y_{\text{step } t})$ = 每一中间步一个标量。

% \textbf{PRM is better when}:
\textbf{PRM 更好的场景}：

\begin{enumerate}
%   \item \textbf{Long reasoning chains}: Math problems with 10+ steps. ORM can’t tell which step went wrong; PRM provides per-step credit assignment.
  \item \textbf{长推理链}：10+ 步的数学题。ORM 无法判断哪一步错了；PRM 提供逐步信用分配。
%   \item \textbf{Search/verification}: PRM enables tree search (beam search over reasoning steps, prune branches with low step-reward).
  \item \textbf{搜索/验证}：PRM 支持树搜索（对推理步做 beam search，剪掉步 Reward 低的分支）。
%   \item \textbf{Training signal density}: PRM gives $T$ rewards per trajectory (one per step) vs.~ORM’s single reward $\to$ lower variance advantage estimates.
  \item \textbf{训练信号密度}：PRM 每条轨迹给 $T$ 个 Reward（每步一个），而 ORM 只有一个 $\to$ 更低方差的优势估计。
\end{enumerate}

% \textbf{ORM is better when}: Tasks are short (single-turn); step boundaries are unclear; annotation cost for per-step labels is prohibitive.
\textbf{ORM 更好的场景}：任务短（单轮）；步骤边界不清晰；每步标注成本过高。

% \textbf{PRM annotation}: Can be automated via ``Math-Shepherd'' approach: for each step, complete the solution from that point multiple times. If completions from step $t$ succeed but completions from step $t+1$ fail, step $t+1$ is likely wrong.
\textbf{PRM 标注}：可通过"Math-Shepherd"方法自动化：对每一步，从该点多次补全整个解。如果从步骤 $t$ 出发的补全成功而从步骤 $t+1$ 出发的补全失败，则步骤 $t+1$ 很可能是错的。

% \emph{\textbf{Review:} Chapters~9 and~13 (Reward Model Training; RL for Large Reasoning Models).}
\emph{\textbf{复习：}第 9 章与第 13 章（Reward 模型训练；大型推理模型的 RL）。}
\end{questionbox}

\section{大型推理模型的 RL 问题}
\label{rl-for-large-reasoning-models-questions}

% \begin{questionbox}[Q: Why does DeepSeek-R1 not use a Process Reward Model despite training on long reasoning chains?]
\begin{questionbox}[Q：DeepSeek-R1 在长推理链上训练，却为何不使用过程奖励模型？]
% \textbf{Answer}: DeepSeek-R1 uses only \textbf{outcome-based rewards} (accuracy + format) for several reasons:
\textbf{答}：DeepSeek-R1 仅使用\textbf{基于结果的 Reward}（准确率 + 格式），原因如下：

\begin{enumerate}
%   \item \textbf{Verifiable tasks}: Math and code have deterministic ground-truth answers. The binary accuracy reward provides sufficient signal even for long chains.
  \item \textbf{可验证任务}：数学和代码有确定的真值答案。即便是长链，二元的准确率 Reward 也能提供足够信号。
%   \item \textbf{PRM failure modes}: Step-level reward models introduce their own reward hacking --- the model can learn to produce steps that ``look correct'' to the PRM without actually being correct.
  \item \textbf{PRM 失效模式}：步骤级 Reward 模型自身会引入 Reward hacking——模型可学到产出“在 PRM 看来正确”但实际不正确的步骤。
%   \item \textbf{GRPO’s group normalization}: By sampling $G$ completions per prompt and normalizing advantages within each group, GRPO naturally provides relative signal about which reasoning \emph{strategies} work, even without per-step rewards.
  \item \textbf{GRPO 的组归一化}：通过每个 Prompt 采样 $G$ 个补全并在组内归一化优势，GRPO 即使没有每步 Reward 也能自然提供关于哪些推理\emph{策略}有效的相对信号。
%   \item \textbf{Emergent self-correction}: With outcome-only rewards, the model learns to self-correct within chains (the ``aha moment''), which wouldn’t emerge if per-step rewards micromanage the reasoning process.
  \item \textbf{涌现的自我纠正}：在仅有结果 Reward 时，模型学会在链内自我纠正（“aha 时刻”），如果每步 Reward 对推理过程微观管理，这种行为将无法涌现。
\end{enumerate}

% \textbf{Key insight}: The verifiability of the task domain is what makes PRMs unnecessary --- for subjective tasks (creative writing), outcome-only rewards may not suffice.
\textbf{关键洞见}：任务领域的可验证性使 PRM 变得不必要——对于主观任务（创意写作），仅有结果 Reward 可能不够。

% \emph{\textbf{Review:} Chapter~13 (RL for Large Reasoning Models).}
\emph{\textbf{复习：}第 13 章（大型推理模型的 RL）。}
\end{questionbox}

% \begin{questionbox}[Q: Explain the test-time compute scaling law and its implications for model deployment]
\begin{questionbox}[Q：解释测试时计算扩展律（Test-Time Compute Scaling Law）及其对模型部署的启示]
% \textbf{Answer}: The test-time compute scaling law states:
\textbf{答}：测试时计算扩展律表述为：
\[
\text{Accuracy}(C_{\text{train}}, C_{\text{test}}) \approx f(\alpha \log C_{\text{train}} + \beta \log C_{\text{test}})
\]

% \textbf{Implications}:
\textbf{启示}：

\begin{enumerate}
%   \item \textbf{Compute equivalence}: A 7B model with 64$\times$ more inference tokens can match a 70B model with 1$\times$ tokens on reasoning tasks.
  \item \textbf{算力等价}：在推理任务上，使用 64$\times$ 推理 token 的 7B 模型可以与使用 1$\times$ token 的 70B 模型相匹敌。
%   \item \textbf{Adaptive allocation}: Easy questions get short chains (cheap); hard questions get long chains (expensive). Average cost is lower than always using a large model.
  \item \textbf{自适应分配}：简单问题用短链（便宜）；困难问题用长链（贵）。平均成本低于始终使用大模型。
%   \item \textbf{Deployment flexibility}: Instead of one large model, deploy a smaller reasoning model and scale inference compute per-query based on difficulty.
  \item \textbf{部署灵活性}：不再部署一个大模型，而是部署一个较小的推理模型，并按查询难度对推理算力进行扩展。
%   \item \textbf{Diminishing returns}: The log relationship means doubling test-time compute gives diminishing accuracy gains --- there’s an optimal allocation between training and inference compute.
  \item \textbf{收益递减}：对数关系意味着测试时计算翻倍带来的准确率提升递减——训练与推理算力之间存在一个最优分配。
\end{enumerate}

% \textbf{The ``overthinking'' failure}: Very long chains can \emph{decrease} accuracy due to error accumulation and attention dilution. Optimal chain length depends on problem difficulty.
\textbf{“过度思考”失败模式}：非常长的链可能因错误累积和 Attention 稀释而\emph{降低}准确率。最优链长取决于问题难度。

\emph{\textbf{复习：}第 13 章（大型推理模型的 RL）。}
\end{questionbox}

% \begin{questionbox}[Q: How does MCTS (Monte Carlo Tree Search) apply to LLM reasoning?]
\begin{questionbox}[Q：MCTS（蒙特卡洛树搜索，Monte Carlo Tree Search）如何应用于 LLM 推理？]
% \textbf{Answer}: MCTS for reasoning treats each partial solution as a tree node:
\textbf{答}：用于推理的 MCTS 将每个部分解视为一个树节点：

% \textbf{Four phases per iteration}:
\textbf{每次迭代四个阶段}：

\begin{enumerate}
%   \item \textbf{Selection}: Navigate from root using UCB: $\text{UCB}(s) = Q(s) + c\sqrt{\frac{\ln N(\text{parent})}{N(s)}}$
  \item \textbf{选择}：用 UCB 从根开始导航：$\text{UCB}(s) = Q(s) + c\sqrt{\frac{\ln N(\text{parent})}{N(s)}}$
%   \item \textbf{Expansion}: Generate new reasoning steps (child nodes) from the LLM
  \item \textbf{扩展}：从 LLM 生成新的推理步（子节点）
%   \item \textbf{Simulation}: Complete the solution from the new node (rollout)
  \item \textbf{模拟}：从新节点完成解（Rollout）
%   \item \textbf{Backpropagation}: Update Q-values along the path based on final correctness
  \item \textbf{反向传播}：根据最终正确性更新路径上的 Q 值
\end{enumerate}

% \textbf{Key differences from game MCTS}:
\textbf{与博弈 MCTS 的关键差异}：

\begin{itemize}
%   \item \textbf{Branching factor}: Reasoning has enormous branching factor (any next sentence is possible). Practical implementations use the LLM’s top-k outputs to limit branches.
  \item \textbf{分支因子}：推理具有巨大的分支因子（任何下一句都可能）。实际实现使用 LLM 的 top-k 输出来限制分支。
%   \item \textbf{Value function}: A trained PRM estimates partial solution quality, replacing random rollouts.
  \item \textbf{Value 函数}：训练好的 PRM 估计部分解质量，替代随机 Rollout。
%   \item \textbf{Step granularity}: Each ``step'' might be one sentence, one equation, or one logical inference --- choosing granularity matters.
  \item \textbf{步骤粒度}：每“步”可能是一句话、一个方程或一次逻辑推断——粒度选择很关键。
\end{itemize}

% \textbf{Used by}: AlphaProof (math olympiad), and hypothesized for OpenAI o1/o3’s hidden reasoning.
\textbf{使用案例}：AlphaProof（数学奥林匹克），以及据推测 OpenAI o1/o3 的隐藏推理。

\emph{\textbf{复习：}第 13 章（大型推理模型的 RL）。}
\end{questionbox}

% \begin{questionbox}[Q: Compare distillation vs direct RL for creating small reasoning models]
\begin{questionbox}[Q：比较通过蒸馏 vs 直接 RL 来打造小型推理模型]
\textbf{答}：

% \textbf{Distillation} (DeepSeek-R1-Distill approach):
\textbf{蒸馏}（DeepSeek-R1-Distill 路线）：

\begin{itemize}
%   \item Generate reasoning chains from large model (R1-671B)
  \item 从大模型（R1-671B）生成推理链
%   \item SFT small model on these chains
  \item 在这些链上对小模型做 SFT
%   \item Result: Small model mimics large model’s reasoning \emph{format}
  \item 结果：小模型模仿大模型的推理\emph{格式}
%   \item Pro: Cheap (just SFT). Con: May learn surface patterns not true reasoning.
  \item 优点：便宜（仅 SFT）。缺点：可能学到表层模式而非真正推理。
\end{itemize}

% \textbf{Direct RL} on small model:
\textbf{直接对小模型做 RL}：

\begin{itemize}
%   \item Train small model with GRPO/PPO against verifiable rewards
  \item 用 GRPO/PPO 在可验证 Reward 上训练小模型
%   \item Model discovers its own reasoning strategies
  \item 模型发现自己的推理策略
%   \item Pro: Genuine capability. Con: Much more compute; may not converge for very small models.
  \item 优点：真实能力。缺点：算力消耗大得多；对于非常小的模型可能不收敛。
\end{itemize}

% \textbf{Empirical finding}: R1-Distill-7B (distilled) outperforms direct-RL-7B on most benchmarks. The reasoning chains from the large model provide such strong supervision that SFT alone is competitive. However, distilled models show less generalization to truly novel problem types.
\textbf{经验发现}：R1-Distill-7B（蒸馏）在大多数基准上优于 direct-RL-7B。大模型的推理链提供了非常强的监督信号，单 SFT 就具有竞争力。然而，蒸馏模型对真正新颖的问题类型泛化较差。

% \textbf{Best practice}: Distill first (cheap baseline), then optionally run RL on the distilled model for further gains (``distill + RL'' combo used by Qwen).
\textbf{最佳实践}：先蒸馏（便宜的基线），再可选地在蒸馏模型上跑 RL 以获得进一步收益（Qwen 使用的“蒸馏 + RL”组合）。

\emph{\textbf{复习：}第 13 章（大型推理模型的 RL）。}
\end{questionbox}

\section{LLM 评估问题}
\label{llm-evaluation-questions}

% \begin{questionbox}[Q: Derive the ELO rating update rule and explain why Chatbot Arena uses it]
\begin{questionbox}[Q：推导 ELO 评分更新规则，并解释 Chatbot Arena 为何使用它]
% \textbf{Answer}: \textbf{ELO derivation}:
\textbf{答}：\textbf{ELO 推导}：

% Expected score of player A vs B: $E_A = \frac{1}{1 + 10^{(R_B - R_A)/400}}$ (logistic model).
玩家 A 对 B 的期望得分：$E_A = \frac{1}{1 + 10^{(R_B - R_A)/400}}$（逻辑斯蒂模型）。

% After a match with actual score $S_A \in \{0, 0.5, 1\}$: $R_A' = R_A + K(S_A - E_A)$
对局后，实际得分 $S_A \in \{0, 0.5, 1\}$：$R_A' = R_A + K(S_A - E_A)$

% The $K$-factor controls update magnitude (higher $K$ = more reactive to recent results).
$K$ 因子控制更新幅度（$K$ 越大对近期结果反应越敏感）。

% \textbf{Why Chatbot Arena uses ELO}:
\textbf{为什么 Chatbot Arena 使用 ELO}：

\begin{enumerate}
%   \item \textbf{Transitivity}: If model A beats B and B beats C, ELO predicts A beats C. This gives a total ordering from pairwise comparisons.
  \item \textbf{传递性}：若 A 胜 B 且 B 胜 C，ELO 预测 A 胜 C。这从成对比较给出全序。
%   \item \textbf{Online updates}: New models can be added without re-evaluating all pairs. Each new comparison updates ratings incrementally.
  \item \textbf{在线更新}：新增模型无需重新评估所有对。每次新比较都增量更新评分。
%   \item \textbf{Confidence}: After $N$ comparisons, rating uncertainty shrinks as $O(1/\sqrt{N})$. Standard error: $\text{SE} \approx \frac{400}{\sqrt{N}}$.
  \item \textbf{置信度}：在 $N$ 次比较后，评分不确定性以 $O(1/\sqrt{N})$ 收缩。标准误：$\text{SE} \approx \frac{400}{\sqrt{N}}$。
%   \item \textbf{Human preference capture}: Real users provide honest preferences without needing to articulate criteria. The aggregate reveals true model quality.
  \item \textbf{捕捉人类偏好}：真实用户无需阐明标准即可提供诚实偏好。聚合揭示了真实模型质量。
\end{enumerate}

% \textbf{Chatbot Arena specifics}: Uses Bradley-Terry MLE (equivalent to ELO at convergence) with bootstrap confidence intervals. Style-controlled ELO removes length/formatting bias.
\textbf{Chatbot Arena 细节}：使用 Bradley-Terry MLE（收敛时等价于 ELO），并配合 Bootstrap 置信区间。风格控制 ELO 去除长度/格式偏倚。

% \emph{\textbf{Review:} Chapter~14 (LLM Evaluation).}
\emph{\textbf{复习：}第 14 章（LLM 评估）。}
\end{questionbox}

% \begin{questionbox}[Q: What is the pass@k metric for code generation and why is the unbiased estimator important?]
\begin{questionbox}[Q：代码生成的 pass@k 指标是什么，为什么无偏估计很重要？]
% \textbf{Answer}: \textbf{pass@k} = probability that at least one of $k$ generated samples passes all test cases.
\textbf{答}：\textbf{pass@k} = 生成的 $k$ 个样本中至少有一个通过全部测试用例的概率。

% \textbf{Naive (biased) estimator}: Generate $k$ samples, check if any passes. Problem: high variance, expensive (need many trials per problem).
\textbf{朴素（有偏）估计器}：生成 $k$ 个样本，看是否有任意一个通过。问题：高方差、昂贵（每题需要很多次试验）。

% \textbf{Unbiased estimator} (Chen et al., 2021): Generate $n \geq k$ samples, count $c$ that pass:
\textbf{无偏估计器}（Chen 等，2021）：生成 $n \geq k$ 个样本，统计其中通过的数量 $c$：
\[
\text{pass@}k = 1 - \frac{\binom{n-c}{k}}{\binom{n}{k}}
\]

% \textbf{Why unbiased matters}:
\textbf{无偏为何重要}：

\begin{enumerate}
%   \item Generates $n$ samples once (e.g., $n=200$) and computes pass@1, pass@10, pass@100 from the same samples
  \item 一次性生成 $n$ 个样本（例如 $n=200$），从同一批样本算出 pass@1、pass@10、pass@100
%   \item No need to repeat the entire evaluation $k$ times
  \item 无需将整个评估重复 $k$ 次
%   \item Statistically exact (combinatorial argument: fraction of $k$-subsets with no correct sample)
  \item 统计上精确（组合论证：不含任何正确样本的 $k$-子集占比）
%   \item Numerically stable computation via log-space: $\text{pass@}k = 1 - \exp\left(\sum_{i=0}^{k-1} \log(n-c-i) - \log(n-i)\right)$
  \item 通过对数空间数值稳定计算：$\text{pass@}k = 1 - \exp\left(\sum_{i=0}^{k-1} \log(n-c-i) - \log(n-i)\right)$
\end{enumerate}

% \textbf{Intuition}: If 50/200 samples pass ($c=50$, $n=200$), pass@1 $\approx 0.25$, pass@10 $\approx 0.94$. The estimator counts what fraction of $k$-sized draws would contain at least one success.
\textbf{直觉}：若 50/200 个样本通过（$c=50$，$n=200$），pass@1 $\approx 0.25$，pass@10 $\approx 0.94$。该估计器统计大小为 $k$ 的抽取中至少包含一个成功的比例。

% \emph{\textbf{Review:} Chapters~14 and~19 (LLM Evaluation; Agentic Environments).}
\emph{\textbf{复习：}第 14 章与第 19 章（LLM 评估；Agent 环境）。}
\end{questionbox}

% \begin{questionbox}[Q: How do you detect and mitigate benchmark contamination?]
\begin{questionbox}[Q：如何检测和缓解基准污染（Benchmark Contamination）？]
% \textbf{Answer}: \textbf{Contamination}: Training data contains benchmark test examples (or close paraphrases), inflating scores.
\textbf{答}：\textbf{污染}：训练数据中包含基准测试样例（或其近似改写），使分数虚高。

% \textbf{Detection methods}:
\textbf{检测方法}：

\begin{enumerate}
%   \item \textbf{N-gram overlap}: Check if training data contains exact or near-exact matches to test items. 8-gram overlap with $>$80\% coverage is suspicious.
  \item \textbf{N-gram 重叠}：检查训练数据是否与测试条目精确或近似匹配。8-gram 重叠且覆盖率 $>$80\% 即可疑。
%   \item \textbf{Canary strings}: Insert unique identifiers in test sets; check if model can reproduce them.
  \item \textbf{Canary 字符串}：在测试集中插入唯一标识符；检查模型是否能复现它们。
%   \item \textbf{Rephrased benchmarks}: Create semantically equivalent but textually different versions of benchmarks. Large accuracy drops suggest memorization.
  \item \textbf{改写基准}：创建语义等价但文本不同的基准版本。准确率大幅下降表明存在记忆。
%   \item \textbf{Temporal analysis}: Model performance on pre-training-cutoff vs post-cutoff test items. Unusually high performance on old items suggests contamination.
  \item \textbf{时序分析}：比较模型在训练截止前 vs 截止后测试条目上的表现。在旧条目上异常高表现暗示污染。
%   \item \textbf{Membership inference}: Statistical tests for whether specific examples were in training data.
  \item \textbf{成员推断}：统计检验判断特定样例是否出现在训练数据中。
\end{enumerate}

% \textbf{Mitigation}:
\textbf{缓解}：

\begin{itemize}
%   \item \textbf{Dynamic benchmarks}: Regularly generate new test items (LiveCodeBench, Chatbot Arena)
  \item \textbf{动态基准}：定期生成新测试条目（LiveCodeBench、Chatbot Arena）
%   \item \textbf{Private test sets}: Keep test items secret (LMSYS)
  \item \textbf{私有测试集}：保持测试条目保密（LMSYS）
%   \item \textbf{Decontamination during training}: Remove detected overlaps from training data
  \item \textbf{训练期去污染}：从训练数据中删除检测到的重叠
%   \item \textbf{Report contamination analysis}: Disclose overlap metrics alongside benchmark scores
  \item \textbf{报告污染分析}：在基准分数旁披露重叠度指标
\end{itemize}

\emph{\textbf{复习：}第 14 章（LLM 评估）。}
\end{questionbox}

% \begin{questionbox}[Q: Explain position bias in LLM-as-Judge and how to mitigate it]
\begin{questionbox}[Q：解释 LLM-as-Judge 中的位置偏倚（Position Bias），以及如何缓解]
% \textbf{Answer}: \textbf{Position bias}: When using an LLM to judge two responses (A vs B), the model systematically prefers the response in a particular position (usually first or last), regardless of quality.
\textbf{答}：\textbf{位置偏倚}：当用 LLM 评判两个响应（A vs B）时，模型系统性地偏好特定位置的响应（通常是第一个或最后一个），与质量无关。

% \textbf{Empirical magnitude}: GPT-4 shows 10--15\% position bias; Claude shows 5--10\%. Smaller models show larger bias.
\textbf{经验幅度}：GPT-4 表现出 10\%--15\% 的位置偏倚；Claude 为 5\%--10\%。较小的模型偏倚更大。

% \textbf{Mitigation strategies}:
\textbf{缓解策略}：

\begin{enumerate}
%   \item \textbf{Position swapping}: Judge each pair twice (A-B and B-A). Final decision = majority. If disagreement, mark as ``tie.'' This eliminates systematic position bias but doubles cost.
  \item \textbf{位置交换}：每对评判两次（A-B 与 B-A）。最终决定 = 多数。若不一致，则标为“平局”。这能消除系统性位置偏倚，但成本翻倍。
%   \item \textbf{Multi-judge panels}: Use 3+ different models as judges. Majority vote reduces individual model biases.
  \item \textbf{多评委合议}：使用 3 个及以上不同模型作为评委。多数投票减少单一模型的偏倚。
%   \item \textbf{Reference-guided}: Provide a rubric or reference answer. Judges score each response independently against the rubric, then compare scores (eliminates pairwise comparison entirely).
  \item \textbf{参考引导}：提供评分量表或参考答案。评委按量表独立为每个响应打分，再比较分数（完全消除成对比较）。
%   \item \textbf{Calibrated prompting}: Add explicit instruction: ``The order of presentation is random and should not influence your judgment.''
  \item \textbf{校准 Prompt}：加入明确指令：“呈现顺序是随机的，不应影响你的判断。”
\end{enumerate}

% \textbf{Additional biases}: Verbosity bias (prefers longer responses), self-enhancement bias (models prefer their own outputs), authority bias (defers to responses that cite sources).
\textbf{其他偏倚}：冗长偏倚（偏好更长响应）、自我增强偏倚（模型偏好自身输出）、权威偏倚（顺从于引用来源的响应）。

\emph{\textbf{复习：}第 14 章（LLM 评估）。}
\end{questionbox}

\section{Agent 记忆问题}
\label{agentic-memory-questions}

% \begin{questionbox}[Q: Compare the four types of agentic memory and when each is critical]
\begin{questionbox}[Q：比较 Agent 的四种记忆类型，以及各自在何时至关重要]
\textbf{答}：

\begin{tabularx}{\textwidth}{@{}lXXX@{}}
\toprule
% \textbf{Type} & \textbf{What it stores} & \textbf{Access pattern} & \textbf{Critical when} \\
\textbf{类型} & \textbf{存储内容} & \textbf{访问方式} & \textbf{何时关键} \\
\midrule
% Working & Current context/scratchpad & Always in context & Complex multi-step reasoning \\
工作记忆 & 当前上下文/草稿板 & 始终在上下文中 & 复杂多步推理 \\
% Episodic & Past experiences & Retrieved by similarity & Learning from past mistakes \\
情景记忆 & 过往经验 & 按相似度检索 & 从过去错误中学习 \\
% Semantic & Facts and knowledge & Retrieved by concept & Domain-specific tasks \\
语义记忆 & 事实与知识 & 按概念检索 & 领域特定任务 \\
% Procedural & Skills and patterns & Triggered by task type & Repeated tool-use \\
程序性记忆 & 技能与模式 & 按任务类型触发 & 重复性工具调用 \\
\bottomrule
\end{tabularx}

% \textbf{Key insight}: These aren’t independent --- they interact. Episodic memory feeds semantic memory (generalizing from episodes to facts). Procedural memory is refined by episodic feedback (learning which tool sequences work). Working memory orchestrates retrieval from all other types.
\textbf{关键洞见}：它们并非独立——彼此交互。情景记忆喂养语义记忆（从事件泛化为事实）。程序性记忆被情景反馈细化（学习哪些工具序列有效）。工作记忆协调对所有其他类型的检索。

% \textbf{MemGPT analogy}: Working = hot (in-context), Episodic/Semantic = warm (vector store), Procedural = cold (archived policies). The agent itself decides when to page information in/out.
\textbf{MemGPT 类比}：工作记忆 = 热（在上下文），情景/语义记忆 = 温（向量库），程序性记忆 = 冷（归档的 Policy）。Agent 自身决定何时换入/换出信息。

% \emph{\textbf{Review:} Chapter~16 (Agentic Memory Systems).}
\emph{\textbf{复习：}第 16 章（Agent 记忆系统）。}
\end{questionbox}

% \begin{questionbox}[Q: How does temporal decay work in memory retrieval and why is it important?]
\begin{questionbox}[Q：记忆检索中的时间衰减如何工作，为何重要？]
% \textbf{Answer}: \textbf{Temporal decay} down-weights older memories during retrieval:
\textbf{答}：\textbf{时间衰减}在检索时给较旧记忆降权：
\[
\text{score}(m) = \alpha \cdot \text{similarity}(q, m) + (1 - \alpha) \cdot \text{recency}(m)
\]
% where $\text{recency}(m) = e^{-\lambda \cdot \Delta t}$ with $\Delta t$ = time since last access.
其中 $\text{recency}(m) = e^{-\lambda \cdot \Delta t}$，$\Delta t$ = 自上次访问以来的时间。

% \textbf{Why it’s important}:
\textbf{为什么重要}：

\begin{enumerate}
%   \item \textbf{Relevance decay}: User preferences change. A preference from 6 months ago may be outdated.
  \item \textbf{相关性衰减}：用户偏好会变化。6 个月前的偏好可能已过时。
%   \item \textbf{Contradiction resolution}: When old and new information conflict, recency bias naturally prefers current truth.
  \item \textbf{矛盾消解}：当新旧信息冲突时，近期偏好自然倾向当前真相。
%   \item \textbf{Retrieval efficiency}: Without decay, memory grows unbounded and retrieval returns increasingly irrelevant ancient items.
  \item \textbf{检索效率}：没有衰减时记忆无界增长，检索会返回越来越不相关的远古条目。
%   \item \textbf{Cognitive plausibility}: Humans forget too --- recent events are more accessible. This mirrors the spacing effect.
  \item \textbf{认知合理性}：人类也会遗忘——近期事件更易获取。这与间隔效应（spacing effect）相符。
\end{enumerate}

% \textbf{Access-based refresh}: When a memory is retrieved and used, its timestamp updates (similar to LRU caching). Frequently-accessed memories stay ``fresh'' regardless of creation date.
\textbf{基于访问的刷新}：当记忆被检索并使用时，其时间戳更新（类似 LRU 缓存）。频繁访问的记忆无论创建日期如何都保持“新鲜”。

% \textbf{Decay rate tuning}: $\lambda$ depends on domain. Customer service: high decay (preferences change fast). Legal/medical: low decay (facts persist). Can be learned via RL.
\textbf{衰减率调优}：$\lambda$ 取决于领域。客户服务：高衰减（偏好变化快）。法律/医疗：低衰减（事实持久）。可通过 RL 学习。

\emph{\textbf{复习：}第 16 章（Agent 记忆系统）。}
\end{questionbox}

% \begin{questionbox}[Q: How can RL be used to train memory operations?]
\begin{questionbox}[Q：如何用 RL 训练记忆操作？]
% \textbf{Answer}: Memory operations (write/read/update/delete) can be actions in the agent’s MDP:
\textbf{答}：记忆操作（写入/读取/更新/删除）可作为 Agent MDP 中的动作：

% \textbf{Formulation}:
\textbf{建模}：

\begin{itemize}
%   \item \textbf{State}: Current context + memory state
  \item \textbf{状态}：当前上下文 + 记忆状态
%   \item \textbf{Actions}: Standard actions + \texttt{memory\_write(key/value)}, \texttt{memory\_read(query)}, \texttt{memory\_delete(key)}
  \item \textbf{动作}：标准动作 + \texttt{memory\_write(key/value)}、\texttt{memory\_read(query)}、\texttt{memory\_delete(key)}
%   \item \textbf{Reward}: Task success (did memory help?) + memory efficiency penalty (fewer reads = better)
  \item \textbf{Reward}：任务成功（记忆是否有帮助？）+ 记忆效率惩罚（读取越少越好）
\end{itemize}

% \textbf{What RL learns}:
\textbf{RL 学到什么}：

\begin{enumerate}
%   \item \textbf{What to store}: Important information (API keys/user preferences) vs ephemeral details
  \item \textbf{存什么}：重要信息（API key / 用户偏好）vs 临时细节
%   \item \textbf{When to retrieve}: Before answering domain questions vs during general chat
  \item \textbf{何时检索}：回答领域问题前 vs 一般聊天中
%   \item \textbf{Compression policy}: When to summarize old memories vs keep verbatim
  \item \textbf{压缩策略}：何时对旧记忆做摘要 vs 原样保留
%   \item \textbf{Forgetting}: When old information is stale and should be removed
  \item \textbf{遗忘}：何时旧信息陈旧应被删除
\end{enumerate}

% \textbf{Training signal}: Counterfactual --- ``would the agent have succeeded if it hadn’t stored/retrieved this memory?'' Implemented via trajectory comparison: trajectories with good memory use get higher rewards.
\textbf{训练信号}：反事实——“如果 Agent 没有存储/检索这条记忆，它还会成功吗？”通过轨迹比较实现：记忆使用得当的轨迹获得更高 Reward。

% \textbf{Challenge}: Delayed reward --- storing information now may only help 100 steps later. Requires long-horizon credit assignment (GAE with high $\lambda$).
\textbf{挑战}：延迟 Reward——现在存储的信息可能 100 步后才有用。需要长期信用分配（高 $\lambda$ 的 GAE）。

\emph{\textbf{复习：}第 12 章与第 16 章（LLM Agent 训练；Agent 记忆系统）。}
\end{questionbox}

\section{Agent 编排问题}
\label{agent-orchestration-questions}

% \begin{questionbox}[Q: Explain the context budget problem and how to solve it with dynamic allocation]
\begin{questionbox}[Q：解释上下文预算问题，以及如何用动态分配解决]
% \textbf{Answer}: \textbf{The problem}: An agent has context window $L$ tokens but needs space for:
\textbf{答}：\textbf{问题}：Agent 有上下文窗口 $L$ 个 token，但需要为以下部分留空间：
\[
C = S + M + T + H + R \leq L
\]
% where $S$ = system prompt, $M$ = memory/retrieved context, $T$ = tool descriptions, $H$ = conversation history, $R$ = reserved for response.
其中 $S$ = 系统 Prompt，$M$ = 记忆/检索上下文，$T$ = 工具描述，$H$ = 对话历史，$R$ = 为响应预留。

% As conversations grow, $H$ increases and pushes out other components.
随着对话增长，$H$ 增大并挤掉其他组件。

% \textbf{Dynamic allocation strategy}:
\textbf{动态分配策略}：

\begin{enumerate}
%   \item \textbf{Fixed minimums}: $S_{\min}$, $R_{\min}$ are non-negotiable
  \item \textbf{固定下限}：$S_{\min}$、$R_{\min}$ 不可妥协
%   \item \textbf{Adaptive history}: Summarize old turns when $H > H_{\max}$. Keep last $k$ turns verbatim; summarize the rest.
  \item \textbf{自适应历史}：当 $H > H_{\max}$ 时摘要旧轮次。保留最近 $k$ 轮原样；其余做摘要。
%   \item \textbf{On-demand tools}: Only include tool descriptions relevant to current query (not all 50 tools). Use a classifier or embedding similarity to select top-$k$ tools.
  \item \textbf{按需工具}：仅包含与当前查询相关的工具描述（而非全部 50 个）。用分类器或 Embedding 相似度选 top-$k$ 工具。
%   \item \textbf{Lazy memory}: Retrieve memory only when needed (after analyzing the query) rather than pre-loading.
  \item \textbf{惰性记忆}：仅在需要时（分析查询之后）检索记忆，而非预先加载。
\end{enumerate}

% \textbf{Overflow handling}: When total exceeds $L$ even after compression:
\textbf{溢出处理}：当压缩后总量仍超过 $L$：

\begin{itemize}
%   \item Drop least-important tool descriptions
  \item 丢弃最不重要的工具描述
%   \item Aggressively summarize history to 1-sentence-per-turn
  \item 激进地把历史摘要到每轮一句
%   \item Reduce memory slots
  \item 减少记忆槽
%   \item If still over: truncate with warning to user
  \item 若仍超：截断并向用户提示
\end{itemize}

% \textbf{Pre-flight check}: Always count tokens BEFORE calling the LLM. Never discover overflow at inference time.
\textbf{飞行前检查}：调用 LLM 之前始终先统计 token 数。绝不要在推理时才发现溢出。

% \emph{\textbf{Review:} Chapter~17 (Agent Harness -- Context Management and Orchestration).}
\emph{\textbf{复习：}第 17 章（Agent Harness——上下文管理与编排）。}
\end{questionbox}

% \begin{questionbox}[Q: Compare ReAct vs Plan-and-Execute orchestration patterns]
\begin{questionbox}[Q：比较 ReAct 与 Plan-and-Execute 两种编排模式]
\textbf{答}：

% \textbf{ReAct} (Reason + Act):
\textbf{ReAct}（Reason + Act）：

\begin{itemize}
%   \item Loop: Thought $\to$ Action $\to$ Observation $\to$ Thought $\to$ \ldots{}
  \item 循环：Thought $\to$ Action $\to$ Observation $\to$ Thought $\to$ \ldots{}
%   \item Each step decides the next action based on all previous observations
  \item 每一步基于此前所有观察决定下一动作
%   \item \textbf{Pro}: Adaptive --- can change direction based on tool outputs
  \item \textbf{优点}：自适应——可根据工具输出改变方向
%   \item \textbf{Con}: Myopic --- no upfront planning; can get stuck in loops; each LLM call sees entire history (expensive)
  \item \textbf{缺点}：短视——无前置规划；可能陷入循环；每次 LLM 调用都看到完整历史（昂贵）
\end{itemize}

% \textbf{Plan-and-Execute}:
\textbf{Plan-and-Execute}：

\begin{itemize}
%   \item Phase 1: Generate full plan (list of steps)
  \item 阶段 1：生成完整计划（步骤列表）
%   \item Phase 2: Execute steps sequentially (simpler executor; possibly cheaper model)
  \item 阶段 2：顺序执行步骤（更简单的执行器；可能用更便宜的模型）
%   \item Phase 3: Re-plan if execution fails
  \item 阶段 3：执行失败则重新规划
%   \item \textbf{Pro}: Efficient (planning once is cheaper than reasoning every step); parallelizable independent steps
  \item \textbf{优点}：高效（一次规划比每步推理便宜）；独立步骤可并行
%   \item \textbf{Con}: Brittle plans --- if early steps fail the plan may be invalid. Re-planning adds latency.
  \item \textbf{缺点}：计划脆弱——若早期步骤失败，整个计划可能失效。重规划增加延迟。
\end{itemize}

% \textbf{When to use which}:
\textbf{何时选择}：

\begin{itemize}
%   \item ReAct: Exploratory tasks; unknown environments; tasks where each step’s result determines the next
  \item ReAct：探索性任务；未知环境；每步结果决定下一步的任务
%   \item Plan-and-Execute: Well-defined tasks; known tool set; parallelizable sub-tasks; cost-sensitive deployments
  \item Plan-and-Execute：定义清晰的任务；已知工具集；可并行子任务；成本敏感的部署
%   \item \textbf{Hybrid}: Plan at high level then ReAct within each plan step (LangGraph’s recommended pattern)
  \item \textbf{混合}：高层做规划，在每个规划步内用 ReAct（LangGraph 推荐模式）
\end{itemize}

% \emph{\textbf{Review:} Chapters~17 and~18 (Agent Harness; Agent Design Patterns).}
\emph{\textbf{复习：}第 17 章与第 18 章（Agent Harness；Agent 设计模式）。}
\end{questionbox}

% \begin{questionbox}[Q: How do you detect and prevent infinite loops in agent execution?]
\begin{questionbox}[Q：如何检测并防止 Agent 执行中的无限循环？]
% \textbf{Answer}: Agents can enter infinite loops when they repeat the same action expecting different results.
\textbf{答}：当 Agent 重复同一动作期待不同结果时，会进入无限循环。

% \textbf{Detection methods}:
\textbf{检测方法}：

\begin{enumerate}
%   \item \textbf{Max iteration guard}: Hard limit (e.g., 25 steps). Simple but loses work on genuinely long tasks.
  \item \textbf{最大迭代守卫}：硬上限（例如 25 步）。简单但会在真正长任务上丢失进展。
%   \item \textbf{Action hash window}: Hash the last $k$ (action/observation) pairs. If current hash matches a hash from the last $w$ steps then loop detected.
  \item \textbf{动作哈希窗口}：对最近 $k$ 个（动作/观察）对求哈希。若当前哈希与最近 $w$ 步内的哈希匹配则检测到循环。
%   \item \textbf{Semantic similarity}: Embed recent actions. If cosine similarity between consecutive actions exceeds threshold ($>$0.95) then likely stuck.
  \item \textbf{语义相似度}：对近期动作做 Embedding。若相邻动作余弦相似度超过阈值（$>$0.95）则可能卡住。
%   \item \textbf{Progress monitoring}: Define task-specific progress metrics. If no progress in $N$ steps then intervene.
  \item \textbf{进度监控}：定义任务特定的进度指标。若 $N$ 步内无进展则干预。
\end{enumerate}

% \textbf{Recovery strategies}:
\textbf{恢复策略}：

\begin{enumerate}
%   \item \textbf{Inject hint}: Add system message: ``You seem to be repeating actions. Try a different approach.''
  \item \textbf{注入提示}：加入系统消息：“你似乎在重复动作。尝试不同方法。”
%   \item \textbf{Force different action}: Mask the repeated action from the action space for the next step.
  \item \textbf{强制不同动作}：下一步在动作空间中屏蔽重复动作。
%   \item \textbf{Escalate}: Return to user with partial results and ask for guidance.
  \item \textbf{升级}：返回部分结果给用户并请求指导。
%   \item \textbf{Backtrack}: Reset to a checkpoint before the loop began and try alternative path.
  \item \textbf{回溯}：重置到循环开始前的 checkpoint 并尝试备选路径。
\end{enumerate}

% \textbf{Best practice}: Combine max iterations (safety net) + hash-based detection (early intervention) + graceful escalation (preserve user trust).
\textbf{最佳实践}：组合最大迭代（安全网）+ 基于哈希的检测（早期干预）+ 优雅升级（维护用户信任）。

\emph{\textbf{复习：}第 17 章与第 18 章（Agent Harness；Agent 设计模式）。}
\end{questionbox}

\newpage
\section{MCP 协议问题}
\label{mcp-protocol-questions}

% \begin{questionbox}[Q: Explain MCP’s N+M architecture and why it matters for the agent ecosystem]
\begin{questionbox}[Q：解释 MCP 的 N+M 架构，以及它对 Agent 生态系统为何重要]
% \textbf{Answer}: \textbf{The N$\times$M problem}: Without MCP, $N$ agent frameworks must each implement integrations with $M$ tools = $N \times M$ total integrations. Adding one new tool requires $N$ implementations.
\textbf{答}：\textbf{N$\times$M 问题}：没有 MCP 时，$N$ 个 Agent 框架必须各自实现与 $M$ 个工具的集成 = 总共 $N \times M$ 次集成。新增一个工具需要 $N$ 次实现。

% \textbf{MCP’s N+M solution}: Standardize the interface. Each agent implements one MCP client ($N$ total). Each tool implements one MCP server ($M$ total). Total integrations = $N + M$.
\textbf{MCP 的 N+M 方案}：标准化接口。每个 Agent 实现一个 MCP 客户端（共 $N$ 个）。每个工具实现一个 MCP 服务器（共 $M$ 个）。总集成数 = $N + M$。

% \textbf{Concrete example}: 5 agent frameworks (LangChain/AutoGen/CrewAI/Claude/custom) $\times$ 20 tools (GitHub/Slack/DB/filesystem/\ldots{}) = 100 integrations without MCP. With MCP: 5 clients + 20 servers = 25 implementations.
\textbf{具体示例}：5 个 Agent 框架（LangChain/AutoGen/CrewAI/Claude/custom）$\times$ 20 个工具（GitHub/Slack/DB/filesystem/\ldots{}）= 没有 MCP 时 100 次集成。有 MCP 时：5 个客户端 + 20 个服务器 = 25 次实现。

% \textbf{Why it matters}:
\textbf{为何重要}：

\begin{enumerate}
%   \item \textbf{Tool reuse}: Build a tool server once; use from any MCP-compatible agent
  \item \textbf{工具复用}：构建一次工具服务器；可被任何兼容 MCP 的 Agent 使用
%   \item \textbf{Agent portability}: Switch from Claude to a custom agent without rewriting tool integrations
  \item \textbf{Agent 可移植性}：从 Claude 切换到自定义 Agent 而无需重写工具集成
%   \item \textbf{Ecosystem growth}: Lower barrier to adding new tools incentivizes the community to build more
  \item \textbf{生态成长}：降低新增工具的门槛激励社区构建更多
%   \item \textbf{Composability}: Connect multiple servers to one agent dynamically at runtime
  \item \textbf{可组合性}：运行时将多个服务器动态连接到一个 Agent
\end{enumerate}

% \textbf{Analogy}: USB standardized peripheral connections. Before USB: every device had a proprietary connector. After USB: one port fits all. MCP does the same for agent-tool connections.
\textbf{类比}：USB 标准化了外设连接。USB 之前：每种设备都有专有接口。USB 之后：一个口适配所有。MCP 对 Agent–工具连接做了同样的事。

% \emph{\textbf{Review:} Chapter~20 (Model Context Protocol).}
\emph{\textbf{复习：}第 20 章（模型上下文协议）。}
\end{questionbox}

% \begin{questionbox}[Q: What are MCP’s four core primitives and when do you use each?]
\begin{questionbox}[Q：MCP 的四个核心原语是什么，何时分别使用？]
\textbf{答}：

\begin{tabularx}{\textwidth}{@{}lXXX@{}}
\toprule
% \textbf{Primitive} & \textbf{Direction} & \textbf{Purpose} & \textbf{Example} \\
\textbf{原语} & \textbf{方向} & \textbf{用途} & \textbf{示例} \\
\midrule
% Tools & Client $\to$ Server & Execute actions & \texttt{create\_issue}; \texttt{query\_db} \\
Tools & 客户端 $\to$ 服务器 & 执行动作 & \texttt{create\_issue}；\texttt{query\_db} \\
% Resources & Client $\to$ Server & Read data & File contents; DB records \\
Resources & 客户端 $\to$ 服务器 & 读取数据 & 文件内容；数据库记录 \\
% Prompts & Client $\to$ Server & Get templates & ``Summarize this PR'' template \\
Prompts & 客户端 $\to$ 服务器 & 获取模板 & “总结这个 PR”模板 \\
% Sampling & Server $\to$ Client & Request LLM gen & Server asks LLM to classify \\
Sampling & 服务器 $\to$ 客户端 & 请求 LLM 生成 & 服务器请求 LLM 分类 \\
\bottomrule
\end{tabularx}

% \textbf{Key distinctions}:
\textbf{关键区分}：

\begin{itemize}
%   \item \textbf{Tools vs Resources}: Tools have \emph{side effects} (create/modify/delete). Resources are \emph{read-only}. This distinction matters for safety --- an agent can freely read resources but must get approval for tools.
  \item \textbf{Tools vs Resources}：Tools 有\emph{副作用}（创建/修改/删除）。Resources 是\emph{只读}的。这对安全很重要——Agent 可自由读取 Resources，但调用 Tools 必须获得批准。
%   \item \textbf{Sampling} reverses the direction: normally the client (agent) calls the server (tool). With Sampling the server asks the client’s LLM for help. Use case: a code analysis server needs the LLM to interpret a code snippet.
  \item \textbf{Sampling} 方向反转：通常是客户端（Agent）调用服务器（工具）。在 Sampling 中，服务器请求客户端的 LLM 帮忙。用例：代码分析服务器需要 LLM 解释一段代码。
%   \item \textbf{Prompts} are metadata (reusable templates) not execution. They help the agent formulate better tool calls.
  \item \textbf{Prompts} 是元数据（可复用模板），不是执行。它们帮助 Agent 构造更好的工具调用。
\end{itemize}

\emph{\textbf{复习：}第 20 章（模型上下文协议）。}
\end{questionbox}

\section{Agent 通信（A2A）问题}
\label{agent-communication-a2a-questions}

% \begin{questionbox}[Q: How does Google’s A2A protocol differ from MCP and when do you need both?]
\begin{questionbox}[Q：Google 的 A2A 协议与 MCP 有何不同，何时两者都需要？]
% \textbf{Answer}: \textbf{Core distinction}:
\textbf{答}：\textbf{核心区分}：

\begin{itemize}
%   \item \textbf{MCP}: Agent $\leftrightarrow$ Tool (structured function calls with defined schemas)
  \item \textbf{MCP}：Agent $\leftrightarrow$ 工具（具有定义 schema 的结构化函数调用）
%   \item \textbf{A2A}: Agent $\leftrightarrow$ Agent (opaque task delegation --- you don’t know how the other agent works)
  \item \textbf{A2A}：Agent $\leftrightarrow$ Agent（不透明任务委派——你不知道对方 Agent 如何工作）
\end{itemize}

% \textbf{Key A2A concepts}:
\textbf{A2A 关键概念}：

\begin{itemize}
%   \item \textbf{Agent Cards}: JSON describing what an agent can do (like a resume). Discovery mechanism.
  \item \textbf{Agent Cards}：描述 Agent 能力的 JSON（类似简历）。发现机制。
%   \item \textbf{Opaque execution}: Requester doesn’t see internal reasoning of the delegate. Just sends task and gets result.
  \item \textbf{不透明执行}：请求者看不到受托方的内部推理。只是发送任务并获取结果。
%   \item \textbf{Task lifecycle}: submitted $\to$ working $\to$ completed/failed (with streaming updates via SSE)
  \item \textbf{任务生命周期}：submitted $\to$ working $\to$ completed/failed（通过 SSE 提供流式更新）
\end{itemize}

% \textbf{When you need both}:
\textbf{何时两者都需要}：

\begin{enumerate}
%   \item An orchestrator agent uses \textbf{A2A} to delegate ``research this topic'' to a research agent
  \item 编排 Agent 用 \textbf{A2A} 将“研究这个话题”委派给一个研究 Agent
%   \item The research agent uses \textbf{MCP} to call web search and file read and database tools
  \item 研究 Agent 用 \textbf{MCP} 调用网页搜索、文件读取和数据库工具
%   \item Results flow back via A2A to the orchestrator
  \item 结果经 A2A 流回编排 Agent
\end{enumerate}

% \textbf{Architecture}: A2A sits at the \emph{inter-agent} layer; MCP sits at the \emph{agent-tool} layer. A complete system uses both: A2A for coordination between agents and MCP for each agent’s tool access.
\textbf{架构}：A2A 位于\emph{Agent 间}层；MCP 位于\emph{Agent–工具}层。完整系统两者并用：A2A 用于 Agent 间协调，MCP 用于每个 Agent 的工具访问。

% \emph{\textbf{Review:} Chapters~20 and~22 (MCP; Agent-to-Agent Communication).}
\emph{\textbf{复习：}第 20 章与第 22 章（MCP；Agent 到 Agent 通信）。}
\end{questionbox}

% \begin{questionbox}[Q: What is the Contract Net Protocol and how does it apply to LLM agents?]
\begin{questionbox}[Q：什么是合同网协议（Contract Net Protocol），它如何应用于 LLM Agent？]
% \textbf{Answer}: The \textbf{Contract Net Protocol} (CNP) is a task allocation mechanism from distributed AI:
\textbf{答}：\textbf{合同网协议}（Contract Net Protocol，CNP）是来自分布式 AI 的任务分配机制：

% \textbf{Steps}:
\textbf{步骤}：

\begin{enumerate}
%   \item \textbf{Announce}: Manager broadcasts task description to all available agents
  \item \textbf{宣告}：管理者向所有可用 Agent 广播任务描述
%   \item \textbf{Bid}: Agents assess their capability and submit bids (confidence; estimated cost; estimated time)
  \item \textbf{投标}：Agent 评估自身能力并提交投标（置信度；预估成本；预估时间）
%   \item \textbf{Award}: Manager selects best bid(s) based on criteria (capability/cost/availability)
  \item \textbf{授标}：管理者按标准（能力/成本/可用性）选出最佳投标
%   \item \textbf{Execute}: Winning agent(s) perform the task
  \item \textbf{执行}：中标 Agent 执行任务
%   \item \textbf{Report}: Agent reports results back to manager
  \item \textbf{回报}：Agent 将结果报告给管理者
\end{enumerate}

% \textbf{For LLM agents}:
\textbf{对 LLM Agent 而言}：

\begin{itemize}
%   \item \textbf{Bidding = self-assessment}: Each agent LLM evaluates ``can I do this task well?'' and provides a confidence score. This requires calibrated self-knowledge.
  \item \textbf{投标 = 自我评估}：每个 Agent LLM 评估“我能把这个任务做好吗？”并给出置信度分数。这需要校准过的自我认知。
%   \item \textbf{Specialization emerges}: Code agents bid high on code tasks; research agents bid high on research tasks. No central routing logic needed.
  \item \textbf{专业化涌现}：代码 Agent 在代码任务上出高价；研究 Agent 在研究任务上出高价。无需中心路由逻辑。
%   \item \textbf{Load balancing}: If one agent is busy (high estimated time) others win the contract.
  \item \textbf{负载均衡}：若一个 Agent 忙（预估时间长），其他 Agent 胜出。
%   \item \textbf{Failure handling}: If awarded agent fails then re-announce to remaining agents (automatic failover).
  \item \textbf{失败处理}：若中标 Agent 失败，则向剩余 Agent 重新宣告（自动故障切换）。
\end{itemize}

% \textbf{Limitation for LLMs}: LLMs often overestimate their capabilities (hallucinate confidence). Bids should incorporate track record (historical success rate on similar tasks) not just self-reported confidence.
\textbf{LLM 的局限}：LLM 常常高估自身能力（幻觉置信度）。投标应纳入历史记录（相似任务的历史成功率），而不仅是自报置信度。

% \emph{\textbf{Review:} Chapters~22 and~23 (A2A; Multi-Agent Systems).}
\emph{\textbf{复习：}第 22 章与第 23 章（A2A；多 Agent 系统）。}
\end{questionbox}

\section{多 Agent 系统问题}
\label{multi-agent-systems-questions}

% \begin{questionbox}[Q: Compare centralized vs decentralized multi-agent architectures for LLMs]
\begin{questionbox}[Q：比较 LLM 多 Agent 系统的中心化 vs 去中心化架构]
\textbf{答}：

% \textbf{Centralized (Supervisor)}:
\textbf{中心化（Supervisor）}：

\begin{itemize}
%   \item One orchestrator LLM routes tasks to specialist workers
  \item 一个编排 LLM 将任务路由给专家 worker
%   \item Clear control flow; easy to debug (inspect supervisor decisions)
  \item 控制流清晰；易于调试（检查 supervisor 决策）
%   \item Single point of failure; supervisor becomes token bottleneck
  \item 单点故障；supervisor 成为 token 瓶颈
%   \item Best for: well-defined workflows; small agent teams (3--5 agents)
  \item 最适合：定义清晰的工作流；小型 Agent 团队（3--5 个 Agent）
\end{itemize}

% \textbf{Decentralized (Peer-to-Peer)}:
\textbf{去中心化（点对点）}：

\begin{itemize}
%   \item Agents communicate directly; no central coordinator
  \item Agent 直接通信；无中央协调者
%   \item Resilient (no single point of failure); scales horizontally
  \item 弹性（无单点故障）；可水平扩展
%   \item Hard to debug (emergent behavior); potential for conflicts and deadlocks
  \item 难以调试（涌现行为）；可能出现冲突和死锁
%   \item Communication scales $O(n^2)$ without structure
  \item 通信无结构时按 $O(n^2)$ 扩展
%   \item Best for: resilient systems; large agent populations; creative tasks where emergent behavior is desired
  \item 最适合：弹性系统；大量 Agent 群体；期望涌现行为的创造性任务
\end{itemize}

% \textbf{Hybrid (Hierarchical)}: Tree structure with sub-managers. Combines benefits: local autonomy within groups and global coordination at the top. Communication scales $O(n \log n)$.
\textbf{混合（层次化）}：带子管理者的树形结构。结合优点：组内本地自治 + 顶层全局协调。通信按 $O(n \log n)$ 扩展。

% \textbf{Decision framework}: Use centralized if you need predictability and auditability. Use decentralized if you need resilience and creativity. Use hierarchical for large ($>$10 agent) systems.
\textbf{决策框架}：需要可预测性和可审计性时用中心化。需要弹性和创造性时用去中心化。大型（$>$10 个 Agent）系统用层次化。

% \emph{\textbf{Review:} Chapter~23 (Multi-Agent Systems).}
\emph{\textbf{复习：}第 23 章（多 Agent 系统）。}
\end{questionbox}

% \begin{questionbox}[Q: What is CTDE and why is it important for training multi-agent LLM systems?]
\begin{questionbox}[Q：什么是 CTDE，为什么它对训练多 Agent LLM 系统重要？]
% \textbf{Answer}: \textbf{CTDE} = Centralized Training; Decentralized Execution.
\textbf{答}：\textbf{CTDE} = 中心化训练；去中心化执行（Centralized Training; Decentralized Execution）。

% \textbf{The problem}: In multi-agent RL each agent’s environment is non-stationary (other agents are changing their policies simultaneously). This makes independent training unstable.
\textbf{问题}：多 Agent RL 中，每个 Agent 的环境都是非平稳的（其他 Agent 同时在改变各自的 Policy）。这使独立训练不稳定。

% \textbf{CTDE solution}:
\textbf{CTDE 方案}：

\begin{itemize}
%   \item \textbf{During training}: A centralized critic has access to all agents’ observations and actions: $V(s_1, s_2, \ldots, s_n, a_1, a_2, \ldots, a_n)$. This stabilizes training by removing non-stationarity from the value function.
  \item \textbf{训练时}：一个中心化 critic 可访问所有 Agent 的观察和动作：$V(s_1, s_2, \ldots, s_n, a_1, a_2, \ldots, a_n)$。这通过把非平稳性从 Value 函数中剔除来稳定训练。
%   \item \textbf{During execution}: Each agent acts based only on its own observation: $a_i = \pi_i(o_i)$. No communication overhead at inference time.
  \item \textbf{执行时}：每个 Agent 仅基于自身观察行动：$a_i = \pi_i(o_i)$。推理时没有通信开销。
\end{itemize}

% \textbf{For LLM agents}: The centralized critic can be a reward model that evaluates the \emph{joint} output of all agents (e.g., did the team of agents produce a correct software system?) while each agent is trained to maximize its contribution to the team reward using counterfactual credit assignment.
\textbf{对 LLM Agent 而言}：中心化 critic 可以是一个 Reward 模型，评估所有 Agent 的\emph{联合}输出（例如，Agent 团队是否产出了一个正确的软件系统？），而每个 Agent 通过反事实信用分配被训练以最大化其对团队 Reward 的贡献。

% \textbf{Practical challenge}: Full CTDE requires all agents to train simultaneously with shared state --- expensive for LLMs. Approximations: train agents in rounds (freeze others and train one) or use population-based training with periodic synchronization.
\textbf{实际挑战}：完整 CTDE 要求所有 Agent 在共享状态下同时训练——对 LLM 而言代价昂贵。近似方法：分轮训练 Agent（冻结其余 Agent 训练一个），或使用具有周期同步的种群训练。

\emph{\textbf{复习：}第 23 章（多 Agent 系统）。}
\end{questionbox}

\section{Agent 开发框架问题}
\label{agent-development-framework-questions}

% \begin{questionbox}[Q: Compare LangGraph vs AutoGen vs CrewAI for building multi-agent systems]
\begin{questionbox}[Q：比较 LangGraph、AutoGen、CrewAI 用于构建多 Agent 系统]
\textbf{答}：

\begin{tabularx}{\textwidth}{@{}lXX@{}}
\toprule
% \textbf{Dimension} & \textbf{LangGraph} & \textbf{AutoGen / CrewAI} \\
\textbf{维度} & \textbf{LangGraph} & \textbf{AutoGen / CrewAI} \\
\midrule
% Orchestration & Explicit state graph (nodes + edges) & Implicit (conversation / role-based) \\
编排 & 显式状态图（节点 + 边） & 隐式（基于对话/角色） \\
% State mgmt & TypedDict schemas; checkpointing & Conversation history as state \\
状态管理 & TypedDict schema；checkpoint & 以对话历史为状态 \\
% Multi-agent & Graph with conditional routing & GroupChat / Crew \\
多 Agent & 带条件路由的图 & GroupChat / Crew \\
% Debugging & Graph visualization; step replay & Chat logs \\
调试 & 图可视化；步骤回放 & 聊天日志 \\
% HITL & First-class (interrupt nodes) & Via approval tools \\
人机协同（Human-in-the-Loop） & 一等公民（中断节点） & 通过审批工具 \\
% Production & LangGraph Cloud; persistence & Limited (AutoGen); growing (CrewAI) \\
生产部署 & LangGraph Cloud；持久化 & 有限（AutoGen）；成长中（CrewAI） \\
% Learning curve & High (graph concepts) & Low (AutoGen); Very low (CrewAI) \\
学习曲线 & 高（图概念） & 低（AutoGen）；极低（CrewAI） \\
\bottomrule
\end{tabularx}

% \textbf{Choose LangGraph when}: You need fine-grained control; complex conditional flows; production deployment with persistence and human-in-the-loop.
\textbf{选 LangGraph 当}：需要细粒度控制；复杂条件流；带持久化和人机协同的生产部署。

% \textbf{Choose AutoGen when}: Rapid prototyping of multi-agent conversations; code execution agents; research experimentation.
\textbf{选 AutoGen 当}：多 Agent 对话的快速原型；代码执行 Agent；研究实验。

% \textbf{Choose CrewAI when}: Simple role-based teams; sequential task execution; quick demos; minimal code.
\textbf{选 CrewAI 当}：简单的基于角色的团队；顺序任务执行；快速演示；代码量最小。

% \textbf{Choose none (custom) when}: You need maximum performance/control; don’t want framework lock-in; or have non-standard orchestration patterns.
\textbf{都不选（自定义）当}：需要最大性能/控制；不希望被框架锁定；或有非标准的编排模式。

% \emph{\textbf{Review:} Chapter~24 (Agent Development Frameworks).}
\emph{\textbf{复习：}第 24 章（Agent 开发框架）。}
\end{questionbox}

% \begin{questionbox}[Q: How do you test and evaluate an agent system in production?]
\begin{questionbox}[Q：如何在生产中测试和评估一个 Agent 系统？]
% \textbf{Answer}: Agent testing follows a \textbf{testing pyramid}:
\textbf{答}：Agent 测试遵循一个\textbf{测试金字塔}：

% \textbf{Level 1 --- Unit Tests} (fast; many):
\textbf{第 1 层——单元测试}（快；多）：

\begin{itemize}
%   \item Test individual tools in isolation (mock LLM; verify tool logic)
  \item 独立测试单个工具（mock LLM；验证工具逻辑）
%   \item Test prompt templates (given context; verify correct prompt construction)
  \item 测试 Prompt 模板（给定上下文；验证正确的 Prompt 构造）
%   \item Test parsers (given LLM output; verify correct extraction)
  \item 测试解析器（给定 LLM 输出；验证正确提取）
\end{itemize}

% \textbf{Level 2 --- Integration Tests} (medium speed):
\textbf{第 2 层——集成测试}（中等速度）：

\begin{itemize}
%   \item Test complete agent loops with deterministic inputs
  \item 用确定性输入测试完整的 Agent 循环
%   \item ``Golden trajectory'' tests: known-good execution traces that must reproduce
  \item “黄金轨迹”测试：必须能复现的已知良好执行轨迹
%   \item Tool chain tests: verify multi-tool sequences work end-to-end
  \item 工具链测试：验证多工具序列端到端可工作
\end{itemize}

% \textbf{Level 3 --- Behavioral Tests} (slow; few):
\textbf{第 3 层——行为测试}（慢；少）：

\begin{itemize}
%   \item Does the agent follow safety constraints? (adversarial inputs)
  \item Agent 是否遵守安全约束？（对抗性输入）
%   \item Does it ask for clarification when appropriate?
  \item 它是否在合适时请求澄清？
%   \item Does it stay within token/cost budgets?
  \item 它是否保持在 token/成本预算内？
\end{itemize}

% \textbf{Production evaluation}:
\textbf{生产评估}：

\begin{itemize}
%   \item \textbf{A/B testing}: Route 5\% of traffic to new agent version
  \item \textbf{A/B 测试}：将 5\% 流量路由到新 Agent 版本
%   \item \textbf{Shadow mode}: Run new agent alongside old; compare outputs without serving
  \item \textbf{影子模式}：与旧 Agent 并行运行新 Agent；比较输出但不对外服务
%   \item \textbf{LLM-as-judge}: Automated quality scoring of agent responses
  \item \textbf{LLM-as-judge}：自动化打分评估 Agent 响应质量
%   \item \textbf{User satisfaction}: Thumbs up/down; task completion rate; time-to-resolution
  \item \textbf{用户满意度}：点赞/点踩；任务完成率；解决时间
\end{itemize}

% \textbf{Key metric}: \textbf{Task Success Rate} (TSR) --- fraction of tasks the agent completes correctly without human intervention.
\textbf{关键指标}：\textbf{任务成功率}（Task Success Rate, TSR）——Agent 在无人干预下正确完成任务的比例。

% \emph{\textbf{Review:} Chapters~14 and~24 (LLM Evaluation; Agent Development Frameworks).}
\emph{\textbf{复习：}第 14 章与第 24 章（LLM 评估；Agent 开发框架）。}
\end{questionbox}

\section{Agent 环境问题}
\label{agentic-environments-questions}

% \begin{questionbox}[Q: Design a reward function for a web browsing agent environment]
\begin{questionbox}[Q：为网页浏览 Agent 环境设计一个 Reward 函数]
% \textbf{Answer}: For WebArena-style tasks (e.g., ``find the cheapest flight from NYC to SF on Dec 15''):
\textbf{答}：对于 WebArena 风格任务（例如“找到 12 月 15 日从 NYC 到 SF 的最便宜航班”）：

% \textbf{Sparse reward} (simple but hard to learn from):
\textbf{稀疏 Reward}（简单但难以学习）：
\[
r = \begin{cases} 1 & \text{if final page/state matches ground truth} \\ 0 & \text{otherwise} \end{cases}
\]

% \textbf{Dense reward} (better for training; harder to design):
\textbf{密集 Reward}（更利于训练；设计更难）：

\begin{enumerate}
%   \item \textbf{Progress reward}: $+0.1$ for each page that brings agent closer to goal (measured by text similarity to target state)
  \item \textbf{进度 Reward}：每个让 Agent 更接近目标的页面 $+0.1$（用与目标状态的文本相似度衡量）
%   \item \textbf{Efficiency penalty}: $-0.01$ per action (encourages shorter trajectories)
  \item \textbf{效率惩罚}：每个动作 $-0.01$（鼓励更短轨迹）
%   \item \textbf{Milestone rewards}: $+0.3$ for reaching intermediate goals (e.g., navigating to flight search page)
  \item \textbf{里程碑 Reward}：到达中间目标（如导航到航班搜索页）$+0.3$
%   \item \textbf{Invalid action penalty}: $-0.05$ for actions that produce errors (404; form validation failures)
  \item \textbf{无效动作惩罚}：产生错误的动作（404；表单校验失败）$-0.05$
\end{enumerate}

% \textbf{Potential-based shaping} (preserves optimal policy):
\textbf{基于势能的 reward shaping}（保留最优 Policy）：
\[
r_{\text{shaped}}(s, a, s') = r(s, a, s') + \gamma \Phi(s') - \Phi(s)
\]
% where $\Phi(s) = -\text{min\_steps\_to\_goal}(s)$ (estimated by heuristic or learned value function).
其中 $\Phi(s) = -\text{min\_steps\_to\_goal}(s)$（由启发式或学习到的 Value 函数估计）。

% \textbf{Challenges}: Partial observability (can’t always tell if you’re closer to goal); stochastic environments (page content changes); reward hacking (agent finds shortcuts that satisfy reward but not user intent).
\textbf{挑战}：部分可观察（无法总是判断是否更接近目标）；随机环境（页面内容会变）；Reward hacking（Agent 找到满足 Reward 但不满足用户意图的捷径）。

\emph{\textbf{复习：}第 12 章与第 19 章（LLM Agent 训练；Agent 环境）。}
\end{questionbox}

% \begin{questionbox}[Q: What makes SWE-bench a particularly challenging agent benchmark?]
\begin{questionbox}[Q：是什么让 SWE-bench 成为一个特别具有挑战性的 Agent 基准？]
% \textbf{Answer}: SWE-bench tests agents on real GitHub issues from popular Python repositories:
\textbf{答}：SWE-bench 在来自流行 Python 仓库的真实 GitHub issue 上测试 Agent：

% \textbf{Why it’s hard}:
\textbf{为什么难}：

\begin{enumerate}
%   \item \textbf{Repository-scale context}: Agent must understand codebases with 100K+ lines. Cannot fit in context window --- must explore and search and navigate.
  \item \textbf{仓库级上下文}：Agent 必须理解 10 万+ 行的代码库。无法装入上下文窗口——必须探索、搜索和导航。
%   \item \textbf{Underspecified tasks}: Issues are written by humans with implicit context. Agent must infer what’s actually needed.
  \item \textbf{规范不充分的任务}：Issue 由人类带着隐含上下文写就。Agent 必须推断真正需要什么。
%   \item \textbf{Multi-file edits}: Solutions often span multiple files with cascading dependencies.
  \item \textbf{跨文件编辑}：解决方案通常跨越多个文件，并有级联依赖。
%   \item \textbf{Test verification}: Must pass existing tests AND new tests that verify the fix.
  \item \textbf{测试验证}：必须通过现有测试 \emph{以及}验证修复的新测试。
%   \item \textbf{No hand-holding}: Unlike HumanEval (single function) SWE-bench requires full software engineering workflow: read issue $\to$ explore code $\to$ localize bug $\to$ implement fix $\to$ verify.
  \item \textbf{无人辅助}：不像 HumanEval（单函数），SWE-bench 需要完整的软件工程工作流：读 issue $\to$ 探索代码 $\to$ 定位 bug $\to$ 实现修复 $\to$ 验证。
\end{enumerate}

% \textbf{State of the art} (2024--2025): Best agents solve $\sim$50\% of SWE-bench Verified (curated subset). Full SWE-bench: $\sim$30\%.
\textbf{当前 SOTA}（2024--2025）：最佳 Agent 解决 SWE-bench Verified（精选子集）的约 $\sim$50\%。完整 SWE-bench：约 $\sim$30\%。

% \textbf{Key insight for training}: SWE-bench exposes the gap between ``coding ability'' (writing correct functions) and ``software engineering ability'' (understanding systems; navigating codebases; making minimal changes). RL training on SWE-bench-style environments teaches agents exploration and planning strategies not just code generation.
\textbf{对训练的关键洞见}：SWE-bench 揭示了“编程能力”（编写正确函数）与“软件工程能力”（理解系统；导航代码库；做最小改动）之间的差距。在 SWE-bench 风格环境上的 RL 训练教会 Agent 探索和规划策略，而不仅是代码生成。

% \emph{\textbf{Review:} Chapter~19 (Agentic Environments and Benchmarks).}
\emph{\textbf{复习：}第 19 章（Agent 环境与基准）。}
\end{questionbox}

\section{Agent UI 框架问题}
\label{agentic-ui-framework-questions}

% \begin{questionbox}[Q: Compare chat-based vs canvas-based UI paradigms for agents]
\begin{questionbox}[Q：比较 Agent 的聊天式 vs 画布式 UI 范式]
\textbf{答}：

% \textbf{Chat-based} (ChatGPT; Claude default):
\textbf{聊天式}（ChatGPT；Claude 默认）：

\begin{itemize}
%   \item Linear message stream: user $\to$ assistant $\to$ user $\to$ \ldots{}
  \item 线性消息流：用户 $\to$ 助手 $\to$ 用户 $\to$ \ldots{}
%   \item \textbf{Pro}: Familiar UX; natural for exploration and Q\&A; easy to implement
  \item \textbf{优点}：UX 熟悉；适合探索和问答；易于实现
%   \item \textbf{Con}: Generated artifacts (code/documents) are buried in conversation. Hard to iterate on a specific artifact. Context gets lost in long conversations.
  \item \textbf{缺点}：生成的产物（代码/文档）埋没在对话中。难以针对特定产物迭代。长对话中上下文易丢失。
\end{itemize}

% \textbf{Canvas/Artifact-based} (Claude Artifacts; ChatGPT Canvas; Cursor):
\textbf{画布/Artifact 式}（Claude Artifacts；ChatGPT Canvas；Cursor）：

\begin{itemize}
%   \item Side panel displays generated content; chat panel for instructions
  \item 侧边栏展示生成内容；聊天栏发送指令
%   \item Agent can create and edit and iterate on persistent artifacts
  \item Agent 可以对持久化的产物创建、编辑和迭代
%   \item \textbf{Pro}: Artifacts persist independently of chat. Direct editing by user. Version history.
  \item \textbf{优点}：产物独立于聊天持久化。用户可直接编辑。具有版本历史。
%   \item \textbf{Con}: More complex UI; requires artifact type detection; harder to implement streaming to both panels.
  \item \textbf{缺点}：UI 更复杂；需要 artifact 类型检测；对两个面板的流式输出实现更难。
\end{itemize}

% \textbf{When to use which}:
\textbf{何时选择}：

\begin{itemize}
%   \item Chat: brainstorming; Q\&A; quick tasks; mobile interfaces
  \item 聊天：头脑风暴；问答；快速任务；移动端界面
%   \item Canvas: code generation; document writing; data analysis --- any task with persistent output that needs iteration
  \item 画布：代码生成；文档写作；数据分析——任何需要迭代、有持久化输出的任务
%   \item \textbf{Hybrid} (most modern UIs): Chat by default; auto-elevate to canvas when detecting code/document/visualization output
  \item \textbf{混合}（多数现代 UI）：默认聊天；检测到代码/文档/可视化输出时自动提升到画布
\end{itemize}

% \textbf{For agent training}: The UI paradigm affects the reward signal. Canvas UIs provide explicit edit feedback (user modifies the artifact) which can be used for online learning.
\textbf{对 Agent 训练而言}：UI 范式影响 Reward 信号。画布式 UI 提供显式的编辑反馈（用户修改 artifact），可用于在线学习。

% \emph{\textbf{Review:} Chapter~25 (Agentic UI Frameworks).}
\emph{\textbf{复习：}第 25 章（Agent UI 框架）。}
\end{questionbox}

% \begin{questionbox}[Q: How do you design approval gates for human-in-the-loop agent systems?]
\begin{questionbox}[Q：如何为人机协同（Human-in-the-Loop）的 Agent 系统设计审批关卡？]
% \textbf{Answer}: Approval gates pause agent execution at critical points for human review.
\textbf{答}：审批关卡在关键点暂停 Agent 执行以供人工审核。

% \textbf{Three-tier model}:
\textbf{三层模型}：

\begin{enumerate}
%   \item \textbf{Auto-approve} (no gate): Safe reversible actions. Read operations; searches; calculations.
  \item \textbf{自动批准}（无关卡）：安全可逆动作。读取操作；搜索；计算。
%   \item \textbf{Notify} (soft gate): Potentially impactful but recoverable. Send email; create draft; modify file. Agent proceeds but user is notified and can undo.
  \item \textbf{通知}（软关卡）：可能有影响但可恢复。发送邮件；创建草稿；修改文件。Agent 继续，但用户被通知，可撤销。
%   \item \textbf{Block} (hard gate): Irreversible or high-stakes. Delete data; send money; publish content; execute code with side effects. Agent MUST wait for explicit approval.
  \item \textbf{阻断}（硬关卡）：不可逆或高风险。删数据；汇款；发布内容；执行有副作用的代码。Agent 必须等待显式批准。
\end{enumerate}

% \textbf{Design principles}:
\textbf{设计原则}：

\begin{itemize}
%   \item \textbf{Minimize interruptions}: Too many gates = user abandons the agent. The 3-tier model lets most actions flow while catching dangerous ones.
  \item \textbf{最小化打断}：关卡过多 = 用户放弃 Agent。三层模型让大多数动作通行，同时拦截危险动作。
%   \item \textbf{Show context}: At approval gate display: what action; why (agent’s reasoning); what will change; how to undo.
  \item \textbf{显示上下文}：审批关卡展示：什么动作；为什么（Agent 的推理）；将改变什么；如何撤销。
%   \item \textbf{Batch approvals}: If agent needs 5 file writes present them together not one by one.
  \item \textbf{批量审批}：如果 Agent 需要 5 次文件写入，一起呈现而不是逐个。
%   \item \textbf{Timeout handling}: If user doesn’t respond within $T$ minutes either retry notification or proceed with safe default or abort gracefully.
  \item \textbf{超时处理}：如果用户在 $T$ 分钟内未响应，则重试通知、按安全默认继续或优雅中止。
%   \item \textbf{Learning from approvals}: Track approval/rejection patterns. If users always approve a certain action type consider auto-promoting it.
  \item \textbf{从审批中学习}：跟踪审批/拒绝模式。若用户总是批准某类动作，考虑自动提升其级别。
\end{itemize}

% \textbf{Implementation}: Tool annotations (MCP’s \texttt{destructiveHint} and \texttt{readOnlyHint}) drive automatic gate assignment. Custom rules can override based on context.
\textbf{实现}：工具注解（MCP 的 \texttt{destructiveHint} 和 \texttt{readOnlyHint}）驱动关卡的自动分配。自定义规则可根据上下文覆盖。

% \emph{\textbf{Review:} Chapters~17 and~25 (Agent Harness; Agentic UI Frameworks).}
\emph{\textbf{复习：}第 17 章与第 25 章（Agent Harness；Agent UI 框架）。}
\end{questionbox}

\section{RAG 与 Agent 化 RAG 问题}
\label{rag-and-agentic-rag-questions}

% \begin{questionbox}[Q: Explain Reciprocal Rank Fusion (RRF) and why it works for hybrid retrieval]
\begin{questionbox}[Q：解释互逆排名融合（Reciprocal Rank Fusion, RRF），以及它为何对混合检索有效]
% \textbf{Answer}: RRF combines rankings from multiple retrieval systems without needing score calibration:
\textbf{答}：RRF 合并多个检索系统的排名，无需对分数做校准：
\[
\text{RRF}(d) = \sum_{r \in R} \frac{1}{k + r(d)}
\]
% where $r(d)$ is the rank of document $d$ in retriever $r$, and $k=60$ is a constant that prevents high-ranked documents from dominating.
其中 $r(d)$ 是文档 $d$ 在检索器 $r$ 中的排名，$k=60$ 是防止高排名文档主导的常数。

% \textbf{Why it works}:
\textbf{为何有效}：

\begin{enumerate}
%   \item \textbf{No score normalization needed}: BM25 scores are unbounded; dense similarity is in $[-1, 1]$. RRF uses only ranks, making them directly comparable.
  \item \textbf{无需分数归一化}：BM25 分数无界；稠密相似度在 $[-1, 1]$。RRF 只使用排名，使二者可直接比较。
%   \item \textbf{Robust to outliers}: A single retriever giving anomalously high scores doesn’t dominate because $1/(k+1) \approx 0.016$ even for rank 1.
  \item \textbf{对离群值鲁棒}：单个检索器给出异常高分也不会主导，因为即便是排名 1，$1/(k+1) \approx 0.016$。
%   \item \textbf{Complementary signals}: BM25 catches exact keyword matches; dense retrieval catches semantic similarity. Documents ranked highly by both get boosted.
  \item \textbf{互补信号}：BM25 捕捉精确关键词匹配；稠密检索捕捉语义相似度。两者都排名高的文档得到加成。
\end{enumerate}

% \textbf{Example}: Document $d$ is rank 3 in BM25 and rank 7 in dense. RRF score $= 1/(60+3) + 1/(60+7) = 0.0159 + 0.0149 = 0.0308$. A document at rank 1 in one but rank 100 in the other gets $1/61 + 1/160 = 0.0226$ --- lower despite having a top-1 ranking.
\textbf{示例}：文档 $d$ 在 BM25 中排名 3，在稠密中排名 7。RRF 分数 $= 1/(60+3) + 1/(60+7) = 0.0159 + 0.0149 = 0.0308$。在某一边排名 1 但在另一边排名 100 的文档得到 $1/61 + 1/160 = 0.0226$——尽管有排名 1，但更低。

% \textbf{In practice}: Hybrid (BM25 + dense + RRF) outperforms either alone on 85\%+ of benchmarks.
\textbf{实践中}：混合（BM25 + 稠密 + RRF）在 85\%+ 的基准上优于单独使用任一种。

% \emph{\textbf{Review:} Chapter~15 (Retrieval-Augmented Generation).}
\emph{\textbf{复习：}第 15 章（检索增强生成）。}
\end{questionbox}

% \begin{questionbox}[Q: What is Agentic RAG and how does it differ from standard RAG?]
\begin{questionbox}[Q：什么是 Agent 化 RAG，它与标准 RAG 有何不同？]
% \textbf{Answer}: \textbf{Standard RAG} follows a fixed pipeline: query $\to$ retrieve $\to$ generate. It has no ability to:
\textbf{答}：\textbf{标准 RAG} 遵循固定流水线：查询 $\to$ 检索 $\to$ 生成。它没有能力：

\begin{itemize}
%   \item Decide whether retrieval is needed at all
  \item 判断是否根本需要检索
%   \item Evaluate if retrieved documents are sufficient
  \item 评估检索到的文档是否足够
%   \item Reformulate queries when retrieval fails
  \item 在检索失败时重新表述查询
%   \item Combine information from multiple retrieval steps
  \item 整合多步检索的信息
\end{itemize}

% \textbf{Agentic RAG} treats retrieval as an \emph{action} in the agent’s MDP:
\textbf{Agent 化 RAG} 将检索视为 Agent MDP 中的一个\emph{动作}：

\begin{itemize}
%   \item \textbf{Retrieve-or-not decision}: Agent assesses if it already knows the answer (skip retrieval for factual questions in its training data)
  \item \textbf{是否检索决策}：Agent 评估它是否已经知道答案（对训练数据内的事实性问题跳过检索）
%   \item \textbf{Query planning}: Decomposes complex questions into sub-queries (``What year did X happen?'' + ``Who was president then?'')
  \item \textbf{查询规划}：把复杂问题分解为子查询（“X 发生在哪一年？” + “那时谁是总统？”）
%   \item \textbf{Self-evaluation}: After retrieval, grades relevance. If insufficient, reformulates query or tries different source.
  \item \textbf{自我评估}：检索后评估相关性。若不足，则重新表述查询或尝试不同来源。
%   \item \textbf{Multi-hop reasoning}: Retrieves $\to$ reasons $\to$ identifies knowledge gaps $\to$ retrieves again
  \item \textbf{多跳推理}：检索 $\to$ 推理 $\to$ 识别知识缺口 $\to$ 再次检索
%   \item \textbf{Source routing}: Routes queries to appropriate knowledge bases (web for current events; internal docs for company info; code search for programming)
  \item \textbf{来源路由}：将查询路由到合适的知识库（时事查网络；公司信息查内部文档；编程查代码搜索）
\end{itemize}

% \textbf{Key architectural difference}: Standard RAG = deterministic pipeline. Agentic RAG = state machine with conditional transitions (LangGraph pattern with retrieve/grade/rewrite/generate nodes).
\textbf{关键架构差异}：标准 RAG = 确定性流水线。Agent 化 RAG = 带条件转移的状态机（LangGraph 的 retrieve/grade/rewrite/generate 节点模式）。

% \textbf{Trade-off}: Agentic RAG is more accurate on complex queries but adds latency (multiple LLM calls for routing/grading). Use standard RAG for simple factual lookups; agentic RAG for multi-hop or ambiguous queries.
\textbf{权衡}：Agent 化 RAG 在复杂查询上更准确，但增加延迟（多次 LLM 调用用于路由/评分）。简单事实查询用标准 RAG；多跳或歧义查询用 Agent 化 RAG。

% \emph{\textbf{Review:} Chapters~15 and~17 (RAG; Agent Harness).}
\emph{\textbf{复习：}第 15 章与第 17 章（RAG；Agent Harness）。}
\end{questionbox}

% \begin{questionbox}[Q: Compare Self-RAG and CRAG approaches to improving retrieval quality]
\begin{questionbox}[Q：比较 Self-RAG 与 CRAG 两种改进检索质量的方法]
\textbf{答}：

% \textbf{Self-RAG} (Asai et al., 2023):
\textbf{Self-RAG}（Asai 等，2023）：

\begin{itemize}
%   \item Trains special \emph{reflection tokens} into the LLM vocabulary
  \item 将特殊的\emph{反思 token}训练进 LLM 词表
%   \item At inference, model outputs tokens like [Retrieve], [IsRel], [IsSup], [IsUse]
  \item 推理时，模型输出 [Retrieve]、[IsRel]、[IsSup]、[IsUse] 等 token
%   \item Model decides \emph{when} to retrieve (not every query needs it)
  \item 模型决定\emph{何时}检索（并非每个查询都需要）
%   \item After retrieval, model self-grades: Is the retrieved passage relevant? Does my answer follow from it?
  \item 检索后，模型自我评分：检索段落是否相关？我的答案是否从中得出？
%   \item \textbf{Training}: SFT on data augmented with reflection labels from GPT-4
  \item \textbf{训练}：在用 GPT-4 反思标签增强的数据上做 SFT
%   \item \textbf{Pro}: Single model handles everything. \textbf{Con}: Requires custom training.
  \item \textbf{优点}：单一模型处理一切。\textbf{缺点}：需要定制训练。
\end{itemize}

% \textbf{CRAG} (Corrective RAG, Yan et al., 2024):
\textbf{CRAG}（纠正式 RAG，Corrective RAG，Yan 等，2024）：

\begin{itemize}
%   \item Uses a lightweight \emph{retrieval evaluator} (separate model) to grade retrieved docs
  \item 使用一个轻量的\emph{检索评估器}（独立模型）来给检索到的文档评分
%   \item Three actions based on confidence: \texttt{Correct} (use as-is), \texttt{Ambiguous} (augment with web search), \texttt{Incorrect} (discard; fallback to web)
  \item 根据置信度有三种动作：\texttt{Correct}（按原样使用）、\texttt{Ambiguous}（用网络搜索增强）、\texttt{Incorrect}（丢弃；回退到网络）
%   \item Adds a \emph{knowledge refinement} step: extract only relevant sentences from retrieved docs
  \item 增加\emph{知识精炼}步骤：仅从检索文档中抽取相关句子
%   \item \textbf{Pro}: Works with any frozen LLM. \textbf{Con}: Extra model for evaluation; added latency.
  \item \textbf{优点}：可与任何冻结 LLM 配合。\textbf{缺点}：需要额外的评估模型；增加延迟。
\end{itemize}

% \textbf{Key difference}: Self-RAG embeds retrieval decisions into the LLM itself (requires training). CRAG is a pipeline approach that wraps around any LLM (no training needed). Self-RAG is more elegant; CRAG is more practical for production with existing models.
\textbf{关键差异}：Self-RAG 把检索决策嵌入 LLM 自身（需要训练）。CRAG 是包裹在任意 LLM 周围的流水线方法（无需训练）。Self-RAG 更优雅；CRAG 在现有模型的生产环境中更实用。

\emph{\textbf{复习：}第 15 章（检索增强生成）。}
\end{questionbox}

% \begin{questionbox}[Q: What is the lost-in-the-middle problem and how do you mitigate it?]
\begin{questionbox}[Q：什么是“中间迷失”（lost-in-the-middle）问题，如何缓解？]
% \textbf{Answer}: \textbf{The problem}: When retrieved context is long (many passages), LLMs disproportionately attend to information at the \emph{beginning} and \emph{end} of the context, ignoring information in the middle. If the answer is in passage 5 of 10, the model may miss it.
\textbf{答}：\textbf{问题}：当检索到的上下文很长（很多段落）时，LLM 不成比例地关注上下文的\emph{开头}和\emph{结尾}，而忽略中间的信息。如果答案在 10 段中的第 5 段，模型可能漏掉。

% \textbf{Empirical evidence}: Liu et al. (2023) showed that for 20-document retrieval, accuracy drops by 15--20\% when the relevant document is in positions 5--15 vs positions 1--3.
\textbf{经验证据}：Liu 等（2023）表明，对于 20 文档检索，相关文档位于位置 5--15 时，相比位置 1--3，准确率下降 15\%--20\%。

% \textbf{Mitigation strategies}:
\textbf{缓解策略}：

\begin{enumerate}
%   \item \textbf{Re-rank and truncate}: Use a cross-encoder to re-rank, then only include top-3 most relevant passages (fewer = less lost-in-middle).
  \item \textbf{重排序并截断}：用 cross-encoder 重排序，然后只保留 top-3 最相关段落（更少 = 更少中间迷失）。
%   \item \textbf{Strategic ordering}: Place highest-relevance passages at the start AND end of context, low-relevance in the middle.
  \item \textbf{策略性排序}：把最相关的段落同时放在上下文的开头和结尾，低相关的放中间。
%   \item \textbf{Contextual compression}: Summarize each passage to 1--2 sentences before insertion. Less text = less position bias.
  \item \textbf{上下文压缩}：插入前将每段摘要到 1--2 句。文本越少 = 位置偏倚越小。
%   \item \textbf{Map-reduce}: Process each passage independently (map), then combine answers (reduce). Eliminates position effects entirely.
  \item \textbf{Map-Reduce}：独立处理每段（map），再合并答案（reduce）。完全消除位置效应。
%   \item \textbf{Citation prompting}: Ask model to cite which passage it used. This forces attention to all passages.
  \item \textbf{引用 Prompt}：要求模型引用其使用的段落。这迫使 Attention 关注所有段落。
%   \item \textbf{Chunk size reduction}: Smaller chunks mean fewer total chunks needed to cover the answer.
  \item \textbf{减小块大小}：更小的块意味着覆盖答案所需的总块数更少。
\end{enumerate}

% \textbf{Best practice}: Retrieve many (20+), re-rank to top 3--5, order by relevance (best first). This sidesteps the problem entirely for most use cases.
\textbf{最佳实践}：检索很多（20+），重排序到 top 3--5，按相关性排序（最佳在前）。这在大多数用例中完全绕开此问题。

\emph{\textbf{复习：}第 15 章（检索增强生成）。}
\end{questionbox}

\chapter{速查手册}
\label{quick-reference}

% This chapter consolidates key equations, architecture specifications, API references, and failure mode diagnostics for rapid lookup during development and debugging.
本章汇总了关键公式、架构规格、API 参考以及失效模式诊断，便于开发与调试时快速查阅。

\section{核心 RL 与对齐公式}
\label{core-rl-alignment-equations}

\begin{align}
\text{PPO Clip:}&\quad L = \mathbb{E}[\min(r_t\hat{A}_t, \text{clip}(r_t,1{\pm}\epsilon)\hat{A}_t)], \quad r_t = \pi_\theta(a_t|s_t)/\pi_{\text{old}}(a_t|s_t) \\
\text{DPO:}&\quad L = -\mathbb{E}[\log\sigma(\beta\log\tfrac{\pi_\theta(y_w|x)}{\pi_\text{ref}(y_w|x)} - \beta\log\tfrac{\pi_\theta(y_l|x)}{\pi_\text{ref}(y_l|x)})] \\
\text{GRPO:}&\quad \hat{A}_i = (r_i - \mu_G)/\sigma_G, \quad \text{then PPO clip update (no critic)} \\
\text{KTO:}&\quad L = \lambda_w(1 - v(y_w)) + \lambda_l \cdot v(y_l), \quad v = \sigma(\beta\log(\pi_\theta/\pi_\text{ref}) - z) \\
\text{IPO:}&\quad L = \mathbb{E}[(\log(\pi_\theta(y_w)/\pi_\text{ref}(y_w)) - \log(\pi_\theta(y_l)/\pi_\text{ref}(y_l)) - 1/(2\beta))^2] \\
\text{ORPO:}&\quad L = L_\text{SFT}(y_w) - \lambda\log\sigma(\log(\text{odds}(y_w)/\text{odds}(y_l))) \\
\text{GAE:}&\quad \hat{A}_t = \textstyle\sum_{l=0}^{T-t}(\gamma\lambda)^l\delta_{t+l}, \quad \delta_t = r_t + \gamma V(s_{t+1}) - V(s_t) \\
\text{KL Penalty:}&\quad R_\text{total} = r_\phi(x,y) - \beta D_\text{KL}[\pi_\theta(y|x)\|\pi_\text{ref}(y|x)] \\
\text{RM (Bradley-Terry):}&\quad L = -\mathbb{E}[\log\sigma(r_\phi(x,y_w)-r_\phi(x,y_l))] \\
\text{Best-of-N:}&\quad y^* = \arg\max_{y_i \sim \pi_\theta(\cdot|x),\, i=1..N} r_\phi(x, y_i)
\end{align}

\section{Transformer 与架构公式}
\label{transformer-architecture-formulas}

\begin{align}
\text{Self-Attention:}&\quad \text{Attn}(Q,K,V) = \text{softmax}(QK^\top / \sqrt{d_k}) \cdot V \\
\text{Multi-Head:}&\quad \text{MHA}(X) = \text{Concat}(\text{head}_1, \ldots, \text{head}_h)W^O,\quad \text{head}_i = \text{Attn}(XW_i^Q, XW_i^K, XW_i^V) \\
\text{RoPE:}&\quad f(x_m, m) = x_m e^{im\theta_j}, \quad \theta_j = 10000^{-2j/d} \\
\text{LoRA:}&\quad W' = W_0 + (\alpha/r) \cdot BA, \quad B \in \mathbb{R}^{d \times r},\; A \in \mathbb{R}^{r \times k} \\
\text{KD (soft targets):}&\quad L_\text{KD} = (1{-}\alpha)L_\text{CE}(y, \hat{y}) + \alpha\, T^2 \cdot \text{KL}(p_T^\text{teacher} \| p_T^\text{student}) \\
\text{FFN (SwiGLU):}&\quad \text{FFN}(x) = (\text{Swish}(xW_1) \odot xW_3) W_2
\end{align}

\section{解码方法}
\label{decoding-methods}

\begin{tabularx}{\textwidth}{@{}l X X@{}}
\toprule
\textbf{方法} & \textbf{公式 / 规则} & \textbf{关键参数} \\
\midrule
Greedy & $y_t = \arg\max_v P(v|y_{<t})$ & --- \\
Beam search & 按联合概率保留 top-$B$ 个部分序列 & $B=4$--$8$ \\
Temperature & $P'(v) = \text{softmax}(\text{logit}_v / T)$ & $T \in [0.1, 1.5]$ \\
Top-$k$ & 仅保留 top-$k$ 个 logit，其余置零后重新归一化 & $k=40$--$100$ \\
Top-$p$ (nucleus) & 保留最小集合 $V'$ 使 $\sum_{v \in V'} P(v) \geq p$ & $p=0.9$--$0.95$ \\
Min-$p$ & 保留满足 $P(v) \geq p_\text{min} \cdot P(v_\text{max})$ 的 token & $p_\text{min}=0.05$--$0.1$ \\
Repetition penalty & 若 $v$ 已出现，则 $\text{logit}_v \leftarrow \text{logit}_v / \theta$ & $\theta=1.1$--$1.3$ \\
\bottomrule
\end{tabularx}

\section{系统与并行}
\label{systems-parallelism}

\begin{tabularx}{\textwidth}{@{}l X X@{}}
\toprule
\textbf{公式} & \textbf{取值 (70B, BF16)} & \textbf{说明} \\
\midrule
模型显存 & $2P$ 字节 & $140$ GB（仅权重） \\
Adam 优化器 & $2P \times 4$ 字节 (m + v) & $280$ GB \\
完整训练占用 & $\sim 8P$ 字节 & $560$ GB（权重 + 优化器 + 梯度） \\
FSDP 每 GPU 显存 & $8P / N_\text{GPUs}$ & 8 卡时为 $70$ GB \\
生成算术强度 & $2P / 2P = 1$ FLOP/byte & 严重内存瓶颈 \\
Token 速率（生成） & HBM\_BW $/ (2P)$ & $\sim$14 tok/s (A100, batch=1) \\
TP AllReduce / 层 & $2 \times 2 \cdot \frac{T-1}{T} \cdot bsd$ 字节 & $\sim$188 MB (70B, TP=8) \\
PP 气泡占比 & $(P-1)/(P+M-1)$ & $P$=阶段数，$M$=微批数 \\
MFU & observed\_toks $\times$ 6$P$ / peak\_FLOPS & 目标：$>40\%$ \\
\bottomrule
\end{tabularx}

\section{GPU 硬件规格}
\label{gpu-hardware-specs}

{\footnotesize
\begin{tabularx}{\textwidth}{@{}l l l l l X@{}}
\toprule
\textbf{GPU} & \textbf{显存} & \textbf{带宽 (HBM)} & \textbf{BF16 TFLOPS} & \textbf{NVLink} & \textbf{备注} \\
\midrule
A100-80GB & 80 GB HBM2e & 2.0 TB/s & 312 & 600 GB/s & 主力型号，广泛可用 \\
H100-80GB & 80 GB HBM3 & 3.35 TB/s & 989 & 900 GB/s & 当前一代，支持 FP8 \\
H200-141GB & 141 GB HBM3e & 4.8 TB/s & 989 & 900 GB/s & 大上下文 / 更少 GPU \\
B200 & 192 GB HBM3e & 8.0 TB/s & 2250 & 1800 GB/s & 下一代（2025） \\
\bottomrule
\end{tabularx}
}

\section{超参数范围}
\label{hyperparameter-ranges}

\begin{tabularx}{\textwidth}{@{}l X X X@{}}
\toprule
\textbf{参数} & \textbf{典型范围} & \textbf{默认值} & \textbf{说明} \\
\midrule
$\beta$ (DPO/KTO) & 0.05--0.5 & 0.1 & 越大越保守 \\
$\epsilon$ (PPO clip) & 0.1--0.3 & 0.2 & 越大更新越激进 \\
$\gamma$ (GAE 折扣) & 0.99--1.0 & 1.0 & 情景式任务使用 1.0 \\
$\lambda$ (GAE) & 0.9--0.99 & 0.95 & 越小偏差越大、方差越小 \\
KL 系数 ($\beta_\text{KL}$) & 0.01--0.2 & 0.05 & 自适应目标 KL $\approx$ 5--8 \\
学习率 (RLHF) & 1e-7 -- 5e-6 & 5e-7 & 远低于预训练 \\
学习率 (SFT) & 1e-5 -- 5e-5 & 2e-5 & 标准微调范围 \\
LoRA 秩 $r$ & 8--128 & 16--64 & 越大容量越大、显存越多 \\
LoRA alpha $\alpha$ & $r$ -- $2r$ & $2r$ & 缩放因子；$\alpha/r$ 为有效尺度 \\
温度（生成） & 0.6--1.0 & 0.7 & 越低候选越同质 \\
生成数量 $K$ & 4--64 & 4--16 & 用于 GRPO / Online DPO / Best-of-N \\
梯度裁剪范数 & 0.5--2.0 & 1.0 & 防止梯度爆炸 \\
\bottomrule
\end{tabularx}

\section{TRL API 速查}
\label{trl-api-quick-reference}

\begin{tabularx}{\textwidth}{@{}l X X X@{}}
\toprule
\textbf{Trainer} & \textbf{方法} & \textbf{关键配置} & \textbf{数据格式} \\
\midrule
\texttt{SFTTrainer} & 监督微调 & \texttt{packing, max\_seq\_length} & prompt + completion \\
\texttt{RewardTrainer} & 奖励模型 & \texttt{center\_rewards\_coefficient} & prompt + chosen + rejected \\
\texttt{PPOTrainer} & PPO & \texttt{init\_kl\_coef, target\_kl, cliprange} & prompts（在线生成） \\
\texttt{DPOTrainer} & DPO / IPO & \texttt{beta, loss\_type="sigmoid"/"ipo"} & prompt + chosen + rejected \\
\texttt{GRPOTrainer} & GRPO & \texttt{num\_generations, beta, use\_vllm} & prompts + reward\_fn \\
\texttt{OnlineDPOTrainer} & Online DPO & \texttt{num\_generations, reward\_model\_path} & prompts（在线生成） \\
\texttt{KTOTrainer} & KTO & \texttt{desirable\_weight, undesirable\_weight} & prompt + completion + label \\
\texttt{ORPOTrainer} & ORPO & \texttt{beta} & prompt + chosen + rejected \\
\texttt{Best-of-N (manual)} & Best-of-N & \texttt{n\_samples} & prompts（推理） \\
\bottomrule
\end{tabularx}

\section{RAG 流水线公式}
\label{rag-pipeline-formulas}

\begin{align}
\text{Cosine similarity:}&\quad \text{sim}(q, d) = \frac{q \cdot d}{\|q\| \cdot \|d\|} \\
\text{Retrieval:}&\quad \mathcal{D}_k = \text{top-}k_{d \in \mathcal{C}} \; \text{sim}(\text{embed}(q),\; \text{embed}(d)) \\
\text{RAG generation:}&\quad P(y|q) = P_\text{LLM}(y \;|\; q, \mathcal{D}_k) \\
\text{Chunking overlap:}&\quad \text{stride} = \text{chunk\_size} - \text{overlap} \\
\text{Reranker (cross-enc):}&\quad \text{score}(q, d) = \text{MLP}(\text{BERT}([q; d]))
\end{align}

\section{智能体设计模式}
\label{agentic-design-patterns}

\begin{tabularx}{\textwidth}{@{}l X X@{}}
\toprule
\textbf{模式} & \textbf{结构} & \textbf{最适用于} \\
\midrule
ReAct & Think $\to$ Act $\to$ Observe $\to$ 循环 & 通用工具使用型 Agent \\
Plan-and-Execute & Plan $\to$ 执行步骤 $\to$ 修订 & 长程、结构化任务 \\
Supervisor & 路由器 $\to$ 专家 Agent & 多领域、子任务边界清晰 \\
Swarm (handoffs) & Agent 移交控制权 + 上下文 & 客服、问题升级流程 \\
Hierarchical & 委派 Agent 的树状结构 & 复杂任务分解 \\
Human-in-the-loop & Agent $\to$ 审批关卡 $\to$ 继续 & 高风险、不可逆动作 \\
\bottomrule
\end{tabularx}

\section{智能体通信协议}
\label{agent-communication-protocols}

\begin{tabularx}{\textwidth}{@{}l X X X@{}}
\toprule
\textbf{协议} & \textbf{范围} & \textbf{传输方式} & \textbf{核心概念} \\
\midrule
MCP & 工具集成 & stdio / HTTP+SSE & 服务端暴露工具；客户端发现并调用 \\
A2A & 智能体到智能体 & HTTP + JSON-RPC & 带生命周期的任务（submitted$\to$working$\to$done） \\
OpenAI Function Calling & 工具使用 & API 负载 & 在 \texttt{tools[]} 数组中以 JSON schema 描述 \\
\bottomrule
\end{tabularx}

\section{上下文窗口预算}
\label{context-window-budget}

\begin{equation}
C \geq \underbrace{S}_{\text{system}} + \underbrace{M}_{\text{memory/RAG}} + \underbrace{T}_{\text{tool defs}} + \underbrace{H}_{\text{history}} + \underbrace{R}_{\text{reserved output}}
\end{equation}

% \textbf{Rule of thumb} for 128K context:
\textbf{经验法则}（针对 128K 上下文）：

\begin{itemize}
  % \item System prompt: 1--4K tokens (fixed)
  \item 系统提示：1--4K token（固定）
  % \item Tool definitions: 2--8K (scales with \# tools)
  \item 工具定义：2--8K（随工具数量增长）
  % \item RAG context: 4--16K (top-$k$ chunks)
  \item RAG 上下文：4--16K（top-$k$ 个 chunk）
  % \item History: grows unbounded $\rightarrow$ summarize/truncate
  \item 历史记录：无界增长 $\rightarrow$ 汇总或截断
  % \item Reserved output: 2--8K
  \item 预留输出：2--8K
\end{itemize}

\section{常见失效模式与修复}
\label{common-failure-modes-fixes}

\resizebox{\textwidth}{!}{%
\begin{tabular}{@{}lll@{}}
\toprule
\textbf{症状} & \textbf{可能原因} & \textbf{修复方法} \\
\midrule
奖励上升、质量下降 & 奖励黑客（Reward Hacking） & RM 集成、长度惩罚、增大 $\beta$ \\
KL 爆炸（$>$15） & 学习率过高或模式崩溃 & 降低学习率、回滚 checkpoint \\
熵塌缩 & 过早收敛 & 增大熵系数、提高温度 \\
训练损失 NaN & 梯度爆炸 & 降低学习率、增大梯度裁剪、检查数据 \\
5K 步后无提升 & Prompt 分布不当 & Goldilocks 过滤（通过率 20--80\%） \\
基准性能回退 & 对齐税（Alignment Tax） & 减小 RL 预算、改用 LoRA、混入 SFT 数据 \\
长度单调增长 & RM 中的长度漏洞 & 长度惩罚、用长度控制重训 RM \\
生成时 OOM & KV cache 溢出 & 减小 batch、增大 TP、使用 PagedAttention \\
Agent 永久循环 & 缺少最大迭代守卫 & 设置 \texttt{max\_iterations}、加循环检测 \\
工具调用解析失败 & 输出格式不一致 & 少样本示例、约束解码 \\
RAG 返回无关文档 & embedding / chunking 不佳 & Reranker、混合检索、更小 chunk \\
多智能体死锁 & 循环依赖 & DAG 约束、每个 Agent 设置超时 \\
\bottomrule
\end{tabular}}

\section{方法选择决策树}
\label{method-selection-decision-tree}

\begin{enumerate}
  % \item \textbf{Have paired preferences (chosen + rejected)?}
  \item \textbf{是否有成对偏好数据（chosen + rejected）？}

\begin{itemize}
  % \item Noisy labels $\rightarrow$ \textbf{IPO}
  \item 标签噪声大 $\rightarrow$ \textbf{IPO}
  % \item Memory-constrained, no SFT done yet $\rightarrow$ \textbf{ORPO}
  \item 显存受限、尚未做 SFT $\rightarrow$ \textbf{ORPO}
  % \item Clean data, limited compute $\rightarrow$ \textbf{DPO}
  \item 数据干净、算力有限 $\rightarrow$ \textbf{DPO}
  % \item DPO plateaus, want exploration $\rightarrow$ \textbf{Online DPO}
  \item DPO 进入瓶颈、希望加入探索 $\rightarrow$ \textbf{Online DPO}
\end{itemize}
  % \item \textbf{Have only binary feedback (thumbs up/down)?} $\rightarrow$ \textbf{KTO}
  \item \textbf{只有二值反馈（点赞 / 点踩）？} $\rightarrow$ \textbf{KTO}
  % \item \textbf{Have verifiable rewards (math/code)?} $\rightarrow$ \textbf{GRPO}
  \item \textbf{有可验证奖励（数学 / 代码）？} $\rightarrow$ \textbf{GRPO}
  % \item \textbf{Need maximum quality, any cost?} $\rightarrow$ \textbf{PPO}
  \item \textbf{不惜代价追求最高质量？} $\rightarrow$ \textbf{PPO}
  % \item \textbf{Want training-free improvement?} $\rightarrow$ \textbf{Best-of-N}
  \item \textbf{想要免训练的改进？} $\rightarrow$ \textbf{Best-of-N}
\end{enumerate}

\section{评估指标}
\label{evaluation-metrics}

\begin{tabularx}{\textwidth}{@{}l X X@{}}
\toprule
\textbf{指标} & \textbf{范围} & \textbf{衡量内容} \\
\midrule
Perplexity & $[1, \infty)$ & 模型的“困惑度”；越低语言建模越好 \\
Win Rate（对比基线） & $[0, 1]$ & 评审/人类偏好的输出比例 \\
BLEU & $[0, 1]$ & 与参考文本的 $n$-gram 重叠（偏精确率） \\
ROUGE-L & $[0, 1]$ & 与参考文本的最长公共子序列 \\
Pass@$k$ & $[0, 1]$ & $k$ 个代码样本中至少 1 个通过测试的概率 \\
MMLU / GPQA & $[0, 1]$ & 知识/推理基准上的多选准确率 \\
HumanEval & $[0, 1]$ & 生成代码的功能正确率 \\
Faithfulness (RAG) & $[0, 1]$ & 受检索上下文支持的论断比例 \\
Context Relevancy & $[0, 1]$ & 检索内容中与查询相关的比例 \\
Answer Relevancy & $[0, 1]$ & 答案对问题的回应程度 \\
\bottomrule
\end{tabularx}

\section{推理与测试时扩展}
\label{reasoning-test-time-scaling}

\begin{tabularx}{\textwidth}{@{}l X X@{}}
\toprule
\textbf{方法} & \textbf{算力开销} & \textbf{机制} \\
\midrule
思维链（Chain-of-Thought，CoT） & 1.5--3$\times$ token & 在 prompt 中加入 “step by step” \\
自洽性（Self-Consistency） & $N \times$ 生成 & 采样 $N$ 条 CoT 路径，对最终答案多数表决 \\
思维树（Tree-of-Thought, ToT） & $B \times D \times$ 生成 & 在推理树上 BFS/DFS；对分支评估 \\
Best-of-$N$ & $N \times$ 生成 & 采样 $N$，用 RM 打分，取最高 \\
Beam search（推理上） & $B \times$ 生成 & 维持 top-$B$ 个部分推理链 \\
预算强制（Budget forcing） & 可变 & 动态为更难问题分配更多 token \\
验证（ORM/PRM） & $N \times$ 生成 + 打分 & 生成 $N$ 个解，按结果/过程 RM 排序 \\
\bottomrule
\end{tabularx}

\section{记忆系统类型}
\label{memory-system-types}

\begin{tabularx}{\textwidth}{@{}l X X@{}}
\toprule
\textbf{类型} & \textbf{存储} & \textbf{使用场景} \\
\midrule
工作记忆 & 上下文窗口 & 当前对话、即时工具结果 \\
情景记忆（Episodic Memory） & 向量存储 & 历史交互、用户偏好、会话历史 \\
语义记忆（Semantic Memory） & 知识图谱 / embedding & 事实、概念、领域知识 \\
程序性记忆（Procedural Memory） & 技能库 / 代码 & 操作手册、学到的工作流 \\
\bottomrule
\end{tabularx}

\section{MCP 速查}
\label{mcp-quick-reference}

\begin{tabularx}{\textwidth}{@{}l X X X@{}}
\toprule
\textbf{原语} & \textbf{方向} & \textbf{有副作用？} & \textbf{用途} \\
\midrule
Tools & 客户端 $\to$ 服务端 & 是 & 执行动作（创建、修改、删除） \\
Resources & 客户端 $\to$ 服务端 & 否（只读） & 读取数据（文件、DB 记录、配置） \\
Prompts & 客户端 $\to$ 服务端 & 否 & 常用任务的可复用模板 \\
Sampling & 服务端 $\to$ 客户端 & 否 & 服务端请求客户端调用 LLM 生成 \\
\bottomrule
\end{tabularx}

% \textbf{Transport}: \texttt{stdio} (local subprocess) or \texttt{HTTP+SSE} (remote, streamable).\\
\textbf{传输方式}：\texttt{stdio}（本地子进程）或 \texttt{HTTP+SSE}（远程、可流式）。\\

% \textbf{Discovery}: Client calls \texttt{tools/list}, \texttt{resources/list}, \texttt{prompts/list} at connection init.\\
\textbf{发现机制}：客户端在连接初始化时调用 \texttt{tools/list}、\texttt{resources/list}、\texttt{prompts/list}。\\

% \textbf{Tool annotations}: \texttt{readOnlyHint}, \texttt{destructiveHint}, \texttt{idempotentHint}, \texttt{openWorldHint}.
\textbf{工具注解}：\texttt{readOnlyHint}、\texttt{destructiveHint}、\texttt{idempotentHint}、\texttt{openWorldHint}。

\section{A2A 协议速查}
\label{a2a-protocol-quick-reference}

\begin{tabularx}{\textwidth}{@{}l X@{}}
\toprule
\textbf{概念} & \textbf{描述} \\
\midrule
Agent Card & 位于 \texttt{/.well-known/agent.json} 的 JSON --- 名称、技能、支持的内容类型 \\
Task & 工作单元：\texttt{id}、\texttt{status}、\texttt{artifacts}。生命周期：submitted $\to$ working $\to$ completed/failed \\
Message & 任务内的通信单元（角色：user/agent；parts：text/file/data） \\
Artifact & Agent 产出的输出（结构化数据、文件、生成内容） \\
Push Notifications & 长任务的 webhook 更新（通过 \texttt{tasks/pushNotification/set}） \\
\bottomrule
\end{tabularx}

% \textbf{Key endpoints}: \texttt{tasks/send} (create/update), \texttt{tasks/get} (poll status), \texttt{tasks/sendSubscribe} (SSE stream).
\textbf{关键端点}：\texttt{tasks/send}（创建/更新）、\texttt{tasks/get}（轮询状态）、\texttt{tasks/sendSubscribe}（SSE 流）。

\section{智能体框架对比}
\label{agent-framework-comparison}

\begin{tabularx}{\textwidth}{@{}l X X X@{}}
\toprule
\textbf{框架} & \textbf{编排方式} & \textbf{多 Agent} & \textbf{最适用于} \\
\midrule
LangGraph & 显式状态图 & 条件路由 & 生产环境：持久化、HITL、精细控制 \\
OpenAI Agents SDK & 声明式 handoff & 基于 handoff & 简单易用：护栏、追踪、快速上手 \\
AutoGen (AG2) & 对话驱动 & GroupChat & 原型开发：代码执行、研究场景 \\
CrewAI & 基于角色的团队 & 顺序 / 并行 & 低代码：快速演示、简单流水线 \\
Google ADK & 会话 + 事件 & 原生 A2A & 企业级：artifact 管理、多模态 \\
\bottomrule
\end{tabularx}

\section{智能体 RL 公式}
\label{agentic-rl-formulas}

\begin{align}
\text{Trajectory GRPO:}&\quad \hat{A}_i = (R(\tau_i) - \mu_G)/\sigma_G, \quad R(\tau_i) = \sum_{t} r_t^{(\tau_i)} \\
\text{Agent reward:}&\quad R = w_1 R_\text{task} + w_2 R_\text{efficiency} + w_3 R_\text{safety}, \quad R_\text{eff} = \max(0, 1 - \text{steps}/N_\text{max}) \\
\text{Masking:}&\quad \mathcal{L} = \sum_{t \in \text{agent tokens}} \min(r_t \hat{A}_t,\; \text{clip}(r_t) \hat{A}_t) \quad \text{(mask env outputs)} \\
\text{Pass@}k:&\quad 1 - \frac{\binom{n-c}{k}}{\binom{n}{k}}, \quad n = \text{total samples},\; c = \text{correct}
\end{align}

\section{智能体安全检查清单}
\label{agent-security-checklist}

\begin{tabularx}{\textwidth}{@{}l X X@{}}
\toprule
\textbf{威胁} & \textbf{层级} & \textbf{缓解措施} \\
\midrule
Prompt 注入（直接） & 输入 & 输入校验、指令层级、分隔符 \\
Prompt 注入（间接） & 工具输出 & 将工具输出视为不可信；勿遵循检索文档中的指令 \\
工具滥用 & 执行 & 最小权限；\texttt{destructiveHint} 关卡；沙箱化 \\
数据外泄 & 输出 & 输出过滤；限制工具仅访问白名单域 \\
过度自主 & 架构 & 最大迭代次数；成本预算；人工审批关卡 \\
混淆代理（Confused Deputy） & 多 Agent & 校验任务来源；基于能力的访问控制 \\
\bottomrule
\end{tabularx}

\section{智能体评估指标}
\label{agent-evaluation-metrics}

\begin{tabularx}{\textwidth}{@{}l X X@{}}
\toprule
\textbf{指标} & \textbf{公式 / 定义} & \textbf{目标} \\
\midrule
任务成功率（TSR） & 正确完成数 / 总任务数 & $>85\%$（生产） \\
完成步数 & 单次成功任务的平均动作数 & 越低越高效 \\
单任务成本 & 总 token 数 $\times$ 单价 & 视预算而定 \\
延迟（TTFC） & 从请求到首条有效输出 & 交互场景 $<5$ 秒 \\
工具调用准确率 & 正确工具选择 / 总调用 & $>90\%$ \\
恢复率 & 重试成功 / 初次失败 & $>60\%$ \\
人工升级率 & 需要人工的任务 / 总任务 & $<15\%$ \\
\bottomrule
\end{tabularx}

\section{关键智能体基准}
\label{key-agentic-benchmarks}

\begin{tabularx}{\textwidth}{@{}l X X X@{}}
\toprule
\textbf{基准} & \textbf{领域} & \textbf{指标} & \textbf{SOTA (2025)} \\
\midrule
SWE-bench Verified & 软件工程 & 解决 issue 百分比 & $\sim$70\% \\
WebArena & Web 浏览 & 任务成功率 & $\sim$40\% \\
OSWorld & 桌面计算机使用 & 任务成功率 & $\sim$25\% \\
GAIA & 通用 AI 助手 & 完全匹配准确率 & $\sim$75\%（L1） \\
Tau-bench & 工具使用可靠性 & 通过率（5 次试验） & $\sim$65\% \\
HumanEval / MBPP & 代码生成 & Pass@1 & $>95\%$ \\
\bottomrule
\end{tabularx}

\chapter{总结与未来方向}
\label{ch:conclusion}

\section{总结}
\label{summary-dup2}

% This guide has traced the full arc from transformer foundations through reinforcement learning for alignment to the construction of autonomous agentic systems. The key themes that emerge across all chapters:
本书完整勾勒了从 Transformer 基础、面向对齐的强化学习，到自主智能体系统构建的全过程。跨越各章节，浮现出以下关键主题：

\begin{enumerate}
  % \item \textbf{Alignment is a systems problem.} It is not enough to have a good loss function. Production RLHF requires managing 4+ models, distributing computation across hundreds of GPUs, handling fault tolerance, and monitoring for reward hacking---all simultaneously.
  \item \textbf{对齐是一个系统工程问题。} 仅有一个好的损失函数并不够。生产级 RLHF 需要同时管理 4 个以上的模型、在数百块 GPU 上分布计算、处理容错，并监控奖励黑客行为。
  % \item \textbf{There is no single best method.} PPO remains the gold standard for maximum quality but demands enormous engineering investment. DPO and its variants offer compelling trade-offs for teams with limited infrastructure. GRPO bridges the gap for verifiable-reward domains. The right choice depends on your data, compute budget, and quality bar.
  \item \textbf{不存在单一最优方法。} PPO 仍是追求极致质量的金标准，但需要巨大的工程投入。DPO 及其变体则为基础设施有限的团队提供了极具吸引力的折中。GRPO 弥合了可验证奖励领域的鸿沟。正确的选择取决于你的数据、算力预算与质量门槛。
  % \item \textbf{Reasoning emerges from reward.} DeepSeek-R1 proved that chain-of-thought, self-verification, and backtracking can emerge from simple binary reward signals and group-relative optimization---without explicit demonstrations of reasoning. Test-time compute scaling means smaller models with more thinking can match larger models.
  \item \textbf{推理能力可以从奖励中涌现。} DeepSeek-R1 证明，思维链、自我验证和回溯能力可以从简单的二值奖励信号与组相对优化中涌现，无需对推理过程进行显式示范。测试时算力扩展意味着“更小的模型 + 更多思考”可以匹敌更大模型。
  % \item \textbf{Standards unlock ecosystems.} MCP reduces the tool integration problem from $N \times M$ to $N + M$. A2A enables agents built by different teams to collaborate without shared internals. These protocols are to agentic AI what HTTP was to the web---the enabling infrastructure for an open ecosystem.
  \item \textbf{标准开启生态。} MCP 将工具集成问题从 $N \times M$ 降到 $N + M$。A2A 让不同团队构建的 Agent 在不共享内部实现的前提下协作。这些协议之于智能体 AI，正如 HTTP 之于 Web——是开放生态的基础设施。
  % \item \textbf{Agents are the natural next step.} Once a model is aligned, the frontier shifts from ``how good is a single response?'' to ``can the model solve multi-step problems autonomously?'' This requires new training paradigms (agentic RL with environment rewards), new infrastructure (harnesses, tool protocols, memory systems), and new evaluation methods (trajectory-level benchmarks).
  \item \textbf{Agent 是自然的下一步。} 一旦模型被对齐，前沿就从“单次回复有多好？”转向“模型能否自主解决多步问题？”这需要新的训练范式（带环境奖励的智能体 RL）、新的基础设施（运行框架、工具协议、记忆系统）以及新的评估方法（轨迹级基准）。
  % \item \textbf{Evaluation drives everything.} Without rigorous evaluation---from reward model validation to agent task success rates, from contamination detection to LLM-as-Judge calibration---progress is unmeasurable and regressions are invisible. The benchmarks you choose shape the systems you build.
  \item \textbf{评估驱动一切。} 没有严格的评估——从奖励模型验证到 Agent 任务成功率，从污染检测到 LLM-as-Judge 校准——进步无法度量、回退也不可见。你所选择的基准塑造了你所构建的系统。
  % \item \textbf{Simplicity scales.} The most reliable production agents use the simplest architecture that meets requirements---prompt chaining and routing before autonomous loops, single agents before multi-agent swarms. Complexity should be earned through demonstrated need.
  \item \textbf{简单才能扩展。} 最可靠的生产级 Agent 采用满足需求的最简单架构——优先使用 prompt chaining 与 routing，而非自主循环；优先使用单 Agent，而非多 Agent 集群。复杂度只能通过被证实的需要去赢得。
\end{enumerate}

\section{前路：开放挑战}
\label{the-road-ahead-open-challenges}

\subsection{从交互中学习}
\label{learning-from-interaction}

% Current RLHF pipelines~\cite{ouyang2022training} treat alignment as a one-time training phase. The future points toward \textbf{continuous learning from deployment}: agents that improve from every user interaction, tool failure, and environment observation---without catastrophic forgetting~\cite{kirkpatrick2017overcoming} or reward drift. Key open problems:
当前的 RLHF 流水线~\cite{ouyang2022training} 将对齐视为一次性的训练阶段。未来则指向\textbf{从部署中持续学习}：Agent 能够从每一次用户交互、工具失败与环境观察中改进——同时避免灾难性遗忘（Catastrophic Forgetting）~\cite{kirkpatrick2017overcoming} 与奖励漂移。关键开放问题：

\begin{itemize}
  % \item Online learning with non-stationary reward distributions.
  \item 非平稳奖励分布下的在线学习。
  % \item Safe exploration in production~\cite{garcia2015comprehensive} (avoiding harmful actions while learning).
  \item 生产环境中的安全探索~\cite{garcia2015comprehensive}（在学习过程中避免有害动作）。
  % \item Efficient credit assignment over long agent trajectories (hundreds of tool calls).
  \item 在长 Agent 轨迹（数百次工具调用）上的高效信用分配。
\end{itemize}

\subsection{可扩展监督}
\label{scalable-oversight}

% As agents become more capable, human oversight becomes the bottleneck. Current approaches (RLHF~\cite{ouyang2022training}, Constitutional AI~\cite{bai2022constitutional}) rely on humans evaluating model outputs---but what happens when model outputs exceed human understanding?
当 Agent 能力越来越强时，人类监督本身成为瓶颈。当前方法（RLHF~\cite{ouyang2022training}、Constitutional AI~\cite{bai2022constitutional}）依赖人类评估模型输出——但当模型输出超出人类理解能力时会发生什么？

\begin{itemize}
  % \item \textbf{Recursive reward modeling}~\cite{christiano2017deep}: Use AI to help humans evaluate AI.
  \item \textbf{递归奖励建模（Recursive Reward Modeling）}~\cite{christiano2017deep}：用 AI 协助人类评估 AI。
  % \item \textbf{Debate and amplification}~\cite{irving2018debate}: Two models argue; a human judges which argument is more compelling.
  \item \textbf{辩论与放大（Debate and Amplification）}~\cite{irving2018debate}：两个模型相互辩论，人类裁判哪一方论证更令人信服。
  % \item \textbf{Process-based supervision}~\cite{lightman2023lets}: Reward correct reasoning steps, not just final answers.
  \item \textbf{基于过程的监督（Process-based Supervision）}~\cite{lightman2023lets}：奖励正确的推理步骤，而不仅是最终答案。
  % \item \textbf{Mechanistic interpretability}~\cite{olsson2022context}: Understand \emph{what} the model is doing internally, not just what it outputs.
  \item \textbf{机制可解释性（Mechanistic Interpretability）}~\cite{olsson2022context}：理解模型内部\emph{在做什么}，而不仅是它输出了什么。
\end{itemize}

\subsection{世界模型与规划}
\label{world-models-and-planning}

% Current agents are reactive---they observe and respond one step at a time. Future agents will need \textbf{internal world models}~\cite{hafner2020dream} that enable lookahead planning:
当前 Agent 是反应式的——它们一次只观察并响应一步。未来 Agent 将需要\textbf{内部世界模型（World Models）}~\cite{hafner2020dream} 以支持前瞻式规划：

\begin{itemize}
  % \item Predicting the consequences of actions before executing them.
  \item 在执行动作之前预测其后果。
  % \item Tree search over possible action sequences (à la AlphaGo~\cite{silver2016mastering} and MuZero~\cite{schrittwieser2020mastering} but for open-ended tasks).
  \item 对可能动作序列进行树搜索（类似 AlphaGo~\cite{silver2016mastering} 与 MuZero~\cite{schrittwieser2020mastering}，但面向开放任务）。
  % \item Learning environment dynamics from interaction traces.
  \item 从交互轨迹中学习环境的状态转移。
\end{itemize}

\subsection{多智能体生态}
\label{multi-agent-ecosystems}

% The A2A protocol~\cite{google-a2a-2025} and multi-agent frameworks hint at a future where hundreds of specialized agents collaborate, negotiate, and delegate---forming an ``economy of agents''~\cite{nisan2007algorithmic}. Open challenges:
A2A 协议~\cite{google-a2a-2025} 与多智能体框架预示着这样一种未来：成百上千的专用 Agent 协作、协商并相互委派——形成所谓的“智能体经济（Economy of Agents）”~\cite{nisan2007algorithmic}。开放挑战包括：

\begin{itemize}
  % \item Trust and verification between agents with different principals.
  \item 服务于不同委托方的 Agent 之间的信任与验证。
  % \item Emergent cooperation vs.~emergent deception in competitive settings~\cite{hubinger2024sleeper}.
  \item 竞争性场景中涌现的合作与涌现的欺骗的对比~\cite{hubinger2024sleeper}。
  % \item Market mechanisms for resource allocation (compute, tool access, priority).
  \item 资源分配（算力、工具访问、优先级）的市场机制。
  % \item Governance: who is responsible when a chain of 10 agents produces a harmful outcome?~\cite{amodei2016concrete}
  \item 治理：当一条由 10 个 Agent 组成的链路产生有害结果时，由谁负责？~\cite{amodei2016concrete}
\end{itemize}

\subsection{智能体安全与信任}
\label{agent-security-and-trust}

% Autonomous agents inherit every security vulnerability of the LLMs they are built on---plus new attack surfaces created by tool access, multi-agent delegation, and persistent memory (Chapters~19--21). Critical unsolved problems:
自主 Agent 继承了其底层 LLM 的全部安全漏洞，再加上工具访问、多智能体委派与持久化记忆所带来的新攻击面（第 19--21 章）。关键的未解问题：

\begin{itemize}
  % \item \textbf{Prompt injection at scale}~\cite{greshake2023indirect}: As agents consume untrusted content (web pages, emails, API responses), indirect prompt injection becomes systemic. No robust defense exists today.
  \item \textbf{规模化的 Prompt 注入}~\cite{greshake2023indirect}：随着 Agent 消费不可信内容（网页、邮件、API 响应），间接 prompt 注入成为系统性问题。目前尚无稳健的防御手段。
  % \item \textbf{Confused deputy attacks}: An agent with legitimate credentials can be tricked into misusing them on behalf of an attacker embedded in the data stream~\cite{anthropic-mcp-2024}.
  \item \textbf{混淆代理（Confused Deputy）攻击}：持有合法凭证的 Agent 可能被骗，代表潜伏在数据流中的攻击者滥用其凭证~\cite{anthropic-mcp-2024}。
  % \item \textbf{Sandboxing without crippling}: Least-privilege execution constrains what agents can do, but overly restrictive sandboxes negate agentic value. Finding the right boundary is an open design problem.
  \item \textbf{不削弱能力的沙箱化}：最小权限执行限制了 Agent 能做的事，但过度严格的沙箱会抹杀智能体的价值。如何找到合适边界仍是开放设计问题。
  % \item \textbf{Audit and attribution}: When an agent chain spans multiple organizations (via A2A~\cite{google-a2a-2025}), tracing \emph{who} authorized \emph{what} action remains architecturally unsolved.
  \item \textbf{审计与归因}：当 Agent 链跨越多个组织（通过 A2A~\cite{google-a2a-2025}）时，追溯\emph{谁}授权了\emph{什么}动作在架构上仍未解决。
  % \item \textbf{Trust calibration}: Agents must learn when \emph{not} to trust---whether a tool response is authentic, whether another agent's claim is verified.
  \item \textbf{信任校准}：Agent 必须学会何时\emph{不要}信任——工具响应是否真实、其他 Agent 的声明是否经过验证。
\end{itemize}

\subsection{超越基准的评估}
\label{evaluation-beyond-benchmarks}

% Chapter~14 showed that benchmarks shape the systems we build---yet current evaluation has critical gaps:
第 14 章揭示了基准塑造了我们所构建的系统——然而当前评估仍存在关键缺口：

\begin{itemize}
  % \item \textbf{Real-world deployment metrics}: Benchmarks like SWE-bench~\cite{jimenez2024swebench} and GAIA~\cite{mialon2023gaia} measure isolated tasks; production agents face ambiguous goals, shifting requirements, and multi-turn recovery.
  \item \textbf{真实部署指标}：SWE-bench~\cite{jimenez2024swebench} 与 GAIA~\cite{mialon2023gaia} 等基准衡量的是孤立任务；而生产 Agent 面对的是模糊目标、变动的需求与多轮恢复。
  % \item \textbf{Reward model validity}: RLHF assumes reward models capture human preferences, but reward hacking~\cite{skalse2022defining} and distributional shift undermine this assumption at scale.
  \item \textbf{奖励模型有效性}：RLHF 假设 RM 能捕捉人类偏好，但奖励黑客~\cite{skalse2022defining} 与分布漂移会在规模化时侵蚀这一假设。
  % \item \textbf{Cost-quality frontiers}: Two agents may achieve the same accuracy, but one costs 10$\times$ more tokens. Evaluation must become cost-aware.
  \item \textbf{成本—质量前沿}：两个 Agent 可能取得相同准确率，但其中一个的 token 成本高出 10$\times$。评估必须开始关注成本。
  % \item \textbf{Safety under distribution shift}: An agent safe in testing may behave unsafely on novel inputs. Adversarial evaluation~\cite{perez2022red} and red-teaming at agentic scale remain immature.
  \item \textbf{分布漂移下的安全性}：在测试中安全的 Agent 在新输入下可能行为不安全。智能体规模上的对抗评估~\cite{perez2022red} 与红队测试仍不成熟。
\end{itemize}

\subsection{效率与可获得性}
\label{efficiency-and-accessibility}

% Training a 70B model with RLHF costs $10K--$100K. Running autonomous agents costs $1--$50 per complex task. For agentic AI to achieve broad impact:
对 70B 模型做 RLHF 训练成本在 \$10K--\$100K。运行一个自主 Agent 在每个复杂任务上的成本是 \$1--\$50。要让智能体 AI 产生广泛影响，需要：

\begin{itemize}
  % \item Distillation of agentic capabilities from large to small models~\cite{hinton2015distilling, kim2024distillation}.
  \item 将智能体能力从大模型蒸馏到小模型~\cite{hinton2015distilling, kim2024distillation}。
  % \item More efficient RL algorithms (fewer samples, lower variance)~\cite{schulman2017proximal}.
  \item 更高效的 RL 算法（更少样本、更低方差）~\cite{schulman2017proximal}。
  % \item On-device agents that operate without cloud round-trips.
  \item 无需云端往返即可运行的端侧 Agent。
  % \item Open-weight models that match proprietary quality for agentic tasks~\cite{deepseek2025r1}.
  \item 在智能体任务上达到闭源质量的开放权重模型~\cite{deepseek2025r1}。
\end{itemize}

\section{延伸阅读}
\label{further-reading}

\subsection{基础论文}
\label{foundational-papers}

\begin{itemize}
  % \item \textbf{Attention Is All You Need}~\cite{vaswani2017attention} --- The transformer architecture.
  \item \textbf{Attention Is All You Need}~\cite{vaswani2017attention} --- Transformer 架构。
  % \item \textbf{RLHF / InstructGPT}~\cite{ouyang2022training} --- The first large-scale RLHF deployment.
  \item \textbf{RLHF / InstructGPT}~\cite{ouyang2022training} --- 首个大规模 RLHF 部署。
  % \item \textbf{PPO}~\cite{schulman2017proximal} --- Proximal Policy Optimization.
  \item \textbf{PPO}~\cite{schulman2017proximal} --- 近端策略优化（Proximal Policy Optimization）。
  % \item \textbf{DPO}~\cite{rafailov2023direct} --- Direct Preference Optimization.
  \item \textbf{DPO}~\cite{rafailov2023direct} --- 直接偏好优化（Direct Preference Optimization）。
  % \item \textbf{GRPO / DeepSeek-R1}~\cite{shao2024deepseekmath, deepseek2025r1} --- Group Relative Policy Optimization and emergent reasoning.
  \item \textbf{GRPO / DeepSeek-R1}~\cite{shao2024deepseekmath, deepseek2025r1} --- 组相对策略优化与涌现推理。
  % \item \textbf{ReAct}~\cite{yao2023react} --- Reasoning + Acting framework for LLM agents.
  \item \textbf{ReAct}~\cite{yao2023react} --- LLM Agent 的 Reasoning + Acting 框架。
  % \item \textbf{Toolformer}~\cite{schick2023toolformer} --- Teaching LLMs to use tools.
  \item \textbf{Toolformer}~\cite{schick2023toolformer} --- 教 LLM 学会使用工具。
  % \item \textbf{RAG}~\cite{lewis2020retrieval} --- Retrieval-Augmented Generation.
  \item \textbf{RAG}~\cite{lewis2020retrieval} --- 检索增强生成（Retrieval-Augmented Generation）。
\end{itemize}

\subsection{系统与扩展}
\label{systems-and-scaling}

\begin{itemize}
  % \item \textbf{Megatron-LM}~\cite{shoeybi2019megatron} --- Tensor and pipeline parallelism.
  \item \textbf{Megatron-LM}~\cite{shoeybi2019megatron} --- 张量并行与流水线并行。
  % \item \textbf{DeepSpeed ZeRO}~\cite{rajbhandari2020zero} --- Memory-efficient distributed training.
  \item \textbf{DeepSpeed ZeRO}~\cite{rajbhandari2020zero} --- 显存高效的分布式训练。
  % \item \textbf{vLLM}~\cite{kwon2023efficient} --- PagedAttention for efficient LLM serving.
  \item \textbf{vLLM}~\cite{kwon2023efficient} --- 为高效 LLM 服务而生的 PagedAttention。
  % \item \textbf{FlashAttention}~\cite{dao2022flashattention} --- IO-aware exact attention.
  \item \textbf{FlashAttention}~\cite{dao2022flashattention} --- IO 感知的精确 attention。
\end{itemize}

\subsection{智能体 AI}
\label{agentic-ai}

\begin{itemize}
  % \item \textbf{Building Effective Agents}~\cite{anthropic2024buildingagents} --- Design patterns and principles.
  \item \textbf{Building Effective Agents}~\cite{anthropic2024buildingagents} --- 设计模式与原则。
  % \item \textbf{Voyager}~\cite{wang2023voyager} --- Open-ended agent with skill library in Minecraft.
  \item \textbf{Voyager}~\cite{wang2023voyager} --- Minecraft 中带有技能库的开放式 Agent。
  % \item \textbf{SWE-bench}~\cite{jimenez2024swebench} --- Benchmark for autonomous software engineering.
  \item \textbf{SWE-bench}~\cite{jimenez2024swebench} --- 自主软件工程基准。
  % \item \textbf{OSWorld}~\cite{xie2024osworld} --- Full computer-use benchmarks.
  \item \textbf{OSWorld}~\cite{xie2024osworld} --- 完整的计算机使用基准。
  % \item \textbf{GAIA}~\cite{mialon2023gaia} --- General AI Assistants benchmark for real-world tasks.
  \item \textbf{GAIA}~\cite{mialon2023gaia} --- 面向真实任务的通用 AI 助手基准。
  % \item \textbf{MemGPT}~\cite{packer2023memgpt} --- OS-inspired memory management for unbounded context.
  \item \textbf{MemGPT}~\cite{packer2023memgpt} --- 受操作系统启发、面向无界上下文的记忆管理。
  % \item \textbf{Model Context Protocol}~\cite{anthropic-mcp-2024} --- Open standard for tool integration.
  \item \textbf{Model Context Protocol}~\cite{anthropic-mcp-2024} --- 工具集成的开放标准。
  % \item \textbf{Agent-to-Agent Protocol}~\cite{google-a2a-2025} --- Inter-agent communication standard.
  \item \textbf{Agent-to-Agent Protocol}~\cite{google-a2a-2025} --- 智能体间通信标准。
\end{itemize}

\subsection{对齐与安全}
\label{alignment-and-safety}

\begin{itemize}
  % \item \textbf{Constitutional AI}~\cite{bai2022constitutional} --- Self-supervised alignment.
  \item \textbf{Constitutional AI}~\cite{bai2022constitutional} --- 自监督对齐。
  % \item \textbf{Sleeper Agents}~\cite{hubinger2024sleeper} --- Deceptive alignment concerns.
  \item \textbf{Sleeper Agents}~\cite{hubinger2024sleeper} --- 关于欺骗性对齐的担忧。
  % \item \textbf{Reflexion}~\cite{shinn2023reflexion} --- Learning from verbal self-reflection.
  \item \textbf{Reflexion}~\cite{shinn2023reflexion} --- 从口头自我反思中学习。
  % \item \textbf{Indirect Prompt Injection}~\cite{greshake2023indirect} --- Security risks for LLM-integrated applications.
  \item \textbf{Indirect Prompt Injection}~\cite{greshake2023indirect} --- LLM 集成应用的安全风险。
\end{itemize}

\subsection{在线资源}
\label{online-resources}

\begin{itemize}
  % \item \textbf{HuggingFace TRL}: \href{https://github.com/huggingface/trl}{https://github.com/huggingface/trl} --- Production RL library.
  \item \textbf{HuggingFace TRL}：\href{https://github.com/huggingface/trl}{https://github.com/huggingface/trl} --- 生产级 RL 库。
  % \item \textbf{LangGraph}: \href{https://github.com/langchain-ai/langgraph}{https://github.com/langchain-ai/langgraph} --- Agent workflow graphs.
  \item \textbf{LangGraph}：\href{https://github.com/langchain-ai/langgraph}{https://github.com/langchain-ai/langgraph} --- Agent 工作流图。
  % \item \textbf{OpenAI Agents SDK}: \href{https://github.com/openai/openai-agents-python}{https://github.com/openai/openai-agents-python} --- Official agent framework.
  \item \textbf{OpenAI Agents SDK}：\href{https://github.com/openai/openai-agents-python}{https://github.com/openai/openai-agents-python} --- 官方 Agent 框架。
  % \item \textbf{DeepSpeed-Chat}: \href{https://github.com/microsoft/DeepSpeedExamples}{https://github.com/microsoft/DeepSpeedExamples} --- End-to-end RLHF pipeline.
  \item \textbf{DeepSpeed-Chat}：\href{https://github.com/microsoft/DeepSpeedExamples}{https://github.com/microsoft/DeepSpeedExamples} --- 端到端 RLHF 流水线。
  % \item \textbf{DSPy}: \href{https://github.com/stanfordnlp/dspy}{https://github.com/stanfordnlp/dspy} --- Declarative prompt optimization.
  \item \textbf{DSPy}：\href{https://github.com/stanfordnlp/dspy}{https://github.com/stanfordnlp/dspy} --- 声明式 prompt 优化。
  % \item \textbf{AutoGen}: \href{https://github.com/microsoft/autogen}{https://github.com/microsoft/autogen} --- Multi-agent conversation framework.
  \item \textbf{AutoGen}：\href{https://github.com/microsoft/autogen}{https://github.com/microsoft/autogen} --- 多智能体对话框架。
\end{itemize}

% \emph{``The best way to predict the future is to build it.''}\\
\emph{“预测未来最好的方式，就是把它创造出来。”}\\

--- Alan Kay
